\section{GHL2025 One-Term Robin Block Encoding: Lean Proof Map}

This note is a human-readable companion to
\texttt{QuantumBlockEncoding/GHL2025.lean},
\texttt{QuantumBlockEncoding/RobinMatrix.lean}, and
\texttt{conversion-windows/QBE-AUTO-001.md}.  It is not the checked source of
truth.  The Lean files are the source of truth; this file records the theorem
shape, notation, and proof obligations so that a human can compare the
formalization with the paper.

\subsection{Target Operator}

The one-term Robin construction block-encodes the discretized operator
\[
  A_k = \operatorname{diag}(f(x_0),\ldots,f(x_{N-1})) D_{\mathrm{Robin}},
  \qquad N = 2^n .
\]
Here \(D_{\mathrm{Robin}}\) is the finite-difference derivative matrix after
Robin ghost-point elimination.  In Lean:
\[
  D_{\mathrm{Robin}}
  \quad\leftrightarrow\quad
  \texttt{Examples.RobinHeat.robinDerivativeMatrix n}.
\]
The theorem-facing matrix is the row-scaled target
\[
  A_k
  \quad\leftrightarrow\quad
  \texttt{Examples.RobinHeat.oneTermRobinAkMatrix n},
\]
whose entries are \(f(x_i)D_{ij}\).

\subsection{Robin Boundary Matrix}

The bulk stencil is represented by
\[
  \texttt{Examples.RobinHeat.centralBulkEntries}
\]
and the boundary rows by
\[
  \texttt{leftBoundaryRow0},\quad
  \texttt{leftBoundaryRow1},\quad
  \texttt{rightBoundaryRowNm2},\quad
  \texttt{rightBoundaryRowNm1}.
\]
The Lean constructor
\[
  \texttt{buildRobinMatrix}
\]
lifts these row stencils into a concrete finite matrix
\[
  \texttt{Matrix (gridSize n) (gridSize n) Coeff}.
\]
The tests in \texttt{Tests/Basic.lean} cross-check every nonzero boundary entry
of the paper's displayed Robin matrix for \(n=3\).

\subsection{Normalizer and Register Layout}

The paper's one-term normalizer is
\[
  \alpha = N_D N_f \kappa.
\]
In Lean this is
\[
  \texttt{GHL2025.oneTermRobinNormalizer}.
\]
The theorem-level register layout is
\[
  m = \lceil \log_2 n\rceil
    + \lceil \log_2 G_f\rceil
    + \lceil \log_2 \kappa\rceil
    + 4
\]
signal qubits and \(2n\) pure ancillas.  In Lean this is represented by
\[
  \texttt{GHL2025.oneTermRobinLayout}
  \quad\text{and}\quad
  \texttt{GHL2025.defaultOneTermRobinTheoremData}.
\]

\subsection{Circuit Skeleton}

The paper circuit is recorded as the Lean data structure
\[
  \texttt{GHL2025.RobinCircuitSkeleton}.
\]
The default skeleton contains the labeled operations
\[
  U_{\mathrm{indic}},\quad O_{D^T}^{S},\quad R_y^{\mathrm{boundary}},
  \quad O_D^{BS},\quad O_f,\quad \mathrm{SWAP},\quad (O_D^{BS})^\dagger .
\]
This is still a skeleton until each labeled gate is assigned a certified matrix
semantics.

\subsection{Register Bit Layout}

The circuit-level register ordering (LSB = bit 0) for the Robin circuit
is formalized in \texttt{GHL2025.RobinRegisterPartition}.  The bit positions
from LSB to MSB are:

\begin{center}
\begin{tabular}{llll}
\hline
Register & Bit range & Width & Field \\
\hline
ancilla & [0] & 1 & \texttt{ancillaQubit} \\
system ($j$) & [1, $1{+}n$) & $n$ & \texttt{systemQubits} \\
od\_pure\_ancilla & [$1{+}n$, $1{+}n{+}$odPure) & $n - \lceil\log_2\kappa\rceil$ & \texttt{odPureAncillaQubits} \\
sparse\_index ($s$) & [$1{+}n{+}$odPure, ...) & $\lceil\log_2\kappa\rceil$ & \texttt{sparseIndexQubits} \\
indicator (ind) & at bit $1 + 2n$ & 1 & \texttt{indicatorQubit} \\
mf & remaining MSBs & $\lceil\log_2 n\rceil + \lceil\log_2 G_f\rceil + 3$ & \texttt{mfQubits} \\
\hline
\end{tabular}
\end{center}

The total circuit qubits equals the register partition total:
\[
  \texttt{totalQubits} = 2n + \lceil\log_2 n\rceil + \lceil\log_2 G_f\rceil + 5.
\]
For $n=3$, $G_f=1$, $\kappa=7$: total $= 6 + 2 + 0 + 5 = 13$ qubits.

\subsection{Indicator Oracle $U_{\mathrm{indic}}$ Matrix}

The indicator oracle is a permutation matrix (controlled-X on the indicator
qubit conditioned on the system register being in the bulk window):
\[
  U_{\mathrm{indic}} \,|\mathit{mf}\rangle|ind\rangle|s\rangle|od\rangle|j\rangle|anc\rangle
  =
  |\mathit{mf}\rangle|ind \oplus \mathrm{isBulk}(K_1,K_2,j)\rangle|s\rangle|od\rangle|j\rangle|anc\rangle
\]
where $\mathrm{isBulk}(K_1,K_2,j) = 1$ iff $K_1 \leq j \leq K_2$.
In Lean this is recorded as
\[
  \texttt{GHL2025.indicatorOracleMatrix}.
\]
The unitarity is proved: the image function
$\texttt{indicatorOracleImage}$ is self-inverse (XOR of a single bit
at position $\texttt{indPos} = 1 + 2n$), hence bijective on the
$\texttt{Fin}$ domain, and the matrix has exactly one $\texttt{Coeff.rat}\;1$
per row and per column.  This is recorded as
\texttt{indicatorOracleMatrix\_is\_permutation}.  The gate matrix
obligation is now $\texttt{proved := true}$.

\subsection{SWAP Gate Matrix}

The SWAP gate exchanges two $n$-qubit register blocks in the compound index:
the system register at bits $[1, 1+n)$ and the O$_D^{BS}$ register at bits $[1+n, 1+2n)$.
The permutation is:
\[
  \mathrm{swap}(j) = j \oplus (d \ll 1) \oplus (d \ll (1+n)),
  \qquad
  d = (j \gg 1 \,\&\, \text{mask}) \oplus (j \gg (1{+}n) \,\&\, \text{mask})
\]
where mask $= (1 \ll n) - 1$.
When the two block values are equal, the SWAP acts as the identity.
All bits outside the two blocks are preserved.
In Lean this is recorded as
\[
  \texttt{GHL2025.swapOracleMatrix}.
\]
The unitarity proof follows the same self-inverse template as $U_{\mathrm{indic}}$:
after swap, block1' $=$ block2 and block2' $=$ block1, so
$d' = \text{block1}' \oplus \text{block2}' = d$, and applying the XOR pattern
twice cancels.  The key bit-arithmetic lemmas are:
\begin{itemize}
  \item $(d \ll n) \,\&\, \text{mask} = 0$ when $d < 2^n$
    (reuse \texttt{shiftLeft\_land\_mask\_eq\_zero} from the U\_indic proof-DAG).
  \item $d \gg n = 0$ when $d < 2^n$.
  \item Right-shift distributes over XOR.
\end{itemize}
Lean has proved both register-image blocks:
\texttt{swapOracleImage\_block1\_eq\_block2} proves that after
\texttt{swapOracleImage} the low $n$-bit block equals the old high $n$-bit
block, and \texttt{swapOracleImage\_block2\_eq\_block1} proves the symmetric
high-block statement.  Lean also proves
\texttt{swapOracleImage\_lt\_qubitDim}, which keeps the image inside the full
finite basis.  The accepted self-inverse block adds
\texttt{swapOracleDiff}, \texttt{swapOracleDiff\_preserved},
\texttt{xor\_two\_shifted\_masks\_cancel}, and
\texttt{swapOracleImage\_self\_inverse}.  The finite permutation bridge then
proves \texttt{swapOracleImage\_injective},
\texttt{swapOracleImage\_bijective}, row and column uniqueness for
\texttt{swapOracleMatrix}, and \texttt{swapOracleMatrix\_is\_permutation}.
The SWAP gate obligation is now \texttt{proved := true}.

\subsection{Banded Sparse Access Oracle $O_D^{BS}$ Matrix}

The banded sparse access oracle (Lemma 1 of the GHL paper,
arXiv:2506.20478) maps
\[
  \hat O_D^{BS}|0\rangle^{n-l}|s\rangle^l|i\rangle^n
  =
  |r_{si}\rangle^n|i\rangle^n,
\]
where $r_{si}=r_{s0}+i \bmod 2^n$ is the $s$-th banded-sparse matrix index.
For the Robin one-term circuit, the same register-level contract is used with
$D$ or $D^T$ and then cleaned by $(O_D^{BS})^\dagger$ after the SWAP step.

\textbf{Faithfulness correction.}  The current Lean
\texttt{bandedSparseAccessMatrix} is an interim column-map helper.  It preserves
the sparse-index register and overwrites the system register with
\texttt{robinSparseColumnMap}.  This is useful for testing Robin stencil
columns, but it is not yet the paper's padded sparse-index oracle.  The
unitarity and cleanup obligations for $O_D^{BS}$ must therefore remain open
until the Lean register-level image formula is corrected to the paper contract.

The source contract is now recorded explicitly in Lean as
\[
  \texttt{GHL2025.BandedSparseAccessPaperContract}.
\]
For one-term Robin parameters, the default instance is
\[
  \texttt{GHL2025.defaultBandedSparseAccessPaperContract}.
\]
This record stores the input ket
\[
  |0\rangle^{n-l}|s\rangle^l|i\rangle^n
\]
and the output ket
\[
  |r_{si}\rangle^n|i\rangle^n,
\]
together with three unproved obligations: the padded-width compatibility,
the clean-input domain, the forward oracle correctness, the cleanup by
$(O_D^{BS})^\dagger$, and a full-unitary extension on basis columns whose
padded register is not clean.  The record does not prove unitarity and does not change
\texttt{bandedSparseAccessMatrix}; it prevents the interim helper from being
mistaken for the paper oracle.

The Lean-facing proof DAG for this oracle is:
\begin{center}
\begin{tabular}{p{0.24\textwidth}p{0.40\textwidth}p{0.22\textwidth}}
\hline
Block & Interface & Status \\
\hline
\texttt{odbs\_width\_compatible} &
the padded zero register and sparse-index register have total width $n$ &
unproved \\
\texttt{odbs\_extract\_registers} &
extract the padded sparse-index register and row register from the compound
basis index &
defined skeleton \\
\texttt{odbs\_clean\_domain} &
check whether the padded register is $|0\rangle^{n-l}$ and record that
non-clean columns require a unitary completion &
defined skeleton; obligations unproved \\
\texttt{odbs\_valid\_sparse\_branch\_domain} &
classify row-dependent sparse branches that carry nonzero Robin stencil data &
defined audit predicate; not the full faithful domain \\
\texttt{odbs\_unused\_branch\_image\_rule} &
record the missing reversible image rule for clean invalid sparse branches &
interface defined with no proposed image; obligations unproved \\
\texttt{odbs\_full\_clean\_domain\_extension\_contract} &
combine valid branches and clean unused branches into one paper-level extension
contract &
defined obligation wrapper; nested unused-branch image remains unspecified and
every semantic flag remains false \\
\texttt{odbs\_address\_range} &
check that the written address $r_{si}$ is an $n$-bit value before it is
inserted into the $O_D^{BS}$ register &
executable check defined; obligation unproved \\
\texttt{odbs\_no\_spill} &
check that the image preserves the high tail above the $O_D^{BS}$ register,
including the indicator and $m_f$ bits &
executable check defined; obligation unproved \\
\texttt{odbs\_forward\_image} &
map $|0\rangle^{n-l}|s\rangle^l|i\rangle^n$ to
$|r_{si}\rangle^n|i\rangle^n$ while preserving the row register &
defined skeleton; correctness unproved \\
\texttt{odbs\_forward\_matrix} &
set matrix entries by the paper image, with
$M[\operatorname{image}(j),j]=1$ &
finite-image entry bridge proved under explicit hypotheses; injectivity unproved \\
\texttt{odbs\_dagger\_matrix} &
set transpose-style entries by the paper image, with
$M^\dagger[i,j]=1$ exactly when $j=\operatorname{image}(i)$ &
active skeleton; inverse-on-range unproved \\
\texttt{odbs\_dagger\_cleanup} &
clean the padded sparse-index register using $(O_D^{BS})^\dagger$ after SWAP &
unproved \\
\hline
\end{tabular}
\end{center}

Lean now defines \texttt{bandedSparseAccessPaperRegisters},
\texttt{bandedSparseAccessPaperCleanInput},
\texttt{bandedSparseAccessPaperColumnContract},
\texttt{bandedSparseAccessPaperAddress}, and
\texttt{bandedSparseAccessPaperImage}.  Lean also defines
\texttt{bandedSparseAccessPaperAddressInRange},
\texttt{bandedSparseAccessPaperHighTail}, and
\texttt{bandedSparseAccessPaperImageNoSpill}.  These executable checks record
that the written $r_{si}$ value should fit in $n$ bits and that the image should
not spill into the indicator or $m_f$ bits above the $O_D^{BS}$ register.  Lean
also defines
\texttt{bandedSparseAccessPaperMatrix} and
\texttt{bandedSparseAccessPaperDaggerMatrix} from that image function.  These
declarations preserve the row register and replace the O$_D^{BS}$ register by
the one-term Robin value of $r_{si}$ on the executable skeleton.  The
per-column contract records whether the padded-zero field is clean, because
Lemma 1 specifies only inputs of the form
$|0\rangle^{n-l}|s\rangle^l|i\rangle^n$.  Concrete tests cover one clean bulk
row and one boundary row at $n=3$ and $\kappa=7$, and also record a non-clean
padded-register column at $n=4$ and $\kappa=7$ whose unitary completion remains
unproved.  The active gate pair now uses these paper-image matrices:
\texttt{oneTermRobinGate\_O\_D\_BS} uses
\texttt{bandedSparseAccessPaperMatrix}, and
\texttt{oneTermRobinGate\_O\_D\_BS\_dagger} uses
\texttt{bandedSparseAccessPaperDaggerMatrix}.  The contract remains the
paper's Lemma 1 map
\[
  |0\rangle^{n-l}|s\rangle^l|i\rangle^n
  \longmapsto
  |r_{si}\rangle^n|i\rangle^n.
\]
The cycle 10 source-contract audit checked the arXiv v2 source against the
row-dependent valid-branch predicate.  Lemma
\texttt{Banded-sparse-access-oracle and resource cost} gives the padded-register
map above.  The Robin one-term construction then says that zeros can be
included in the set of nonzero elements and uses diagonal sparsity
\(\kappa\), while Eq.~\texttt{ROBIN clarified} sums over
\(s=0,\ldots,\kappa-1\).  Thus the Lean predicate
\texttt{bandedSparseAccessPaperValidCleanSource} classifies branches carrying
nonzero stencil data, but it is not a faithful standalone source-domain
restriction.  The faithful contract still needs an unused-branch extension that
agrees with Lemma 1 on valid branches and records injectivity, cleanup, and
unitary extension as obligations until proved.

The unused-branch contract is now explicit in Lean.  The classifier
\texttt{bandedSparseAccessPaperUnusedSparseBranch} identifies clean padded
columns outside the row-dependent nonzero-stencil predicate, and
\texttt{bandedSparseAccessUnusedBranchExtensionContract\_of\_unusedBranch}
packages the corresponding false image-rule, injectivity, cleanup, and unitary
extension obligations.  This bridge does not choose a reversible image for
unused branches and does not change the active paper-image matrix pair.

The image-rule interface is now a separate Lean record:
\texttt{BandedSparseAccessUnusedBranchImageRuleContract}.  Its default instance
\texttt{bandedSparseAccessUnusedBranchImageRuleContract} records the source
column, extracted registers, current active image, and branch classification,
but keeps \texttt{proposedImageIndex = none}.  The fields
\texttt{imageSpecified}, \texttt{imageFinite},
\texttt{separatesActiveCollision}, and \texttt{validBranchAgreement} remain
false obligations.  This is only a contract interface; it does not choose a
replacement image or rewire the active O$_D^{BS}$ matrices.

Lean now records the paper-level wrapper
\texttt{BandedSparseAccessFullCleanDomainExtensionContract}.  This wrapper
combines the valid-branch agreement obligation with the unused-branch image-rule
interface, full clean-domain injectivity, dagger cleanup, and unitary extension.
No reversible-completion theorem has been cited for this step, so the default
wrapper keeps the nested unused-branch image unspecified
(\texttt{proposedImageIndex = none}) and keeps every semantic field false.
The local classifier split is now Lean-proved:
\texttt{bandedSparseAccessPaperCleanDomainSplit\_iff} states that padded-clean
input is exactly either a valid row-dependent sparse branch or a clean unused
sparse branch, and
\texttt{bandedSparseAccessPaperCleanDomainSplit\_disjoint} states that these
two cases are disjoint.  The wrapper-facing theorem
\texttt{bandedSparseAccessFullCleanDomainExtensionContract\_localCleanDomainSplit}
keeps \texttt{cleanDomainSplit.proved = false}; the proof does not supply an
unused-branch image, injectivity, cleanup, or unitarity.

Cycle 14 promotes the source decision itself into Lean as
\texttt{BandedSparseAccessUnusedZeroBranchSourceDecision} and
\texttt{bandedSparseAccessUnusedZeroBranchSourceDecision}.  The cited-result
key is \texttt{QBE.ODBS.UnusedZeroBranchExtension}.  The fields
\texttt{paperImageRuleSpecified},
\texttt{externalExtensionTheoremAccepted}, and
\texttt{lowerProofSearchAllowed} are all false, and the dependency record has
\texttt{proved = false}.  The theorem
\texttt{bandedSparseAccessUnusedZeroBranchSourceDecision\_keepsFullDomainFlagsFalse}
ties this source decision to the full clean-domain wrapper without changing the
active $O_D^{BS}$ matrices or promoting cleanup, injectivity, unitarity, or
block-correctness flags.

Lean now proves a narrow address-range block:
\texttt{bandedSparseAccessPaperAddressInRange\_eq\_true\_of\_two\_le}
shows that the executable address check is true under the explicit grid
condition $2 \le n$.  This uses
\texttt{bandedSparseAccessPaperRegisters\_row\_lt\_gridSize} for the extracted
row field and
\texttt{robinSparseColumnMap\_lt\_gridSize\_of\_row\_lt} for the fourth-order
stencil column map.  This theorem is not yet a proof of the paper-level
contract field, because \texttt{OneTermRobinParameters} does not record the
$2 \le n$ side condition.  Lean also proves
\texttt{bandedSparseAccessPaperImageNoSpill\_iff}, which rewrites the no-spill
Boolean as equality of the high tails above the $O_D^{BS}$ register.  The
address-conditional high-tail theorem
\texttt{bandedSparseAccessPaperImage\_highTail\_eq\_of\_address\_lt} now proves
that equality from an $n$-bit written address, and
\texttt{bandedSparseAccessPaperImageNoSpill\_eq\_true\_of\_two\_le} supplies a
concrete route under $2 \le n$.  The paper-level no-spill flag still remains an
obligation because the source contract has not bundled the parameter-family
side condition.  Lean now also proves the image-range and register-roundtrip
block.  The theorem
\texttt{bandedSparseAccessPaperImage\_lt\_qubitDim\_of\_address\_lt} shows that
the executable image is below
\texttt{qubitDim (oneTermRobinTotalQubits p)} when the source column is in that
finite basis and the written address is $n$-bit.  The theorem
\texttt{bandedSparseAccessPaperImage\_rowValue\_eq} proves that register
extraction from the image preserves the row field.  The theorem
\texttt{bandedSparseAccessPaperImage\_odRegisterValue\_eq} proves, under the
same $n$-bit address hypothesis, that extraction from the image reports the
written $r_{si}$ value in the $O_D^{BS}$ register.  These are executable
register-splice facts, not a promotion of the paper-level oracle contract.
Lean now also packages the executable image as a finite basis index through
\texttt{bandedSparseAccessPaperImageFin} under the same source-column and
$n$-bit address hypotheses.  The bridge theorems
\texttt{bandedSparseAccessPaperMatrix\_imageFin\_eq\_one} and
\texttt{oneTermRobinGate\_O\_D\_BS\_imageFin\_eq\_one} prove the forward entry
$M[\operatorname{image}(j),j]=1$ at that finite image index.  The paired
transpose-style theorems
\texttt{bandedSparseAccessPaperDaggerMatrix\_imageFin\_eq\_one} and
\texttt{oneTermRobinGate\_O\_D\_BS\_dagger\_imageFin\_eq\_one} prove the entry
$M^\dagger[j,\operatorname{image}(j)]=1$.  These entry facts are prerequisites
for injectivity and inverse-on-range work; they do not prove either property.
No unitarity, clean-domain, address-range, no-spill, unitary-extension,
forward-correctness, cleanup, or block-correctness flag is promoted.  The
cleanup statement for $(O_D^{BS})^\dagger$ after SWAP remains an unproved
obligation.  The focused activation tests check the forward entries
$M[40,8]=1$ and $M[4,8]=0$, and the dagger entry $M^\dagger[8,40]=1$ for
$n=3$ and $\kappa=7$.  The focused register-safety tests also check
\texttt{bandedSparseAccessPaperAddressInRange} and
\texttt{bandedSparseAccessPaperImageNoSpill} on the clean bulk column
\texttt{j = 8}, high-tail preservation for \texttt{j = 136}, and general
examples for the image-range, row-roundtrip, and address-roundtrip lemmas.
The remaining fixed work for this oracle is injectivity, inverse-on-range, and
post-SWAP cleanup.  The completed roundtrip facts must not promote
\texttt{oneTermRobinGate\_O\_D\_BS.unitary.proved},
\texttt{oneTermRobinGate\_O\_D\_BS\_dagger.unitary.proved},
\texttt{forwardCorrect}, \texttt{daggerCleanup}, or any block-correctness flag.

For the fourth-order central second-derivative stencil (half-bandwidth $l=2$):
\begin{itemize}
  \item \textbf{Bulk rows} ($K_1 \leq i \leq K_2$): 5 entries at offsets $\{-2,-1,0,1,2\}$.
    $\mathrm{col}(s,i) = i - 2 + s$ for $s < 5$; identity for $s \geq 5$.
  \item \textbf{Left boundary row 0} (3 entries, offsets $\{0,1,2\}$):
    $\mathrm{col}(s,0) = s$ for $s < 3$; identity for $s \geq 3$.
  \item \textbf{Left boundary row 1} (4 entries, offsets $\{-1,0,1,2\}$):
    $\mathrm{col}(s,1) = s$ for $s < 4$; identity for $s \geq 4$.
  \item \textbf{Right boundary row $N{-}2$} (4 entries, offsets $\{-2,-1,0,1\}$):
    $\mathrm{col}(s,N{-}2) = N{-}4+s$ for $s < 4$; identity for $s \geq 4$.
  \item \textbf{Right boundary row $N{-}1$} (3 entries, offsets $\{-2,-1,0\}$):
    $\mathrm{col}(s,N{-}1) = N{-}3+s$ for $s < 3$; identity for $s \geq 3$.
\end{itemize}

The interim matrix form replaces the system register bits while preserving the
sparse index and all other bits.  This sentence describes the Lean helper, not
the final faithful oracle contract.

In Lean:
\[
  \texttt{GHL2025.robinSparseColumnMap} \quad (\text{column mapping}),
  \qquad
  \texttt{GHL2025.bandedSparseAccessMatrix}.
\]
The current column-map contract has a boundary-row non-injectivity witness:
\[
  \texttt{robinSparseColumnMap}\;3\;0\;0
  =
  \texttt{robinSparseColumnMap}\;3\;0\;1
  =
  \texttt{robinSparseColumnMap}\;3\;0\;2
  =
  0.
\]
Thus the permutation-matrix proof route is blocked for the simplified helper.
This should be read as a contract-drift finding for the legacy helper, not as a
missing construction in the GHL paper.  The active paper-image skeleton is now
wired into the gate pair, but the unitarity obligation remains unproved
($\texttt{proved := false}$) until address range, no-spill, range,
injectivity, inverse-on-range, and cleanup are proved for that active contract.

\subsection{$(O_D^{BS})^\dagger$ Inverse Matrix}

The active Lean implementation records a transpose-style matrix by computing
$\mathrm{image}(i)$ from the paper-image skeleton and checking whether
$j = \mathrm{image}(i)$.  This is
\texttt{bandedSparseAccessPaperDaggerMatrix}.  It is wired into
\texttt{oneTermRobinGate\_O\_D\_BS\_dagger}, but Lean has not proved that it is
an inverse on the range of the forward image, a unitary matrix, or the
clean-ancilla operation required after SWAP.

In Lean:
\[
  \texttt{GHL2025.bandedSparseAccessPaperDaggerMatrix}.
\]
The unitarity obligation remains unproved ($\texttt{proved := false}$).

The legacy helper \texttt{bandedSparseAccessDaggerMatrix} remains in Lean only
as the transpose-style matrix for the old sparse-index helper.  It is not the
active paper-contract dagger gate.

The bulk window bounds are derived from the stencil radius:
\[
  K_1 = \mathrm{leftRadius} = 2, \qquad
  K_2 = 2^n - \mathrm{rightRadius} - 1 = 2^n - 3.
\]
The fourth-order central second-derivative stencil
(\texttt{fourthOrderSecondDerivative}) has $\mathrm{leftRadius} = \mathrm{rightRadius} = 2$.
The precondition $K_2 \geq K_1$ requires $2^n \geq 5$, i.e.\ $n \geq 3$.

\subsection{Sparse Amplitude Data}

The function
\[
  \texttt{robinSparseAmplitudeValue}(n, s, i) : \texttt{Coeff}
\]
returns the $s$-th nonzero stencil coefficient of row $i$ in the Robin
derivative matrix.  This is the data layer shared by $O_{D^T}^{S}$
(sparse amplitude oracle, Lemma 3) and $R_y^{\mathrm{boundary}}$
(boundary-controlled rotations).

The function mirrors the case structure of \texttt{robinSparseColumnMap}
but returns coefficient values instead of column indices.  For bulk rows
($K_1 \leq i \leq K_2$) the 5 stencil coefficients are
$\{-1/12,\, 4/3,\, -5/2,\, 4/3,\, -1/12\}$ at offsets
$\{-2,\,-1,\,0,\,1,\,2\}$.  Boundary rows include the Robin correction
terms involving $A_1\Delta x$ (left) and $B_1\Delta x$ (right).
Unused sparse indices return $\texttt{Coeff.rat}\;0$.

\textbf{Coherence:} for all valid $(n, s, i)$:
\[
  \texttt{robinSparseAmplitudeValue}(n, s, i)
  \;=\;
  D_{\mathrm{Robin}}[i][\texttt{robinSparseColumnMap}(n, s, i)].
\]
This is tested with \texttt{native\_decide} for concrete $n = 3$ instances
covering bulk and all four boundary row types.

\subsection{Signal-Qubit Counts}

The project uses two signal-qubit counts:

\begin{itemize}
  \item \textbf{Theorem-level} (\texttt{oneTermRobinLayout}):
    $m = \lceil\log_2 n\rceil + \lceil\log_2 G_f\rceil + \lceil\log_2 \kappa\rceil + 4$.
  \item \textbf{Circuit-level} (\texttt{effectiveRobinSignalQubits}):
    $m' = \mathrm{totalQubits} - n$.
\end{itemize}

These differ because the circuit-level count includes all non-system registers
(ancilla, indicator, pure ancillas).  The block extraction uses $m'$ since the
projection $(\langle 0^a| \otimes I)$ zeros out all non-system registers.

\subsection{Block-Encoding Statement}

The desired final theorem is the block extraction equation
\[
  (\langle 0^a| \otimes I) U (|0^a\rangle \otimes I)
  =
  \frac{A_k}{N_D N_f \kappa}.
\]
At the current stage, Lean records the target matrix, circuit labels,
normalizer, resource counts, proof obligations, and a circuit matrix
semantics backend whose full-space matrix is the product of the seven current
gate matrices.  The block-extraction target is structural: it fixes the signal
projection, target matrix, and normalizer, while the block equation and several
gate-level semantic claims remain unproved.

\subsection{Current Lean Correspondence}

\begin{center}
\begin{tabular}{lll}
\hline
Paper object & Lean declaration & Status \\
\hline
\(D_{\mathrm{Robin}}\) &
\texttt{Examples.RobinHeat.robinDerivativeMatrix} &
defined and tested \\
\(A_k=\operatorname{diag}(f)D_{\mathrm{Robin}}\) &
\texttt{Examples.RobinHeat.oneTermRobinAkMatrix} &
defined and wired as theorem target \\
\(\alpha=N_DN_f\kappa\) &
\texttt{GHL2025.oneTermRobinNormalizer} &
defined \\
Theorem tuple \((\alpha,m,a)\) &
\texttt{GHL2025.OneTermRobinTheoremData} &
defined \\
Fig. one-term Robin circuit &
\texttt{GHL2025.RobinCircuitSkeleton} &
defined as skeleton \\
Eq. ROBIN intermediate states &
\texttt{GHL2025.RobinGamma1/2/3} &
typed data \\
Proof obligations &
\texttt{GHL2025.RobinProofObligations} &
tracked, unproved \\
Circuit matrix semantics &
\texttt{QuantumBlockEncoding.CircuitSemantics} &
started \\
Gate matrix placeholders &
\texttt{GHL2025.oneTermRobinGateMatrixPlaceholders} &
defined; active O$_D^{BS}$ pair uses the paper-image skeleton \\
Block projection &
\texttt{signalSystemBlockProjection} &
defined \\
Robin circuit matrix semantics &
\texttt{Examples.RobinHeat.oneTermRobinCircuitSemantics} &
defined as an \texttt{evalGateMatrices} product \\
Robin block extraction target &
\texttt{Examples.RobinHeat.oneTermRobinBlockExtractionTarget} &
defined, unproved \\
Circuit block encoding claim &
\texttt{CircuitBlockEncodingClaim} &
defined \\
Robin circuit block claim &
\texttt{Examples.RobinHeat.oneTermRobinCircuitBlockClaim} &
defined, unproved \\
Dimension compatibility &
\texttt{Examples.RobinHeat.oneTermRobinCircuitDimCompat} &
proved from \texttt{clog2\_gridSize} \\
Default Robin block claim &
\texttt{Examples.RobinHeat.defaultOneTermRobinCircuitBlockClaim} &
defined, unproved \\
Register bit layout &
\texttt{GHL2025.robinIndicatorBitPosition} &
implemented (cycle 2) \\
U\_indic honest matrix &
\texttt{GHL2025.indicatorOracleMatrix} &
implemented (cycle 2), unitarity \textbf{proved} (run 02 cycle 02) \\
Effective signal qubits &
\texttt{GHL2025.effectiveRobinSignalQubits} &
implemented (cycle 2) \\
SWAP honest matrix &
\texttt{GHL2025.swapOracleMatrix} &
implemented (cycle 3), unitarity \textbf{proved} by
\texttt{GHL2025.swapOracleMatrix\_is\_permutation} \\
$O_D^{BS}$ column map &
\texttt{GHL2025.robinSparseColumnMap} &
implemented as stencil helper, reused to compute $r_{si}$ \\
$O_D^{BS}$ legacy helper matrix &
\texttt{GHL2025.bandedSparseAccessMatrix} &
implemented, not the active paper oracle \\
$O_D^{BS}$ paper contract &
\texttt{GHL2025.BandedSparseAccessPaperContract} &
defined, obligations unproved \\
$O_D^{BS}$ default paper contract &
\texttt{GHL2025.defaultBandedSparseAccessPaperContract} &
defined for one-term Robin parameters \\
$O_D^{BS}$ clean-domain predicate &
\texttt{GHL2025.bandedSparseAccessPaperCleanInput} &
defined; clean-domain proof flag false \\
$O_D^{BS}$ per-column audit &
\texttt{GHL2025.bandedSparseAccessPaperColumnContract} &
defined; address-range, no-spill, and unitary-extension proof flags false \\
$O_D^{BS}$ address-range check &
\texttt{GHL2025.bandedSparseAccessPaperAddressInRange} &
defined executable check; proof flag false \\
$O_D^{BS}$ high-tail no-spill check &
\texttt{GHL2025.bandedSparseAccessPaperImageNoSpill} &
defined executable check using \texttt{GHL2025.bandedSparseAccessPaperHighTail};
proof flag false \\
$O_D^{BS}$ unused-branch image-rule interface &
\texttt{GHL2025.bandedSparseAccessUnusedBranchImageRuleContract} &
defined with \texttt{proposedImageIndex = none}; image-rule proof flags false \\
$O_D^{BS}$ full clean-domain extension wrapper &
\texttt{GHL2025.bandedSparseAccessFullCleanDomainExtensionContract} &
defined obligation wrapper; nested \texttt{proposedImageIndex = none}; no
cited reversible-completion theorem yet; all semantic fields false \\
$O_D^{BS}$ paper-image matrix &
\texttt{GHL2025.bandedSparseAccessPaperMatrix} &
active skeleton; finite-image entry bridge proved under explicit hypotheses;
injectivity and unitarity unproved \\
$(O_D^{BS})^\dagger$ paper-image matrix &
\texttt{GHL2025.bandedSparseAccessPaperDaggerMatrix} &
active transpose-style skeleton; inverse-on-range and cleanup unproved \\
$(O_D^{BS})^\dagger$ legacy helper matrix &
\texttt{GHL2025.bandedSparseAccessDaggerMatrix} &
implemented, not the active paper-contract dagger \\
Sparse amplitude data &
\texttt{GHL2025.robinSparseAmplitudeValue} &
implemented (cycle 5) \\
Function value data &
\texttt{GHL2025.robinFunctionValue} &
implemented (cycle 9) \\
$O_f$ paper registers &
\texttt{GHL2025.FunctionOraclePaperRegisters} and
\texttt{GHL2025.functionOraclePaperRegisters} &
defined skeleton for system, \(m_f\) workspace, non-\(m_f\) index, and clean
workspace flag \\
$O_f$ normalized clean-branch value &
\texttt{GHL2025.functionOracleNormalizedValue} &
symbolic \(f(x_i)N_f^{-1}\) contract; division and bound unproved \\
$O_f$ paper image &
\texttt{GHL2025.FunctionOraclePaperImage} and
\texttt{GHL2025.functionOraclePaperImage} &
clean branch, orthogonal component label, and false paper-image obligations
recorded \\
$O_f$ paper-image matrix &
\texttt{GHL2025.functionOraclePaperMatrix} &
active Phase 1 skeleton; clean-input \(m_f\) branch wired and completion
unproved \\
$O_f$ orthogonal-completion symbols &
\texttt{GHL2025.functionOracleOrthogonalEntry} &
symbolic non-clean workspace entries; orthogonality and unitarity unproved \\
$O_f$ gate matrix &
\texttt{GHL2025.oneTermRobinGate_O_f} &
uses \texttt{functionOraclePaperMatrix}; amplitude relation, \(N_f\) bound,
orthogonality, and unitarity unproved \\
$O_{D^T}^S$ diagonal helper &
\texttt{GHL2025.sparseAmplitudeOracleDTMatrix} &
implemented as legacy/data helper; not the active gate \\
$O_{D^T}^S$ active gate matrix &
\texttt{GHL2025.oneTermRobinGate_O_DT_S} &
uses controlled-rotation skeleton, unitary unproved \\
\(O_{D^T}^S\) register extraction &
\texttt{GHL2025.sparseAmplitudeOracleDTPaperRegisters} &
defined skeleton for Lemma 3 register fields \\
\(O_{D^T}^S\) coefficient-normalizer contract &
\texttt{GHL2025.sparseAmplitudeOracleDTCoefficientNormalizerContract} and
\texttt{GHL2025.sparseAmplitudeOracleDTCoefficientNormalizerObligation} &
typed contract recorded; proof flags false \\
\(O_{D^T}^S\) coefficient-normalizer proof route &
\texttt{GHL2025.sparseAmplitudeOracleDTNormalizedCoefficient} and
\texttt{GHL2025.sparseAmplitudeOracleDTCoefficientNormalizerProofRoute} &
normalized-coefficient stand-in recorded; analytic proof flags false \\
shared \(N_D\) derivative normalizer contract &
\texttt{GHL2025.DerivativeNormalizerNDContract} and
\texttt{GHL2025.derivativeNormalizerNDContract} &
common \(D_j^{(s)}/N_D\) source for \(O_{D^T}^S\) and boundary \(R_y\);
analytic proof flags false \\
\(O_{D^T}^S\) paper rotation matrix &
\texttt{GHL2025.sparseAmplitudeOracleDTRotationMatrix} &
active controlled-rotation skeleton \\
\(R_y^{\mathrm{boundary}}\) active gate matrix &
\texttt{GHL2025.oneTermRobinGate\_Ry\_boundary} &
uses symbolic boundary rotation matrix, unitary unproved \\
\(R_y^{\mathrm{boundary}}\) register extraction &
\texttt{GHL2025.boundaryRotationPaperRegisters} &
defined skeleton for the boundary rotation register fields \\
\(R_y^{\mathrm{boundary}}\) angle-normalizer contract &
\texttt{GHL2025.boundaryRotationAngleNormalizerContract} and
\texttt{GHL2025.boundaryRotationAngleNormalizerObligation} &
typed contract recorded; proof flags false \\
\(R_y^{\mathrm{boundary}}\) angle-normalizer proof route &
\texttt{GHL2025.boundaryRotationAngleNormalizerProofRoute} &
division, arccos, half-angle, normalizer-bound, and unitarity gaps separated;
proof flags false \\
\hline
\end{tabular}
\end{center}

\subsection{Boundary Rotation Angle-Normalizer Route}

The paper sets the boundary-rotation angle by
\[
  \theta_j^s=\arccos\!\left(D_j^{(s)}/N_D\right).
\]
Lean now records the formal normalized coefficient in
\texttt{GHL2025.boundaryRotationNormalizedCoefficient}.  The definition is the
Robin sparse coefficient multiplied by the symbolic reciprocal
\texttt{N\_D\_inv}; it is not a proof that \(N_D\) is nonzero or that the
coefficient language has division.

The record \texttt{GHL2025.BoundaryRotationAngleNormalizerProofRoute} refines
the remaining proof obligations for \(R_y^{\mathrm{boundary}}\).  Its false
fields name the missing division semantics, real arccos semantics, half-angle
identities, \(N_D\) bound, and two-by-two unitarity.  The Lean theorems
\texttt{GHL2025.boundaryRotationAngleNormalizerContract\_coefficient} and
\texttt{GHL2025.boundaryRotationAngleNormalizerProofRoute\_arccosArgument}
prove only definitional bridges to the existing contract data.  They do not
promote \texttt{GHL2025.oneTermRobinGate\_Ry\_boundary.unitary.proved}.

\subsection{Function Oracle $O_f$ Source Contract}

The function oracle $O_f$ is the coefficient-function oracle in
Guseynov--Huang--Liu 2025, Lemma 4 and Fig.~1-term Robin,
arXiv:2506.20478.  Its paper role is to encode the grid value
\(f(x_i)/N_f\).  The one-term Robin normalizer records this factor as
\[
  \alpha = N_D N_f \kappa .
\]

The legacy Lean matrix \texttt{GHL2025.functionOracleMatrix} is a Phase 1
function-value data helper.  For each
compound basis state \(|j\rangle\), the declaration
\texttt{GHL2025.functionOracleMatrix} extracts the system register value
\[
  \text{sysVal} = (j \gg 1) \,\&\, ((1 \ll n) - 1)
\]
and places the symbolic value
\[
  \texttt{GHL2025.robinFunctionValue}\;n\;\text{sysVal}
  =
  \texttt{Coeff.symbol}\;\texttt{"f\_\{n\}\_\{sysVal\}"}
\]
on the diagonal.  Off-diagonal entries are zero.  The diagonal helper depends
only on the system grid index; it does not depend on the sparse-index register
or on \texttt{GHL2025.robinSparseAmplitudeValue}.  It is no longer the active
gate matrix.

The active Lean gate now uses a paper-image matrix skeleton.  The paper-level
coordinate-oracle branch is
\[
  |0\rangle^{m_f}|i\rangle
  \longmapsto
  \frac{f(x_i)}{N_f}|0\rangle^{m_f}|i\rangle
  + |\mathrm{orth}_f(i)\rangle,
\]
where the second summand has zero overlap with the clean
\(|0\rangle^{m_f}\) workspace branch.  Lean now records this branch as a
source contract.  The declarations
\texttt{GHL2025.FunctionOraclePaperRegisters} and
\texttt{GHL2025.functionOraclePaperRegisters} extract the system value, the
\(m_f\) workspace value, the non-\(m_f\) basis index, and a clean-workspace
flag.  The declaration \texttt{GHL2025.functionOracleNormalizedValue} records
\texttt{GHL2025.robinFunctionValue} multiplied by the symbolic reciprocal
\texttt{N\_f\_inv}.  The declarations
\texttt{GHL2025.FunctionOraclePaperImage} and
\texttt{GHL2025.functionOraclePaperImage} record the clean branch,
orthogonal-component label, and false obligations for normalized amplitude,
orthogonality, the \(N_f\) bound, unitary completion, and diagonal-helper
isolation.

The declaration \texttt{GHL2025.functionOraclePaperMatrix} exposes this clean
branch at matrix level for clean \(m_f\) input columns: it places
\texttt{GHL2025.functionOracleNormalizedValue} on the clean \(m_f\) branch row,
sets other clean-workspace output rows to zero, and uses
\texttt{GHL2025.functionOracleOrthogonalEntry} for non-clean output rows.  For
non-clean input columns, all entries remain symbolic.  The active gate
\texttt{GHL2025.oneTermRobinGate\_O\_f} now uses this matrix.  This skeleton
does not prove the amplitude relation, the \(N_f\) bound, orthogonality of the
non-clean rows, or unitarity.  The Lean records that must stay false are
\[
  \texttt{GHL2025.oneTermRobinGate\_O\_f.unitary.proved},
  \qquad
  \texttt{GHL2025.FunctionOracleContract.amplitudeCorrect.proved},
  \qquad
  \texttt{GHL2025.RobinProofObligations.functionOracleCorrect.proved}.
\]

\begin{center}
\begin{tabular}{p{0.22\linewidth}p{0.30\linewidth}p{0.35\linewidth}}
\hline
Contract item & Lean declaration & Status \\
\hline
grid register \(i\) &
\texttt{functionOraclePaperRegisters};
\texttt{functionOracleMatrix} uses the same system bits &
paper-register extractor defined; concrete tests pass \\
function values \(f(x_i)\) &
\texttt{GHL2025.robinFunctionValue} &
symbolic data source defined \\
clean \(m_f\) workspace &
\texttt{FunctionOraclePaperRegisters.mfWorkspaceValue} and
\texttt{cleanWorkspace} &
workspace field defined; action and cleanup unproved \\
normalizer \(N_f\) &
\texttt{FunctionOracleContract.normalizerBound} and
\texttt{GHL2025.oneTermRobinNormalizer} &
recorded; analytic bound unproved \\
amplitude relation \(f(x_i)/N_f\) &
\texttt{functionOracleNormalizedValue} and
\texttt{FunctionOracleContract.amplitudeCorrect} &
symbolic normalized-value contract defined; proof flag false \\
orthogonal component &
\texttt{FunctionOraclePaperImage.orthogonalComponent},
\texttt{functionOracleOrthogonalEntry}, and
\texttt{orthogonalComponentCorrect} &
component label and symbolic non-clean rows recorded; zero-overlap proof false \\
\(m_f\) workspace action and cleanup &
\texttt{FunctionOraclePaperImage.cleanWorkspaceBranch} and future workspace lemma &
width recorded; action unproved \\
active gate skeleton &
\texttt{GHL2025.oneTermRobinGate\_O\_f} &
uses \texttt{functionOraclePaperMatrix}; unitarity false \\
\hline
\end{tabular}
\end{center}

The completed proof-DAG block \texttt{of\_paper\_image} introduced the
paper-register extractor, normalized clean-branch value, and paper-image
record.  The companion block \texttt{of\_paper\_matrix} introduced
\texttt{functionOraclePaperMatrix} and rewired the active gate to it.  The
bridge theorem
\texttt{functionOraclePaperMatrix\_nonCleanInput\_entry} records that a
non-clean input column stays in the symbolic completion branch even when the
output row equals the clean-branch basis row for the same non-\(m_f\) index.
The
block \texttt{of\_diagonal\_helper\_isolation} now keeps
\texttt{functionOracleMatrix} as a helper-only data path and prevents a
unitarity proof for that diagonal helper from being mistaken for Lemma 4.

The current tests pin the active gate matrix to
\texttt{GHL2025.functionOraclePaperMatrix}, pin the Robin default normalizer bound to
\texttt{Coeff.symbol "N\_f"}, pin the general diagonal rule
\texttt{functionOracleMatrix p j j =
robinFunctionValue p.n sysVal} for the extracted system index, and pin
\texttt{amplitudeCorrect.proved = false}.  New tests also pin
\texttt{functionOraclePaperRegisters} on a clean \(n=3,\kappa=7\) column,
the nonclean \(m_f\) extraction example, the normalized value
\texttt{robinFunctionValue 3 2} times \texttt{N\_f\_inv}, the clean-branch
fields of \texttt{functionOraclePaperImage}, the clean-input branch matrix
entry \(\texttt{functionOraclePaperMatrix}\;36\;36\), the zero entry for
another clean-workspace output row, a symbolic non-clean output-row completion
entry, a symbolic non-clean input-column entry at row \(4\), and all five false
proof flags.  These are contract tests, not proofs of the amplitude-oracle
relation.

\subsection{Sparse Amplitude Oracle $O_{D^T}^S$ Contract}

The sparse amplitude oracle $O_{D^T}^S$ is governed by Lemma 3 of
Guseynov--Huang--Liu 2025, arXiv:2506.20478.  The paper oracle uses the sparse
coefficient $D^{(s)}/N_D$ as the amplitude of the $\ket{0}$ branch:
\[
  \hat O_D^S\ket{0}\ket{s}
  =
  \frac{D^{(s)}}{N_D}\ket{0}\ket{s}
  +
  \sqrt{1-\frac{|D^{(s)}|^2}{N_D^2}}\ket{1}\ket{s}.
\]
For the one-term Robin circuit the operator is used for $D^T$, and the
row-dependent coefficient source in Lean is
\(\texttt{GHL2025.robinSparseAmplitudeValue}\).  The registers needed by the
Lean contract are:
\begin{itemize}
  \item the ancilla bit at bit 0;
  \item the row register in bits $[1,1+n)$;
  \item the sparse index inside the padded $O_D$ register block;
  \item the indicator bit at \texttt{GHL2025.robinIndicatorBitPosition}.
\end{itemize}
The indicator bit is set by $U_{\mathrm{indic}}$ before $O_{D^T}^S$.  The
Lean-facing skeleton acts as identity when that bit is 0 and applies a symbolic
two-by-two rotation on ancilla bit 0 when that bit is 1.  The symbols
\(\texttt{GHL2025.sparseAmplitudeOracleDTCosHalf}\) and
\(\texttt{GHL2025.sparseAmplitudeOracleDTSinHalf}\) are placeholders for the
two entries in that block; their relation to the displayed formula is not yet
proved.

The Lean declaration \(\texttt{GHL2025.sparseAmplitudeOracleDTMatrix}\) remains
a diagonal helper:
\[
  M[j][j] =
  \begin{cases}
    1, & \text{indicator bit } = 0,\\
    \texttt{robinSparseAmplitudeValue}(n,s,i), & \text{indicator bit } = 1,
  \end{cases}
  \qquad
  M[i][j] = 0 \quad (i \ne j).
\]
This helper exercises the sparse-amplitude data path, but it is not the
faithful active oracle.  The active gate
\(\texttt{GHL2025.oneTermRobinGate\_O\_DT\_S}\) now points to
\(\texttt{GHL2025.sparseAmplitudeOracleDTRotationMatrix}\), which preserves all
non-ancilla bits and changes only ancilla bit 0.

The lower packet added:
\[
  \texttt{GHL2025.SparseAmplitudeOracleDTPaperRegisters},
  \qquad
  \texttt{GHL2025.sparseAmplitudeOracleDTPaperRegisters},
  \qquad
  \texttt{GHL2025.sparseAmplitudeOracleDTCoefficientNormalizerObligation},
  \qquad
  \texttt{GHL2025.sparseAmplitudeOracleDTCoefficientNormalizerContract},
  \qquad
  \texttt{GHL2025.sparseAmplitudeOracleDTRotationMatrix}.
\]
For \(n=3\) and \(\kappa=7\), column 132 has indicator bit 1, row value 2,
and sparse index 0.  The focused tests check the symbolic
entries \(M[132,132]=\texttt{odts\_cos\_half\_2\_0}\) and
\(M[133,132]=\texttt{odts\_sin\_half\_2\_0}\).  Column 4 has indicator bit 0,
so the tests check identity on row 4 and no ancilla flip to
row 5.

The explicit coefficient-normalizer obligation is:
\[
  \texttt{GHL2025.sparseAmplitudeOracleDTCoefficientNormalizerObligation.proved}
  =
  \texttt{false}.
\]
This obligation requires a future Lean statement proving that the symbolic
\(\ket{0}\) entry represents \(D^{(s)}/N_D\), that the complementary entry
matches the square-root normalizer term above, and that the resulting
two-by-two block is unitary.  It must not be encoded as
\(\texttt{proved := true}\) until those facts are formalized.

The per-row contract
\[
  \texttt{GHL2025.sparseAmplitudeOracleDTCoefficientNormalizerContract}
\]
records the same obligation with the row and sparse index exposed.  For
parameters \(p\), row \(j\), and sparse index \(s\), its coefficient field is
\[
  \texttt{GHL2025.robinSparseAmplitudeValue}\ p.n\ s\ j,
\]
its normalizer field is \(\texttt{Coeff.symbol "N\_D"}\), and its two symbolic
entry fields are
\(\texttt{GHL2025.sparseAmplitudeOracleDTCosHalf}\ j\ s\) and
\(\texttt{GHL2025.sparseAmplitudeOracleDTSinHalf}\ j\ s\).  The contract has
three false proof flags: the \(\ket{0}\) coefficient relation, the
complementary square-root relation, and the two-by-two unitarity relation.  For
\(n=3\), \(\kappa=7\), \(j=2\), and \(s=0\), Lean tests pin the coefficient to
\(-1/12\), the normalizer to \(N_D\), and the symbolic entries to
\(\texttt{odts\_cos\_half\_2\_0}\) and
\(\texttt{odts\_sin\_half\_2\_0}\).

The refined proof-route record
\[
  \texttt{GHL2025.sparseAmplitudeOracleDTCoefficientNormalizerProofRoute}
\]
keeps the same Eq.~(20) contract but exposes the next proof blocks explicitly.
It uses
\[
  \texttt{GHL2025.sparseAmplitudeOracleDTNormalizedCoefficient}
\]
as the symbolic stand-in for \(D_j^{(s)}/N_D\), implemented as the Robin sparse
coefficient multiplied by \(\texttt{Coeff.symbol "N\_D\_inv"}\).  The route
has false obligations for interpreting the division by \(N_D\), proving the
\(N_D\) bound, interpreting the absolute-square term
\(|D_j^{(s)}|^2/N_D^2\), proving the square-root complement, and proving the
two-by-two unitary relation.  These fields are proof-route memory only; they
do not promote \(\texttt{GHL2025.oneTermRobinGate\_O\_DT\_S.unitary.proved}\).
The Lean route also has definitional bridge lemmas showing that its normalizer
and symbolic \(\ket{0}\), \(\ket{1}\) rotation-entry fields are the same fields
as the per-row Eq.~(20) contract.

In Lean:
\[
  \texttt{GHL2025.sparseAmplitudeOracleDTMatrix},
  \qquad
  \texttt{GHL2025.oneTermRobinGate\_O\_DT\_S},
  \qquad
  \texttt{GHL2025.sparseAmplitudeOracleDTRotationMatrix}.
\]

\subsection{Boundary Rotation $R_y^{\mathrm{boundary}}$ Contract}

The boundary rotation gate in Fig. 1-term Robin applies a controlled
\(R_y(\theta_j^s)\) block on ancilla bit 0 for boundary rows.  The paper angle
relation is
\[
  \theta_j^s = \arccos(D_j^{(s)}/N_D).
\]
The half-angle entries used by the two-by-two block are
\[
  \cos(\theta_j^s/2)
  =
  \sqrt{\frac{1 + D_j^{(s)}/N_D}{2}},
  \qquad
  \sin(\theta_j^s/2)
  =
  \sqrt{\frac{1 - D_j^{(s)}/N_D}{2}}.
\]

The Lean register contract extracts the ancilla bit, indicator bit, row value,
sparse index, and non-ancilla rest value:
\[
  \texttt{GHL2025.BoundaryRotationPaperRegisters},
  \qquad
  \texttt{GHL2025.boundaryRotationPaperRegisters}.
\]
The indicator bit controls the branch: indicator value \(1\) means the matrix
acts as identity, and indicator value \(0\) means the boundary rotation is
active.  The coefficient source for \(D_j^{(s)}\) is
\[
  \texttt{GHL2025.robinSparseAmplitudeValue}\ p.n\ s\ j.
\]

The active matrix uses shared symbolic half-angle names:
\[
  \texttt{GHL2025.boundaryRotationCosHalf}\ j\ s,
  \qquad
  \texttt{GHL2025.boundaryRotationSinHalf}\ j\ s.
\]
These symbols do not prove the arccos relation, the half-angle formulas, or
unitarity.  The per-row contract
\[
  \texttt{GHL2025.boundaryRotationAngleNormalizerContract}\ p\ j\ s
\]
records the coefficient, the normalizer symbol \(N_D\), the two symbolic
entries, and five false proof flags: boundary control, arccos argument,
cosine half-angle, sine half-angle, and two-by-two unitarity.  The global flag
\[
  \texttt{GHL2025.boundaryRotationAngleNormalizerObligation.proved}
  =
  \texttt{false}
\]
keeps the paper-level gap visible.

For \(n=3\), \(\kappa=7\), row \(j=0\), and sparse index \(s=0\), compiled
tests pin the coefficient to \(-5/2+(7/3)A_1dx\), the normalizer to \(N_D\),
and the symbolic entries to \(\texttt{boundary\_cos\_half\_0\_0}\) and
\(\texttt{boundary\_sin\_half\_0\_0}\).  The same tests confirm that all
contract proof fields remain false and that
\(\texttt{oneTermRobinGate\_Ry\_boundary.matrix}\) is definitionally
\(\texttt{boundaryRotationMatrix}\).

The Lean-facing proof DAG is:
\begin{center}
\begin{tabular}{p{0.22\textwidth}p{0.42\textwidth}p{0.22\textwidth}}
\hline
Block & Interface & Status \\
\hline
\texttt{ryb\_extract\_registers} &
extract ancilla, indicator, row, sparse index, and non-ancilla fields &
defined skeleton \\
\texttt{ryb\_symbolic\_entries} &
define the symbolic half-angle entries in the active matrix &
active skeleton; angle relation unproved \\
\texttt{ryb\_angle\_normalizer} &
prove \(\theta_j^s=\arccos(D_j^{(s)}/N_D)\) and the half-angle formulas &
typed contract recorded; proof flags false \\
\texttt{ryb\_rotation\_unitary} &
prove the two-by-two block is unitary under the \(N_D\) bound &
unproved \\
\hline
\end{tabular}
\end{center}

\subsection{Shared \(N_D\) Normalizer Contract}

Both the Lemma~3 oracle \(O_{D^T}^S\) and the boundary rotation
\(R_y^{\mathrm{boundary}}\) use the normalized derivative coefficient
\[
  D_j^{(s)}/N_D.
\]
Lean now records this common dependency in
\[
  \texttt{GHL2025.DerivativeNormalizerNDContract},
  \qquad
  \texttt{GHL2025.derivativeNormalizerNDContract}\ p\ j\ s.
\]
The contract stores the coefficient
\(\texttt{GHL2025.robinSparseAmplitudeValue}\ p.n\ s\ j\), the normalizer
symbol \(\texttt{Coeff.symbol "N\_D"}\), and the formal normalized coefficient
obtained by multiplying by \(\texttt{Coeff.symbol "N\_D\_inv"}\).

This is a shared proof obligation, not an analytic proof.  Its fields for
nonzero \(N_D\), division semantics, coefficient bound, absolute-square
semantics, square-root complement, arccos semantics, and two-by-two unitarity
all have \(\texttt{proved := false}\).  The bridge theorem
\[
  \texttt{GHL2025.sparseAmplitudeOracleDTCoefficientNormalizerProofRoute\_sharedND}
\]
states that the \(O_{D^T}^S\) proof route uses the shared division, bound,
absolute-square, and square-root fields.  The bridge theorem
\[
  \texttt{GHL2025.boundaryRotationAngleNormalizerProofRoute\_sharedND}
\]
states that the boundary-rotation proof route uses the shared division,
arccos-domain, and bound fields.  Neither bridge promotes a gate-unitarity
flag.

Lean also records a source-bound view
\[
  \texttt{GHL2025.derivativeNormalizerNDSourceBound}\ p\ j\ s
\]
for the future inequality
\[
  |D_j^{(s)}| \le N_D .
\]
This view uses the same source coefficient
\(\texttt{GHL2025.robinSparseAmplitudeValue}\ p.n\ s\ j\) and normalizer
\(\texttt{Coeff.symbol "N\_D"}\).  The theorem
\[
  \texttt{GHL2025.derivativeNormalizerNDSourceBound\_coefficientBound}
\]
states that the view reuses the shared coefficient-bound obligation, and
\[
  \texttt{GHL2025.derivativeNormalizerNDSourceBound\_coefficientBound\_false}
\]
states that the proof flag is still false.

\paragraph{Cycle 17 source-contract audit.}
The \(O_D^{BS}\) cleanup route is currently blocked by the missing
unused-branch image rule, so the next faithful Phase~1 packet is the shared
\(N_D\) normalizer transcript for \(O_{D^T}^S\) and
\(R_y^{\mathrm{boundary}}\).  The fixed Lean contract is:
\[
  D_j^{(s)}
  =
  \texttt{GHL2025.robinSparseAmplitudeValue}\ p.n\ s\ j,
  \qquad
  N_D = \texttt{Coeff.symbol "N\_D"} .
\]
The formal quotient is represented by
\[
  \texttt{Coeff.mul}\,
    (\texttt{GHL2025.robinSparseAmplitudeValue}\ p.n\ s\ j)\,
    (\texttt{Coeff.symbol "N\_D\_inv"}) ,
\]
inside \texttt{GHL2025.derivativeNormalizerNDContract}.  This is only a
contract placeholder.  The obligations for nonzero \(N_D\), division
semantics, the coefficient bound, absolute-square semantics, the
square-root complement, arccos semantics, and two-by-two unitarity remain
\(\texttt{proved := false}\).

The next lower packet should add only missing definitional bridges or focused
tests showing that the \(O_{D^T}^S\) and \(R_y^{\mathrm{boundary}}\) routes
reuse this shared contract and the source-bound view.  It should not prove
analytic inequalities, change a gate matrix, promote a unitarity flag, or
resume \(O_D^{BS}\) cleanup proof search.

Cycle~17 lower adds exactly these source-bound route bridges:
\[
  \texttt{GHL2025.sparseAmplitudeOracleDTCoefficientNormalizerProofRoute\_sourceBound},
\]
\[
  \texttt{GHL2025.boundaryRotationAngleNormalizerProofRoute\_sourceBound},
  \qquad
  \texttt{GHL2025.derivativeNormalizerNDSourceBound\_sharedRoutes}.
\]
They state that the two proof-route records reuse the same source coefficient,
the same \(N_D\) symbol, and the same coefficient-bound obligation.  These are
definitional bridges only; the analytic and unitary proof flags remain
\(\texttt{proved := false}\).

\paragraph{Current shared \(N_D\) guard.}
Middle rechecked Lemma~3, Eq.~(20), the boundary \(R_y\) angle equation, and
Fig.~1-term Robin before assigning more proof search.  Lean now records the
guard
\[
  \texttt{GHL2025.derivativeNormalizerNDSharedRoute\_flags\_false}.
\]
This theorem packages the current fixed route state: nonzero \(N_D\), division
semantics, the coefficient bound, absolute-square semantics, square-root
semantics, arccos semantics, half-angle semantics, and two-by-two unitarity
remain false where they occur.  It also keeps
\[
  \texttt{GHL2025.oneTermRobinGate\_O\_DT\_S.unitary.proved}
  =
  \texttt{false},
  \qquad
  \texttt{GHL2025.oneTermRobinGate\_Ry\_boundary.unitary.proved}
  =
  \texttt{false}.
\]
This is proof-state synchronization, not an analytic theorem.

\paragraph{Middle closeout for the packaged guard.}
After the lower guard
\[
  \texttt{GHL2025.derivativeNormalizerNDSharedRoute\_sourceBoundAndFlags}
\]
compiled, middle re-audited Lemma~3, Eq.~(20), Eq.~\texttt{angles for Ry},
and Fig.~1-term Robin.  The source contract is unchanged: both
\(O_{D^T}^{S}\) and \(R_y^{\mathrm{boundary}}\) use the same coefficient
source \(D_j^{(s)}\), the same normalizer \(N_D\), and the same bound
\(N_D \geq \|D\|_{\max}\).  Lean records this only as typed route
synchronization through \texttt{DerivativeNormalizerNDContract} and
\texttt{DerivativeNormalizerNDSourceBound}.  Nonzero division, coefficient
bounds, absolute-square semantics, square-root semantics, arccos semantics,
half-angle identities, and two-by-two unitarity remain proof obligations.
Neither affected gate-unitarity flag, nor any LCU or block-extraction flag, is
promoted by this guard.

\paragraph{Cycle 18 middle packet, superseded \(O_D^{BS}\) status.}
At this point the \(O_D^{BS}\) route was still recorded through the unused
zero-amplitude sparse-branch source decision.  That status is historical.  The
active route after the global sparse-slot correction uses
\texttt{bandedSparseAccessPaperGlobalSlotSource} and the cleanup-scope decision
record; the old row-dependent decision remains rejected-model memory only.

The shared \(N_D\) bridges for \(O_{D^T}^S\) and
\(R_y^{\mathrm{boundary}}\) were already recorded.  The next lower packet at
that point therefore stayed on the function-oracle transcript and added only
an \(O_f\) amplitude-route contract.  The existing Lean anchors were
\[
  \texttt{GHL2025.robinFunctionValue}\ p.n\ i,\qquad
  \texttt{GHL2025.functionOracleNormalizedValue}\ p\ i,
\]
\[
  \texttt{GHL2025.functionOraclePaperImage}\ p\ j,\qquad
  \texttt{GHL2025.functionOraclePaperMatrix}\ p .
\]
The planned route names are
\texttt{GHL2025.FunctionOracleAmplitudeProofRoute} and
\texttt{GHL2025.functionOracleAmplitudeProofRoute}.  They should reuse the
paper normalizer
\[
  \texttt{(Examples.RobinHeat.robinOracleComposition n).functionOracle.normalizerBound}
  = \texttt{Coeff.symbol "N\_f"}
\]
and should keep \(N_f\) nonzero, division semantics, the \(N_f\) bound,
orthogonality of the completion, unitary completion, and
\texttt{FunctionOracleContract.amplitudeCorrect} as false obligations.

\paragraph{Cycle 18 lower packet.}
Lean now records the cited amplitude-oracle theorem as
\texttt{GHL2025.FunctionOracleExternalAmplitudeSourceContract} with default
\texttt{GHL2025.functionOracleExternalAmplitudeSourceContract}.  This transcript
names GHL2025 Theorem ``Amplitude-oracle for piece-wise polynomial function'',
Eq.~``coordinate oracle'', and the cited source arXiv:2411.01131.  The guard
\texttt{GHL2025.functionOracleExternalAmplitudeSourceContract\_flags\_false}
keeps the resource claim, external theorem formalization, nonzero \(N_f\),
division semantics, theorem-amplitude correctness, and closure booleans false.
The bridge
\texttt{GHL2025.functionOracleAmplitudeProofRoute\_externalSourceContract}
ties the per-column \(O_f\) amplitude route to this transcript without changing
\texttt{GHL2025.functionOraclePaperMatrix} or promoting any analytic proof flag.

\subsection{$O_f$ vs $O_{D^T}^S$ Double Encoding Audit (Resolved)}

Both \texttt{functionOracleMatrix} ($O_f$) and
\texttt{sparseAmplitudeOracleDTMatrix} ($O_{D^T}^S$) previously used
\texttt{robinSparseAmplitudeValue} as their diagonal data source,
creating a double encoding of derivative amplitude data.
In the paper:
\begin{itemize}
  \item $O_{D^T}^S$ encodes derivative stencil coefficients $D_j^{(s)}/N_D$
    via controlled rotations on the ancilla qubit.
  \item $O_f$ encodes function values $f(x_j)/N_f$
    (Guseynov--Huang--Liu 2025, Lemma 4, arXiv:2506.20478),
    a completely different quantity.
\end{itemize}
This was corrected in cycle 9: $O_f$ now uses \texttt{robinFunctionValue}
(symbolic function values $f(x_j)$), while the $O_{D^T}^S$ data path continues
to use \texttt{robinSparseAmplitudeValue} (derivative stencil data).  The
active $O_{D^T}^S$ gate now uses the controlled-rotation skeleton, not the
diagonal helper.  Both gates carry $\texttt{unitary.proved := false}$; for
$O_{D^T}^S$ the remaining gap is the Lemma 3 coefficient-normalizer relation
and the induced two-by-two unitarity identity.

\subsection{Open Proof Obligations}

The current unproved mathematical obligations are intentionally data, not fake
proofs:
\[
  \texttt{proved := false}.
\]
The most important remaining obligations are:

\begin{enumerate}
  \item prove the remaining five labeled gate matrices are unitary
    (U\_indic and SWAP are already proved permutation matrices);
  \item prove the composed circuit matrix equals the paper's \(U\);
  \item replace the Ry\_boundary zero placeholder with a real oracle implementation; (done: cycle 8)
  \item correct $O_f$ to encode function values $f(x_j)/N_f$ instead of derivative amplitude data; (done: cycle 9)
  \item prove the $O_{D^T}^S$ coefficient-normalizer relation for the symbolic rotation entries from Lemma 3;
  \item prove the $R_y^{\text{boundary}}$ angle-normalizer relation, including the arccos argument, half-angle formulas, and two-by-two unitarity;
  \item prove the extracted block is \(A_k/(N_DN_f\kappa)\);
  \item prove the resource formula including all internal oracle ancillas.
\end{enumerate}

All seven gates in the one-term Robin circuit now have nonzero matrix
semantics in Lean: U\_indic, \(O_{D^T}^S\), \(R_y^{\text{boundary}}\),
\(O_D^{BS}\), \(O_f\), SWAP, and \((O_D^{BS})^\dagger\).
U\_indic unitarity is now proved via \texttt{indicatorOracleMatrix\_is\_permutation}.
SWAP unitarity is proved via \texttt{swapOracleMatrix\_is\_permutation}.
The remaining five gates are still proof-obligation carriers:
their \texttt{unitary.proved} fields remain false, and the exact paper
rotation/block-extraction theorem is still open.

\subsection{Cycle 10: Circuit Product and Block Extraction Pipeline}

The block-extraction target is no longer a zero placeholder.  In Lean,
\texttt{oneTermRobinCircuitSemantics} defines the full-space matrix as
\[
  \texttt{evalGateMatrices (oneTermRobinGateMatrixPlaceholders p)}.
\]
The target \texttt{oneTermRobinBlockExtractionTarget} then obtains
\texttt{unitaryMatrix} from this product and obtains \texttt{blockMatrix}
by applying \texttt{signalSystemBlockProjection} to the signal-index-zero
block.

The tests intentionally check this pipeline structurally rather than asking
Lean to normalize concrete entries of the full \(13\)-qubit matrix for \(n=3\).
That concrete product has dimension \(8192 \times 8192\), and entry-level
\texttt{native\_decide} checks are too expensive for routine builds.  Concrete
entry correctness for the product is therefore part of the remaining block
correctness proof obligation, while the compiled tests verify that the Lean
objects are wired to the real gate product and projection API.

The generic arithmetic bridge
\[
  \texttt{clog2 (gridSize n) = n}
\]
is now proved in Lean as \texttt{clog2\_gridSize}.  This closes the dimension
compatibility side condition for the Robin circuit block claim, but it does not
prove oracle unitarity or block correctness.

\subsection{Block-Projection and Normalizer Contract}

Define the signal dimension to be
\[
  2^{m}
  =
  \texttt{qubitDim}
  \bigl(\texttt{GHL2025.effectiveRobinSignalQubits}
  (\texttt{Examples.RobinHeat.oneTermParameters}\;n)\bigr).
\]
Define the system dimension to be \(N=2^n=\texttt{gridSize}\;n\).  Lean proves
the dimension split by
\[
  \texttt{Examples.RobinHeat.oneTermRobinCircuitDimCompat}\;n.
\]
The block-extraction target uses signal index \(0\).  Its target matrix is
\texttt{Examples.RobinHeat.oneTermRobinAkMatrix}\(n\), and its normalizer is
\[
  \texttt{GHL2025.oneTermRobinNormalizer}
  =
  N_D N_f \kappa.
\]
The evaluation lemma
\(\texttt{GHL2025.oneTermRobinNormalizer\_eval}\) states this product for every
symbol environment.

The lower packet audited only this structural contract:
\texttt{oneTermRobinBlockExtractionTarget} uses the cast circuit product,
\texttt{signalSystemBlockProjection}, signal index \(0\), the row-scaled
Robin target matrix, and the normalizer above.  The compiled tests also pin
\texttt{defaultOneTermRobinCircuitBlockClaim} to that same target.  The
obligations \texttt{blockProjection.proved} and
\texttt{blockCorrect.proved} remain false.  The cleanup obligation for
\((O_D^{BS})^\dagger\) also remains false through
\texttt{defaultBandedSparseAccessPaperContract}.  This packet did not attempt
entry-level normalization of the full \(n=3\) product, whose dimension is
\(8192 \times 8192\).

\subsection{Theorem-Level Proof-Route Contract}

Define the Phase 1 route
\[
  \texttt{Examples.RobinHeat.oneTermRobinBlockEncodingProofRoute}\;n.
\]
This Lean object packages the theorem tuple
\texttt{GHL2025.defaultOneTermRobinTheoremData}, the circuit product
\texttt{oneTermRobinCircuitSemantics}, the signal-index-zero target
\texttt{oneTermRobinBlockExtractionTarget}, the circuit block claim, the
Robin oracle composition record, the active \(O_D^{BS}\) cleanup-scope
decision, and the cited-source blocker for \(O_f\).

The route records the normalizer agreement
\[
  \texttt{theoremData.alpha}
  =
  \texttt{blockClaim.target.normalizer}
  =
  \texttt{GHL2025.oneTermRobinNormalizer},
\]
which is the symbolic coefficient \(N_D N_f \kappa\).  It also records that
the target matrix is \texttt{oneTermRobinAkMatrix}\(n\) and that the block
projection uses signal index \(0\).

The layout audit separates two counts.  The theorem-level block-encoding tuple
uses
\[
  m_{\mathrm{theorem}}
  =
  \lceil\log_2 n\rceil+\lceil\log_2 G_f\rceil
  +\lceil\log_2\kappa\rceil+4
\]
signal qubits and \(2n\) pure ancillas, as recorded by
\texttt{defaultOneTermRobinTheoremData}.  The concrete matrix projection uses
\texttt{effectiveRobinSignalQubits}, because the projection zeros every
non-system wire in the register partition.  Lean proves
\[
  \texttt{effectiveRobinSignalQubits}
  =
  \texttt{oneTermRobinLayout.signalQubits}
  + \texttt{odPureAncillaQubits} + 1.
\]
This bridge is recorded by
\texttt{effectiveRobinSignalQubits\_eq\_layout\_signal\_plus\_visibleWorkspace}
and the route guard
\texttt{oneTermRobinBlockEncodingProofRoute\_layoutProjectionAudit}.  It is not
a proof of pure-ancilla cleanup or block correctness; the corresponding flags
remain false.

The guard theorem
\[
  \texttt{Examples.RobinHeat.oneTermRobinBlockEncodingProofRoute\_flags\_false}
\]
keeps the theorem in obligation mode.  In particular, the block-projection and
block-correctness obligations remain false, the full circuit-unitarity
obligation remains false, \(O_D^{BS}\) cleanup and unitary extension remain
false, the cleanup scope is restricted to
\texttt{bandedSparseAccessPaperGlobalSlotSource}, full clean-domain cleanup
and full-space unitarity are not selected, the \(O_f\) external theorem remains
an obligation, and the LCU composition field remains false.  This route is a
synchronization contract for future proof work, not a proof of the
block-encoding equation.

\subsection{Historical Source Decision for the Row-Dependent Helper}

The historical row-dependent sparse-branch blocker is the cited-results obligation
\texttt{QBE.ODBS.UnusedZeroBranchExtension}.  It covers clean padded
\(O_D^{BS}\) inputs whose sparse index is included in the Robin
\(\kappa\)-wide enumeration but whose amplitude is zero for the row-dependent
stencil.  This subsection is retained as rejected-model memory; it is no
longer the active lower target for \(O_D^{BS}\).

The current source transcript has three parts.  Lemma 1 of GHL2025 records the
padded sparse-access map
\[
  O_D^{BS}|0\rangle^{n-l}|s\rangle^l|i\rangle^n
  =
  |r_{si}\rangle^n|i\rangle^n.
\]
The one-term Robin text records that zeros may be included in the sparse
enumeration, and Eq.~\texttt{ROBIN clarified} uses
\(s=0,\ldots,\kappa-1\).  The figure caption asks
\((O_D^{BS})^\dagger\) to restore the padded sparse-index register after SWAP.
None of these source fragments gives an injective reversible image for clean
unused zero-amplitude branches.

Older Lean cycles recorded this source decision through
\texttt{bandedSparseAccessUnusedZeroBranchSourceDecision} and
\texttt{bandedSparseAccessRobinZeroInclusionSourceContract}.  Those records
remain available as separate transcript memory, but the theorem route no
longer stores them as active fields.  The current route records the active
cleanup scope through \texttt{bandedSparseAccessCleanupScopeDecision}; its
semantic-cleanup promotion, \(O_D^{BS}\) dagger-cleanup, unitary-extension,
LCU, and block-correctness flags remain false.

The active route below supersedes this guard for new lower work.  New
\(O_D^{BS}\) packets should use
\texttt{bandedSparseAccessPaperGlobalSlotSource} and the global address
\(r_{si}=r_{s0}+i \bmod 2^n\), not the row-dependent valid-branch predicate.

The subsequent source-decision subsections in this historical cycle are
therefore superseded unless they explicitly cite the active global-slot
cleanup scope.  They should be read as rejected-model or transcript memory,
not as lower-agent packets for the current active route.

The active \(O_D^{BS}\) address correction now separates this historical
row-dependent source-decision guard from the paper address used by the matrix
skeleton.  Lean defines \texttt{oneTermRobinGlobalSparseOffset} for the seven
one-term Robin slots and
\texttt{oneTermRobinGlobalSparseAddress} for
\(r_{si}=r_{s0}+i \bmod 2^n\).  The active
\texttt{bandedSparseAccessPaperAddress} uses this global-slot address, while
\texttt{bandedSparseAccessRowDependentPaperAddress} and
\texttt{bandedSparseAccessRowDependentPaperImage} preserve the old helper as a
rejected-model regression.  The theorem
\texttt{oneTermRobinGate\_O\_D\_BS\_globalSparseBoundaryNoCollision\_n3}
checks that the old \(n=3\) boundary columns \(0\) and \(48\) no longer collide
in the corrected active image; no forward-correctness, cleanup, unitarity, LCU,
or block-correctness flag is promoted.

\subsection{Cycle 1 Middle Source-Dependency Audit}

The cycle-1 middle audit did not find a new source-backed image rule for clean
unused zero-amplitude sparse branches.  The public anchors remain Lemma
\texttt{Banded-sparse-access}, Theorem \texttt{1 term robin}, Eq.\
\texttt{ROBIN clarified}, Fig.\ \texttt{1 term ROBIN}, and the cited prior PDE
sparse-access primitive.

The proof map is therefore unchanged.  Lemma 1 supplies the padded
sparse-access equation, the Robin text keeps zeros in the sparse range
\(s=0,\ldots,\kappa-1\), and the figure asks the dagger to restore the padded
register.  None of those source fragments supplies an injective reversible
image for clean unused branches.

The Lean route records this as a guard contract through
\texttt{bandedSparseAccessUnusedZeroBranchSourceDecision},
\texttt{bandedSparseAccessRobinZeroInclusionSourceContract}, and
\texttt{oneTermRobinBlockEncodingProofRoute\_odbsNoLowerProofSearch}.  The
next lower packet is restricted to source-transcript or false-flag guard work
unless a source supplies an exact unused-branch image formula or a named
reversible-extension theorem.

\subsection{Cycle 1 Middle Final Handoff}

The final middle handoff for cycle 1 keeps the same source-dependency
classification.  Lemma~\texttt{Banded-sparse-access} supplies the padded
register equation for \(O_D^{BS}\), Eq.~\texttt{ROBIN clarified} keeps
\(s=0,\ldots,\kappa-1\) in the Robin transcript, and Fig.~\texttt{1 term ROBIN}
requires the dagger cleanup after SWAP.  The inspected source still does not
give an injective reversible image for clean unused zero-amplitude sparse
branches.

The historical Lean proof route was therefore a guard route.  Older
declarations such as
\texttt{oneTermRobinBlockEncodingProofRoute\_odbsPaperContractTranscript},
\texttt{oneTermRobinBlockEncodingProofRoute\_sourceDecisionImageSlotsBlocked},
\texttt{oneTermRobinBlockEncodingProofRoute\_activeOdbsGatePairWiring},
\texttt{oneTermRobinBlockEncodingProofRoute\_activeOdbsGatePairPublicSources},
\texttt{oneTermRobinBlockEncodingProofRoute\_blockProjectionNormalizerAudit},
and
\texttt{oneTermRobinBlockEncodingProofRoute\_sourceBlockerKeepsFinalFlagsFalse}
belong to that superseded guard route.  Current lower work should use the
active global-source declarations recorded in the 2026-05-25 sync instead.
The historical guards did not prove injectivity, dagger cleanup, unitarity, LCU
composition, or the final block-extraction equation.

\subsection{Cycle 1 Post-Lower Source-Gate Freeze (Historical)}

The historical lower guard
\texttt{oneTermRobinBlockEncodingProofRoute\_sourceGateFreeze} packages the
same source decision without adding a new source theorem.  This route has been
retired from the active theorem tuple.  The current replacement is the
cleanup-scope route through
\texttt{oneTermRobinBlockEncodingProofRoute\_odbsActiveGlobalSlotBlockers} and
\texttt{oneTermRobinBlockEncodingProofRoute\_odbsActiveGlobalSlotGateFreeze}.

Two focused tests make the blocker visible at the matrix-entry level.  Source
column \(8\) checks that the active paper-image matrix follows Lemma~1 and has
entry \((40,8)=1\), while the legacy helper still has its historical
\((4,8)=1\) entry.  Source columns \(0\) and \(48\) now check the corrected
active global-slot image: the old row-dependent helper collides on these
columns, while the active paper-image skeleton separates them.  These tests
separate the active paper contract from the legacy helper and preserve the
remaining proof obligations; they do not prove injectivity, cleanup, unitarity,
LCU composition, or block extraction.

\subsection{Cycle 1 Route-Collision Guard}

The route guard
\texttt{oneTermRobinBlockEncodingProofRoute\_activeCollisionBlocked\_n3}
ties the historical row-dependent boundary collision to the theorem-level
false flags while also recording that the corrected active global-slot image
separates the old boundary columns \(0\) and \(48\).  It keeps the
dagger-cleanup, unitary-extension, LCU, and final block-extraction flags false.

This declaration is a guard, not a source theorem.  It may be used only as
rejected-model memory for the old row-dependent helper.

\subsection{Cycle 1 Contract-Drift and Projection Freeze}

The route guard
\texttt{oneTermRobinBlockEncodingProofRoute\_contractDriftColumn8Blocked\_n3}
keeps the active \(O_D^{BS}\) paper-image matrix separated from the legacy
helper.  For the \(n=3,\kappa=7\) regression, source column \(8\) follows
Lemma~\texttt{Banded-sparse-access} and maps to row \(40\) in the active gate,
while the old helper still has its row-\(4\) entry.  This is only a
contract-drift guard; it does not prove injectivity, dagger cleanup, unitarity,
LCU closure, or block extraction.

The historical route guard
\texttt{oneTermRobinBlockEncodingProofRoute\_projectionSourceFreeze} packages
the signal-index-zero projection target with the same source blocker.  It
records the normalizer \(N_DN_f\kappa\), keeps the block-projection and
block-correctness obligations false, keeps the theorem-level circuit-unitary
and block-extraction obligations false, and keeps LCU correctness false.  The
active route no longer uses this source blocker as a proof-search switch.

The source-dependency classification for the active route has changed.  The
old row-dependent unused-branch route remains a historical
\texttt{contract-drift} audit.  The current missing ingredient is a
\texttt{classical-lean-lemma}: finite bit-slice arithmetic and finite-map
injectivity for \texttt{bandedSparseAccessPaperImage} on
\texttt{bandedSparseAccessPaperGlobalSlotSource}.

\subsection{Cycle 1 Wrapper-Slot Freeze (Historical)}

The historical route guard
\texttt{oneTermRobinBlockEncodingProofRoute\_sourceDecisionWrapperSlotsBlocked\_n3}
recorded that the concrete boundary column \(48\) still had no unused-branch
image rule in the row-dependent route.  This is now rejected-model memory, not
the active \(O_D^{BS}\) blocker.  The direct
\texttt{bandedSparseAccessUnusedBranchImageRuleContract} slot and the full
clean-domain wrapper slot both keep \texttt{proposedImageIndex = none}; their
image-specified obligations remain false.  The same guard keeps block
correctness and LCU correctness false.

This is a wrapper-level audit of the existing source-contract gap.  It does
not alter either active \(O_D^{BS}\) matrix and does not prove injectivity,
dagger cleanup, unitarity, LCU closure, or block extraction.

\subsection{Cycle 1 Prior-PDE Source-Anchor Guard}

The route guard
\texttt{oneTermRobinBlockEncodingProofRoute\_priorPDESourceAnchorsTranscript}
keeps the source import from the prior PDE sparse-access construction
synchronized with the one-term Robin route.  It records the public
arXiv:2405.12855v3 anchor, Definition~6, Lemma~1, the appendix source anchor,
and the decomposition text
\[
  O_A^{BS} = U^{SUM}(U_A^{(l)} \otimes I^n).
\]
The same guard keeps the prior resource claim unproved, keeps the
Robin-specific unused-branch image rule absent, and keeps lower proof search
disabled.

This source-anchor guard is not a reversible-extension theorem.  It preserves
the cited transcript for the padded sparse-access primitive, but it does not
specify an injective image for clean unused zero-amplitude sparse branches and
does not promote any \(O_D^{BS}\), LCU, or block-extraction flag.

\subsection{Cycle 1 Seven-Gate Flag Freeze}

The route guard
\texttt{oneTermRobinBlockEncodingProofRoute\_gateUnitaryFlags} keeps the active
seven-gate matrix product synchronized with Fig.~\texttt{fig:1 term ROBIN}.
The frozen flag vector is
\[
  [\texttt{true},\texttt{false},\texttt{false},\texttt{false},
   \texttt{false},\texttt{true},\texttt{false}].
\]
The true entries are the local permutation certificates for
\(U_{\mathrm{indic}}\) and SWAP.  The false entries are
\(O_{D^T}^{S}\), \(R_y^{\mathrm{boundary}}\), \(O_D^{BS}\), \(O_f\), and
\((O_D^{BS})^\dagger\).  The same guard keeps the active global-source cleanup
scope from promoting any \(O_D^{BS}\) semantic theorem and keeps the
theorem-level circuit-unitary and block-correctness flags false.

This is a proof-state guard.  It does not add an unused-branch image rule,
does not alter either active \(O_D^{BS}\) matrix, and does not prove
injectivity, dagger cleanup, unitarity, LCU closure, or block extraction.

\subsection{Cycle 1 Global-Source Cleanup Wrapper}

The active \(O_D^{BS}\) source predicate is now
\texttt{bandedSparseAccessPaperGlobalSlotSource}: the padded input to the
\(O_D^{BS}\) register is clean, and the encoded sparse slot satisfies
\(s<\kappa\).  This is the Lean-facing form of Lemma~\texttt{Banded-sparse-access},
the zero-inclusion paragraph before Theorem~\texttt{1 term robin}, and
Remark~\texttt{sparsity maximum}.

The wrapper
\texttt{bandedSparseAccessPostSwapCleanup\_of\_globalSlotSourceCandidate\_noRange}
feeds this active global-source predicate into the existing conditional
post-SWAP cleanup candidate.  The wrapper first extracts the clean padded-input
fact and then reuses
\texttt{bandedSparseAccessPostSwapCleanup\_of\_cleanSourceCandidate\_noRange}.
It is a conditional cleanup witness only: it does not prove preimage
uniqueness, the semantic \texttt{daggerCleanup} obligation, either
\(O_D^{BS}\) unitarity flag, LCU correctness, or block extraction.

The focused regression uses the \(n=3,\kappa=7\) source column \(48\).  That
column is accepted by the active global-source predicate and rejected by the
old row-dependent valid-source predicate, so the test confirms that the
cleanup route no longer depends on deleting zero-amplitude boundary slots from
the sparse register.

\subsection{Cycle 1 Global-Source Inverse Interface}

Lean now records the active inverse route as the obligation interface
\texttt{BandedSparseAccessGlobalSlotInverseOnRangeContract}, with default
record \texttt{bandedSparseAccessGlobalSlotInverseOnRangeContract}.  The
theorem
\texttt{bandedSparseAccessGlobalSlotInverseOnRangeContract\_of\_globalSlotSource}
proves that a global-source column satisfies the executable candidate checks,
while keeping \texttt{inverseOnRange.proved},
\texttt{uniquePreimage.proved},
\texttt{imageInjectiveOnGlobalSource.proved},
\texttt{daggerCleanup.proved}, and \texttt{unitaryExtension.proved} false.

The next lower packet should prove one fixed injectivity or unique-preimage
lemma over \texttt{bandedSparseAccessPaperGlobalSlotSource}.  It must not
promote \(O_D^{BS}\) cleanup, either \(O_D^{BS}\) unitarity flag, LCU
correctness, or final block extraction.

\subsection{Cycle 1 Post-SWAP Unique-Preimage Block}

Lean now proves the finite active-source uniqueness step for the corrected
\(O_D^{BS}\) route.  The declaration
\texttt{bandedSparseAccessPaperPostSwapPreimageCandidate\_sparseIndex\_eq}
states that the candidate preimage writes the reverse global sparse slot into
the clean padded \(O_D^{BS}\) register.  The declaration
\texttt{bandedSparseAccessPaperPostSwapPreimageCandidate\_globalSlotSource\_of\_globalSlotSource}
then proves that the candidate is itself in
\texttt{bandedSparseAccessPaperGlobalSlotSource}.

The theorem
\texttt{bandedSparseAccessPaperPostSwapPreimageCandidate\_unique\_on\_globalSlotSource}
uses the active image-injectivity block to show that any active global-source
preimage of
\[
  \texttt{swapOracleImage}\ p\,
  (\texttt{bandedSparseAccessPaperImage}\ p\ \texttt{source.val})
\]
is equal to the named candidate.  This is the Lean finite-register proof
needed by the cleanup route after Lemma~\texttt{Banded-sparse-access} and
Fig.~\texttt{fig:1 term ROBIN}.

No semantic obligation flag changes in this step.  The fields
\texttt{inverseOnRange.proved}, \texttt{uniquePreimage.proved},
\texttt{imageInjectiveOnGlobalSource.proved}, \texttt{daggerCleanup.proved},
and \texttt{unitaryExtension.proved} in
\texttt{BandedSparseAccessGlobalSlotInverseOnRangeContract} remain false.
The next proof-map step is a reviewed bridge from this compiled finite theorem
to the cleanup contract; it must not prove \(O_D^{BS}\) unitarity, LCU
correctness, or block extraction.

\subsection{Cycle 1 Cleanup-Contract Map}

Middle re-audited the GHL2025 source anchors before assigning the next
cleanup-contract packet: Lemma~\texttt{Diagonal sparsity},
Lemma~\texttt{Banded-sparse-access-oracle},
Remark~\texttt{sparsity maximum}, the zero-inclusion paragraph before
Theorem~\texttt{1 term robin}, Eq.~\texttt{eq: ROBIN clarified}, and
Fig.~\texttt{fig:1 term ROBIN}.  These anchors give the global sparse-slot
map
\[
  \ket{0}^{n-l}\ket{s}^l\ket{i}^n
  \mapsto
  \ket{r_{si}}^n\ket{i}^n,
  \qquad
  r_{si}=r_{s0}+i \bmod 2^n,
\]
and the circuit-level use of \((O_D^{BS})^\dagger\) after SWAP.  They do not
provide a separate QBE theorem that promotes the semantic cleanup flag.

The lower target for this packet is the non-promoting Lean theorem
\texttt{bandedSparseAccessGlobalSlotInverseOnRangeContract\_cleanupContractMap}.
For an active global-source column, it should combine
\texttt{bandedSparseAccessGlobalSlotInverseOnRangeContract\_daggerCleanupBridge}
with
\texttt{bandedSparseAccessGlobalSlotInverseOnRangeContract\_uniquePreimageBridge}.
The theorem should expose the contract post-SWAP target, the candidate
preimage, the \texttt{BandedSparseAccessPostSwapCleanup} witness, uniqueness
among active global-source preimages, and the transpose-style dagger matrix
entry.  It must keep \texttt{daggerCleanup.proved},
\texttt{unitaryExtension.proved}, both \(O_D^{BS}\) unitarity flags, LCU
correctness, circuit unitarity, block projection, and block correctness false.

\textbf{Lower cleanup-contract wrapper.}
Lean now defines
\texttt{bandedSparseAccessGlobalSlotInverseOnRangeContract\_cleanupContractMap}.
For an active global-source column, the theorem returns the contract
post-SWAP target, the candidate preimage, a
\texttt{BandedSparseAccessPostSwapCleanup} witness, active-source uniqueness
for any preimage of the same post-SWAP target, and the transpose-style dagger
matrix entry.  It is only a wrapper over the existing dagger-cleanup bridge
and unique-preimage bridge; the \texttt{inverseOnRange},
\texttt{uniquePreimage}, \texttt{imageInjectiveOnGlobalSource},
\texttt{daggerCleanup}, and \texttt{unitaryExtension} proof flags remain
false.

\textbf{Default paper-contract cleanup bridge.}
Lean now defines
\texttt{defaultBandedSparseAccessPaperContract\_cleanupRouteBridge}.  The
theorem consumes
\texttt{bandedSparseAccessGlobalSlotInverseOnRangeContract\_cleanupContractMap}
and re-exposes the same post-SWAP cleanup route through
\texttt{defaultBandedSparseAccessPaperContract}.  For an active global-source
column it returns the post-SWAP target, the candidate preimage, a
\texttt{BandedSparseAccessPostSwapCleanup} witness, active-source uniqueness
for the candidate, and the transpose-style dagger matrix entry.

This is still a proof-state bridge, not a semantic cleanup proof.  The default
paper-contract fields \texttt{daggerCleanup.proved} and
\texttt{unitaryExtension.proved} remain false, and both active \(O_D^{BS}\)
gate-unitarity flags remain false.  The next step must decide the exact scope
of any semantic cleanup theorem before attempting to promote a proof flag.
The regression
\texttt{defaultBandedSparseAccessPaperContract\_cleanupRouteBridge\_boundaryColumn\_n3}
specializes the bridge to the historical boundary column \(48\): that column
is active for \texttt{bandedSparseAccessPaperGlobalSlotSource}, rejected by the
old row-dependent source classifier, and still leaves the cleanup and unitarity
flags false.

\textbf{Off-candidate dagger-zero bridge.}
Lean now defines
\texttt{bandedSparseAccessGlobalSlotInverseOnRangeContract\_daggerOffCandidate\_zero}.
For the contract post-SWAP target, the theorem proves that an active
global-source row different from the named candidate has entry \(0\) in the
transpose-style \((O_D^{BS})^\dagger\) matrix.  It uses the compiled
active-source unique-preimage bridge and the definition of
\texttt{bandedSparseAccessPaperDaggerMatrix}.

This is still a restricted active-domain matrix-entry fact.  It is not a
full-space inverse theorem and it does not promote
\texttt{daggerCleanup.proved}, either \(O_D^{BS}\) unitarity flag, LCU
correctness, circuit unitarity, block projection, or block correctness.

\textbf{Restricted dagger-column cleanup.}
Lean now defines
\texttt{bandedSparseAccessGlobalSlotInverseOnRangeContract\_restrictedDaggerColumnCleanup}.
For the contract post-SWAP target, the theorem packages the named candidate
preimage, the \texttt{BandedSparseAccessPostSwapCleanup} witness, the
candidate dagger entry \(1\), and zero dagger entries for every other row in
\texttt{bandedSparseAccessPaperGlobalSlotSource}.  The scope is exactly this
active global-source column.  It does not cover non-clean rows, encoded sparse
values outside \(s<\kappa\), a full clean-domain reversible extension, or a
full-space inverse theorem.

Lean now defines
\texttt{bandedSparseAccessGlobalSlotInverseOnRangeContract\_restrictedDaggerColumnIndicator}
and the theorem-route wrapper
\texttt{oneTermRobinBlockEncodingProofRoute\_odbsRestrictedDaggerColumnIndicator}.
For each active global-source row \texttt{other}, the entry is
\[
  \begin{cases}
    1, & \text{if }\texttt{other.val = pre.val},\\
    0, & \text{otherwise}.
  \end{cases}
\]
The proof reuses the restricted cleanup theorem and keeps
\texttt{daggerCleanup.proved}, \texttt{unitaryExtension.proved}, both
\(O_D^{BS}\) unitarity flags, LCU correctness, circuit unitarity, block
projection, and block correctness false.

\textbf{Cleanup-scope decision.}
Lean now defines
\texttt{BandedSparseAccessCleanupScopeDecision} and
\texttt{bandedSparseAccessCleanupScopeDecision}.  The selected cleanup theorem
scope is the active global-source predicate
\texttt{bandedSparseAccessPaperGlobalSlotSource}, with current evidence
\texttt{bandedSparseAccessGlobalSlotInverseOnRangeContract\_restrictedDaggerColumnIndicator}.
The route guard
\texttt{oneTermRobinBlockEncodingProofRoute\_odbsCleanupScopeDecision}
synchronizes this decision with the one-term Robin proof route.  It explicitly
does not select a full clean-domain theorem or a full-space theorem, and it
keeps semantic cleanup promotion, \(O_D^{BS}\) unitary extension, LCU
correctness, and block correctness false.

\subsection{Cycle 11: Coeff.rat Arithmetic and U\_indic Bijection}

\textbf{Coeff.rat arithmetic lemmas.}
The \texttt{Coeff} type wraps rational arithmetic inside a symbolic expression
tree.  For permutation matrices (entries in $\{0, 1\}$), the matrix product
reduces to counting matches, but Lean must still normalize the \texttt{Coeff}
expression tree to verify the count.  The new lemmas in \texttt{Core.lean}:
\[
  \texttt{evalWith\_rat\_add}: \;
  \texttt{evalWith}\;\texttt{env}\;(\texttt{Coeff.add}\;(\texttt{rat}\;a)\;(\texttt{rat}\;b)) = a + b,
\]
and analogous lemmas for \texttt{rat\_mul}, \texttt{rat\_neg}, \texttt{rat\_zero},
\texttt{rat\_one}, close the normalization gap for concrete \texttt{Coeff.rat}
arithmetic.

\textbf{U\_indic bijection stepping stone.}
The indicator oracle $U_{\mathrm{indic}}$ is a permutation matrix: each basis
state $|j\rangle$ maps to $|\texttt{image}(j)\rangle$, where
\[
  \texttt{image}(j) = j \oplus (\mathrm{isBulk}(j) \ll \texttt{indPos}).
\]
The mapping is self-inverse: $\texttt{image}(\texttt{image}(j)) = j$, because
XOR with a fixed bit is an involution.  In Lean this is verified by
\texttt{native\_decide} for concrete $n = 1$ (space $2^7 = 128$) and
$n = 3$ (space $2^{13} = 8192$).

A self-inverse function on a finite set is a bijection, and a 0/1 matrix
induced by a bijection is a permutation matrix, which is unitary.  The
remaining proof obligation is to formalize this chain in Lean for general~$n$.

\subsection{2026-05-25 Active Route Retirement Sync}

The theorem-route record no longer carries the historical
\texttt{unusedZeroBranchDecision}, \texttt{priorSparseAccessSource}, or
\texttt{robinZeroInclusionSource} fields.  Those source contracts remain
separate cited-results and transcript memory, but the active \(O_D^{BS}\)
route is now the global sparse-slot source selected by
\texttt{bandedSparseAccessCleanupScopeDecision}.

The current active route declarations are
\texttt{oneTermRobinBlockEncodingProofRoute\_odbsActiveGlobalSlotBlockers},
\texttt{oneTermRobinBlockEncodingProofRoute\_odbsActiveGlobalSlotGateFreeze},
\texttt{oneTermRobinBlockEncodingProofRoute\_theoremTranscriptDependencies},
and
\texttt{oneTermRobinBlockEncodingProofRoute\_theoremTranscriptActiveCleanupMap}.
For the \(n=3,\kappa=7\) regression,
\texttt{oneTermRobinBlockEncodingProofRoute\_encodedOutOfRangeSparseSlot\_n3}
records that a clean encoded sparse value \(s=7\) is outside the active source
range, while
\texttt{oneTermRobinBlockEncodingProofRoute\_rejectedRowDependentCollisionRegression\_n3}
keeps the old row-dependent collision as rejected-model memory only.

No semantic flag is promoted by this retirement.  The active route still keeps
\texttt{daggerCleanup.proved}, \texttt{unitaryExtension.proved}, \(O_f\)
amplitude correctness, LCU correctness, circuit unitarity, block projection,
block correctness, resource bounds, ancilla cleanup, and final block
extraction false.

\subsection{2026-05-25 Derivative and Boundary Route Bridge}

The theorem-route guard
\texttt{oneTermRobinBlockEncodingProofRoute\_derivativeBoundaryContractMap}
connects the existing \(O_{D^T}^{S}\) and boundary \(R_y\) source contracts to
\texttt{oneTermRobinBlockEncodingProofRoute}.  It consumes
\texttt{derivativeNormalizerNDSharedRoute\_sourceBoundAndFlags} and records
that the two routes use the same coefficient source \(D_j^{(s)}\), the same
normalizer symbol \(N_D\), and the same normalizer-bound obligation.

The guard also pins the Fig.~1-term Robin gate positions: gate slot 1 is
\(O_{D^T}^{S}\), and gate slot 2 is the boundary \(R_y\) rotation.  It does
not prove the analytic identities behind Eq.~(20) or Eq.~\texttt{angles for
Ry}; nonzero division by \(N_D\), coefficient bounds, absolute-square and
square-root semantics, arccos semantics, half-angle identities, and
two-by-two unitarity remain proof obligations.  The theorem-level circuit
unitarity, \(O_f\) amplitude correctness, LCU correctness, projection,
block-correctness, and final extraction flags remain false.

\subsection{2026-05-25 Theorem Transcript Closure Packet}

The guard
\texttt{oneTermRobinBlockEncodingProofRoute\_theoremTranscriptClosurePacket}
is the current Phase~1 closure checkpoint for Theorem~\texttt{1 term robin}.
It consumes the theorem source transcript, the layout/projection audit, the
block-projection dependency map, and the full Fig.~\texttt{1 term ROBIN}
gate-contract ledger.

The packet records the paper source anchor, the Lemma~\texttt{Banded-sparse-access}
image formula \(r_{si}=r_{s0}+i \bmod 2^n\), the GL2024 cited source anchor
for \(O_f\), the theorem normalizer \(N_DN_f\kappa\), the theorem signal-count
and \(2n\) pure-ancilla ledger, the signal-index-zero block convention, and
the seven-gate proof-state vector
\[
  [\texttt{true},\texttt{false},\texttt{false},
    \texttt{false},\texttt{false},\texttt{true},\texttt{false}].
\]
It also records that the \(O_{D^T}^{S}\) and boundary \(R_y\) routes share the
same coefficient source \(D_j^{(s)}\).

This is not a semantic block-encoding proof.  The closure packet keeps the
resource-bound proof, ancilla cleanup, \(O_D^{BS}\) cleanup and unitary
extension, \(O_f\) amplitude correctness, LCU correctness, circuit unitarity,
block projection, block correctness, and final block extraction false.  Future
lower packets should select one of these false dependencies explicitly rather
than repeating the theorem-transcript wiring work.

\subsection{2026-05-25 Finite Composition Contract}

Lean now defines
\texttt{oneTermRobinFiniteBlockCompositionContract} and the route guard
\texttt{oneTermRobinBlockEncodingProofRoute\_finiteBlockCompositionContractMap}.
The contract is the current Lean-facing form of the finite matrix step behind
Theorem~\texttt{1 term robin}.  It stores the existing circuit claim, the
target \texttt{oneTermRobinBlockExtractionTarget}, the matrix
\texttt{oneTermRobinAkMatrix n}, and the normalizer \(N_DN_f\kappa\).

The source audit for this contract uses Theorem~\texttt{1 term robin},
Eq.~\texttt{ROBIN clarified}, Fig.~\texttt{1 term ROBIN}, and the paper's
block-encoding definition.  These anchors provide the theorem statement, the
\(\gamma_3\) clean coefficient, the seven-gate circuit order, and the
zero-signal projection convention.  The source does not give a separate finite
matrix proof that the product of the seven gate matrices has the required
signal-zero block.  QBE therefore records this step as a contract-only
obligation, with the cited-results row \texttt{LCU.StandardBlockEncoding}
serving only as background until an exact finite theorem is accepted.

The contract fields
\texttt{circuitUnitary}, \texttt{lcuComposition}, \texttt{blockProjection},
\texttt{normalizedBlockEquality}, and \texttt{finalExtraction} all remain
false.  The route-level LCU correctness, circuit unitarity, block projection,
block correctness, resource-bound proof, ancilla cleanup, and final extraction
also remain false.

\subsection{2026-05-25 Derivative-Amplitude Entry Packet}

The next fixed entry bridge concerns the derivative factor in the
\(\gamma_3\) route.  Lemma~\texttt{Sparse-amplitude-oracle for a
banded-sparse matrix} and Eq.~\texttt{amplitude\_oracle\_D} supply the
ket-zero amplitude \(D^{(s)}/N_D\) for \(O_{D^T}^{S}\).  Eq.~\texttt{angles
for Ry} supplies the boundary rotation angle
\(\theta_j^s = \arccos(D_j^{(s)}/N_D)\).  Fig.~\texttt{1 term ROBIN} fixes
these as gate slots 1 and 2 in the Lean route.

The existing guard
\texttt{oneTermRobinBlockEncodingProofRoute\_derivativeBoundaryContractMap}
already pins the two gate slots and shares the same \(D_j^{(s)}\) source and
\(N_D\) normalizer record.  The next lower target is a narrower entry bridge:
\texttt{oneTermRobinBlockEncodingProofRoute\_odtsKetZeroEntry}.  Under the
indicator-one and ancilla-zero hypotheses for a full-basis column, it should
show that route gate slot 1 has entry
\[
  \texttt{sparseAmplitudeOracleDTCoefficientNormalizerProofRoute.ketZeroEntry}.
\]
If this theorem is stable, the same patch may add
\texttt{oneTermRobinBlockEncodingProofRoute\_boundaryKetZeroEntry}, exposing
the gate-slot 2 boundary entry
\[
  \texttt{BoundaryRotationAngleNormalizerProofRoute.cosHalfEntry}.
\]

These entry bridges are transcript links only.  They do not prove division by
\(N_D\), the \(N_D\) bound, square-root semantics, arccos semantics,
half-angle identities, two-by-two unitarity, the full seven-gate product, LCU
composition, block projection, or final extraction.

\subsection{2026-05-25 Boundary-Rotation Entry Packet}

The declaration
\texttt{oneTermRobinBlockEncodingProofRoute\_odtsKetZeroEntry} now covers the
gate-slot 1 derivative-amplitude entry.  The next fixed theorem-facing packet
is the analogous gate-slot 2 bridge for the boundary rotation.  Eq.~\texttt{angles for Ry} defines
\[
  \theta_j^s = \arccos(D_j^{(s)}/N_D)
\]
for boundary rows and sparse slots, while Fig.~\texttt{1 term ROBIN} places
\(\mathrm{Ry}_{\mathrm{boundary}}\) immediately after \(O_{D^T}^{S}\) in the
seven-gate route.

The lower bridge
\texttt{oneTermRobinBlockEncodingProofRoute\_boundaryKetZeroEntry} now compiles.
For full basis indices \(i,j\), let
\[
  \texttt{regs}
  =
  \texttt{GHL2025.boundaryRotationPaperRegisters p j.val}.
\]
Under the hypotheses
\(\texttt{regs.indicatorBit = 0}\),
\(\texttt{regs.ancillaBit = 0}\),
\(\texttt{i.val >>> 1 = regs.nonAncillaValue}\), and
\(\texttt{i.val \&\&\& 1 = 0}\), the bridge shows that route gate slot
2 has entry
\[
  \texttt{BoundaryRotationAngleNormalizerProofRoute.cosHalfEntry}.
\]
This statement is a transcript bridge only.  It leaves the arccos semantics,
half-angle formulas, \(N_D\) bound, two-by-two unitarity, LCU composition,
block projection, normalized block equality, and final extraction false.

The focused Lean regression specializes to \(n=3\), boundary column \(0\), and
boundary row \(0\).  It checks that route gate slot 2 equals
\[
  \texttt{Coeff.symbol "boundary\_cos\_half\_0\_0"}.
\]
The same test records that \(O_{D^T}^{S}\) unitarity, boundary-rotation
unitarity, \(O_D^{BS}\) cleanup and unitarity, \(O_f\) amplitude correctness,
LCU correctness, projection, block correctness, and theorem extraction all
remain unproved obligations.

\subsection{2026-05-25 Gamma3 Factor-Entry Ledger}

The declaration
\texttt{oneTermRobinBlockEncodingProofRoute\_gamma3FactorEntryLedger}
packages the currently compiled entry bridges for the future
\(\gamma_3\) signal-block theorem.  It consumes
\texttt{oneTermRobinBlockEncodingProofRoute\_gamma3TargetEntryData},
\texttt{oneTermRobinBlockEncodingProofRoute\_ofCleanFunctionOracleEntry},
\texttt{oneTermRobinBlockEncodingProofRoute\_odtsKetZeroEntry},
\texttt{oneTermRobinBlockEncodingProofRoute\_boundaryKetZeroEntry}, and
\texttt{oneTermRobinBlockEncodingProofRoute\_odbsActiveGlobalSourceCleanupContractMap}.

The source anchors are Theorem~\texttt{1 term robin},
Eq.~\texttt{ROBIN clarified}, Fig.~\texttt{1 term ROBIN},
Lemma~\texttt{Banded-sparse-access-oracle},
Lemma~\texttt{Sparse-amplitude-oracle for a banded-sparse matrix}, and
Theorem~\texttt{Amplitude-oracle for piece-wise polynomial function}.  The
ledger exposes the target entry and normalizer \(N_DN_f\kappa\), the clean
\(O_f\) entry \(f(x_i)/N_f\), the \(O_{D^T}^{S}\) ket-zero symbolic entry,
the boundary \(R_y\) ket-zero symbolic entry, and the active global-source
\((O_D^{BS})^\dagger\) cleanup entry.

This is still a transcript ledger.  It does not prove the finite product of
the seven gate matrices, the normalized block equality, the LCU composition
step, full \(O_D^{BS}\) cleanup, \(O_f\) amplitude correctness, circuit
unitarity, block projection, block correctness, or final extraction.  The
focused Lean regression at \(n=3\) checks the entries
\texttt{f\_3\_2 * N\_f\_inv},
\texttt{odts\_cos\_half\_2\_0},
\texttt{boundary\_cos\_half\_0\_0}, and the active dagger entry
\texttt{Coeff.rat 1}, while keeping all semantic flags false.

\subsection{2026-05-25 Ak Target-Matrix Rewire}

The theorem target is the one-term matrix \(A_k\), not the derivative matrix
\(D\) alone.  The paper states Theorem~\texttt{1 term robin} for the
discretized operator \(A_k \sim f(x)\partial^m/\partial x^m\), and the
\(\gamma_3\) line of Eq.~\texttt{ROBIN clarified} has coefficient
\[
  \frac{f(x_i)(D)_i^{(s)}}{N_DN_f\kappa}.
\]
Thus the Lean target must expose entries
\[
  (A_k)_{ij}=f(x_i)D_{ij}.
\]

The drift is now corrected in the Lean transcript.  The declaration
\texttt{oneTermRobinAkMatrix} uses the existing
\texttt{GHL2025.robinFunctionValue} and \texttt{robinDerivativeMatrix}, and
\texttt{oneTermRobinAkMatrix\_apply} exposes the entry formula above.  The
block-extraction target, finite-composition contract, theorem route, and
\texttt{oneTermRobinBlockEncodingProofRoute\_gamma3TargetEntryData} now point
to this row-scaled matrix.

This correction is still Phase~1 transcript work.  It does not promote
\(\Uf\) amplitude correctness, LCU composition, normalized block equality,
block projection, block correctness, cleanup, unitarity, or final extraction.
The follow-on contract
\texttt{oneTermRobinBlockEncodingProofRoute\_gamma3AkCoefficientEntryContract}
now compiles and synchronizes this Ak entry with the product-entry and
factor-entry ledgers while leaving all semantic flags false.

\subsection{2026-05-25 Gamma3 Product-To-Coefficient Interface}

The declaration
\texttt{oneTermRobinBlockEncodingProofRoute\_gamma3ProductToCoefficientInterface}
now names the remaining finite entry theorem for the \(\gamma_3\) branch.
It consumes
\texttt{oneTermRobinBlockEncodingProofRoute\_gamma3AkCoefficientEntryContract}
and introduces
\texttt{oneTermRobinGamma3ProductToCoefficientObligation}.  The source target
is the entry equality
\[
  \bigl(\langle 0| \otimes I\bigr)U_{A_k}^{(1)}
  \bigl(|0\rangle \otimes I\bigr)_{ij}
  =
  \frac{(A_k)_{ij}}{N_DN_f\kappa}
  =
  \frac{f(x_i)D_{ij}}{N_DN_f\kappa}.
\]

This interface is not a proof of the equality.  It exposes the signal-zero
product entry, the Ak expansion \(A_{k,ij}=f(x_i)D_{ij}\), the inherited
factor entries for \(O_f\), \(O_{D^T}^{S}\), \(R_y^{\mathrm{boundary}}\), and
the active global-source \(O_D^{BS}\) cleanup entry.  It also records that
LCU composition, normalized block equality, cleanup, unitarity, projection,
block correctness, circuit unitarity, and final extraction remain false.

The next lower packet is a focused proof attempt, not another guard.  It
should use the \(n=3\) specialization with source column \(48\), clean
\(O_f\) column \(36\), \(O_{D^T}^{S}\) column \(132\), boundary column \(0\),
and target entry \((2,5)\).  An accepted result is either a compiled equality
or path-isolation lemma that materially advances the finite product proof.  If
the equality cannot yet be stated or proved, lower must update
\texttt{proof-attempts/QBE-AUTO-002/gamma3-product-to-coefficient.md} with the
exact Lean goal and missing interface.  No semantic flag may be promoted in
that attempt.

\subsection{2026-05-25 Gamma3 Projection Layout Audit}

The focused product attempt exposed a projection-layout drift before the
unique-path lemma could be applied.  The compiled declaration
\texttt{oneTermRobinBlockEncodingProofRoute\_gamma3ProjectionPathAudit\_n3}
shows that the current generic block helper sends the system entry
\((2,5)\) to the full matrix entry \((2,5)\).  Along the executable gate
images, the forward branch from full column \(5\) reaches the final dagger
column \(160\), where row \(2\) has entry \(0\).  The previously ledgered
columns \(48\), \(36\), \(132\), and \(0\) are therefore separate source and
factor witnesses, not one seven-gate path for the projected entry.

The paper ket order in Eq.~\texttt{ROBIN clarified} explains the mismatch.
The clean \(\gamma_3\) ket has the form
\[
  |0\rangle^{m_f+1}|s\rangle^{\lceil\log_2\kappa\rceil}
  |0\rangle^{n-\lceil\log_2\kappa\rceil}|j\rangle^n|0\rangle .
\]
In the Lean register layout, the rightmost paper ancilla is bit \(0\), and
the system register occupies bits \([1,1+n)\).  Thus a clean paper-basis index
has the form
\[
  (s \ll (1+n+odPure)) + (j \ll 1),
  \qquad odPure=n-\lceil\log_2\kappa\rceil .
\]
For \(n=3\) and \(\kappa=7\), this becomes
\[
  (s \ll 4) + (j \ll 1).
\]
The zero-sparse system values \(2\) and \(5\) therefore correspond to full
indices \(4\) and \(10\), not to \(2\) and \(5\).

This is classified as contract drift in the Lean projection convention.  The
Robin-specific layout bridge has now compiled as
\texttt{oneTermRobinBlockEncodingProofRoute\_gamma3PaperBasisLayout\_n3}.  It
proves the paper clean-basis index formula against the existing register
extractors for the focused \(n=3\), \(\kappa=7\), sparse-slot-zero branch, and
it records the full endpoints \((4,10)\).

The next lower packet should audit the Fig.~\texttt{1 term ROBIN} path from
column \(10\) to row \(4\), suggested as
\texttt{oneTermRobinBlockEncodingProofRoute\_gamma3PaperBasisPathAudit\_n3}.
It should list the adjacent states through
\(U_{\mathrm{indic}}\), \(O_{D^T}^{S}\), \(R_y^{\mathrm{boundary}}\),
\(O_D^{BS}\), \(O_f\), SWAP, and \((O_D^{BS})^\dagger\) before applying
\texttt{Matrix.evalWith\_mul\_unique\_path}.  Product-to-coefficient,
signal-block, LCU, cleanup, unitarity, projection, block correctness, circuit
unitarity, normalized equality, and final extraction remain false.

\subsection{2026-05-25 Gamma3 Paper-Basis Path Audit}

The compiled declaration
\texttt{oneTermRobinBlockEncodingProofRoute\_gamma3PaperBasisPathAudit\_n3}
audits the Fig.~\texttt{1 term ROBIN} gate order at the paper-basis endpoint
\((4,10)\) for the focused \(n=3\), \(\kappa=7\), sparse-slot-zero branch.
The ket-zero branch from column \(10\) is
\[
10 \xrightarrow{U_{\mathrm{indic}}} 138
   \xrightarrow{O_{D^T}^{S}} 138
   \xrightarrow{R_y^{\mathrm{boundary}}} 138
   \xrightarrow{O_D^{BS}} 186
   \xrightarrow{O_f} 186
   \xrightarrow{\mathrm{SWAP}} 214
   \xrightarrow{(O_D^{BS})^\dagger} 198 .
\]
The final dagger column is \(214\).  Row \(198\) has dagger entry \(1\), while
the paper clean row \(4\) has dagger entry \(0\).

The first decisive mismatch is the active \(O_D^{BS}\) address.  At column
\(138\), the extracted row is \(5\) and the sparse slot is \(0\), so the
global sparse-slot address writes \(3\).  After SWAP, the visible system field
is \(3\), not the target paper row \(2\).  The next proof obligation is
therefore to map the paper coefficient \(D_{ij}\) to the correct sparse slot
or projection convention before applying
\texttt{Matrix.evalWith\_mul\_unique\_path}.  Product-to-coefficient,
signal-block, LCU, cleanup, unitarity, projection, block correctness, circuit
unitarity, normalized equality, and final extraction remain false.

For the focused entry \(D_{2,5}\), the next lower packet should first audit the
global sparse slot satisfying \(r_{s,5}=2\).  The slot-zero audit is retained as
negative evidence: it proves \(r_{0,5}=3\) in the active global-slot table and
therefore cannot be multiplied into the desired coefficient path.  This is a
finite register-index obligation and a projection-slot contract question; it
is not a new external theorem.

\subsection{2026-05-25 Gamma3 Projection-Slot Convention}

The focused coefficient \(D_{2,5}\) uses sparse slot \(5\), since the active
global sparse-address table gives \(r_{5,5}=2\).  Lean records this
projection-slot handoff in
\texttt{oneTermRobinBlockEncodingProofRoute\_gamma3ProjectionSlotConventionMap\_n3}.
The compiled endpoint facts are
\[
  \text{generic projection endpoints }(2,5),
  \qquad
  \text{slot-}5\text{ clean endpoints }(84,90),
\]
with
\[
  O_D^{BS}(90)=42,
  \qquad
  \mathrm{SWAP}(42)=84,
  \qquad
  U_{\mathrm{indic}}(90)=218 .
\]
The corresponding source-contract obligation is
\texttt{oneTermRobinGamma3ProjectionSlotConventionObligation}.  It asks for
the theorem-level convention that relates a slot-specific clean branch in
Eq.~\texttt{ROBIN clarified} to the block projection or sparse-register
summation in Definition~\texttt{def:block-encoding}.  This obligation is
compiled with \texttt{proved = false}.

The next lower packet should audit the Fig.~\texttt{1 term ROBIN} gate path
for source column \(90\) and candidate row \(84\).  It should not promote the
projection-slot convention, product-to-coefficient equality, LCU composition,
cleanup, unitarity, block projection, block correctness, circuit unitarity,
normalized equality, or final extraction.

\subsection{2026-05-25 Gamma3 Slot-5 Path Audit}

The compiled declaration
\texttt{oneTermRobinBlockEncodingProofRoute\_gamma3Slot5PathAudit\_n3}
audits the Fig.~\texttt{1 term ROBIN} gate order for the slot-\(5\)
coefficient path of \(D_{2,5}\).  The clean sparse-access subchain is
\[
  90 \xrightarrow{O_D^{BS}} 42 \xrightarrow{\mathrm{SWAP}} 84,
\]
but this is not the full circuit path, because
\(U_{\mathrm{indic}}\) is the first gate in Fig.~\texttt{1 term ROBIN}.  The
ket-zero branch through the seven active gates is
\[
  90 \to 218 \to 218 \to 218 \to 170 \to 170 \to 212 \to 228.
\]
The final dagger entry at row \(228\), column \(212\) is \(1\), while the entry
at the clean slot-\(5\) row \(84\), column \(212\) is \(0\).  The final row
has sparse-index value \(6\), while the clean row has sparse-index value \(5\).

Middle re-read Theorem~\texttt{1 term robin}, Eq.~\texttt{ROBIN clarified},
Fig.~\texttt{1 term ROBIN}, Lemma~\texttt{Banded-sparse-access-oracle}, Lemma
\texttt{Sparse-amplitude-oracle for a banded-sparse matrix}, and
Definition~\texttt{def:block-encoding}.  The source gives the clean
\(\gamma_3\) ket and the circuit, but it does not include a finite
basis-index proof that identifies row \(228\) with row \(84\), nor does it
state a sparse-register summation or basis-permutation convention after the
full seven-gate product.  The remaining blocker is therefore a
projection/register source-contract gap plus local register arithmetic, not a
new external theorem and not another product-fold search.

The next lower packet should compile a field-level audit, suggested as
\texttt{oneTermRobinBlockEncodingProofRoute\_gamma3Slot5ProjectionRegisterAudit\_n3}.
It should compare the indicator bit, trailing ancilla bit, system value,
padded \(\ODBS\) zero field, sparse-index value, \(m_f\) clean workspace, and
active-source predicate for states \(90\), \(218\), \(170\), \(212\), \(228\),
and \(84\).  Product-to-coefficient, signal-block, LCU, cleanup, unitarity,
projection, block correctness, circuit unitarity, normalized equality, and
final extraction remain false.

\subsection{2026-05-25 Gamma3 Projection/Register Decision}

The requested field-level audit now compiles as
\texttt{oneTermRobinBlockEncodingProofRoute\_gamma3Slot5ProjectionRegisterAudit\_n3}.
It compares the full Fig.~\texttt{1 term ROBIN} endpoint \(228\) with the
clean Eq.~\texttt{ROBIN clarified} endpoint \(84\) for the focused
slot-\(5\) coefficient of \(D_{2,5}\).  The comparison records that the
system row, padded \(\ODBS\) zero field, trailing ancilla bit, \(m_f\)
clean-workspace predicate, and active-source predicate agree.  The first
mismatch in the requested comparison order is the indicator bit:
row \(228\) has indicator \(1\), while row \(84\) has indicator \(0\).
The sparse-index field also differs, with row \(228\) carrying sparse index
\(6\) and row \(84\) carrying sparse index \(5\).

Middle promoted this result only to a decision record,
\texttt{oneTermRobinGamma3ProjectionRegisterConventionDecision\_n3}, with
transcript theorem
\texttt{oneTermRobinGamma3ProjectionRegisterConventionDecision\_n3\_transcript}.
The record states that the required convention is still a source-contract gap
and that product search is blocked.  It does not prove the
\(\gamma_3\) product-to-coefficient equality, normalized block equality, LCU
composition, cleanup, unitarity, block projection, block correctness, circuit
unitarity, or final extraction.

The next action is not another lower-agent multiplication attempt.  A source
or human decision must first choose one exact Lean-facing convention relating
the full endpoint \(228\) to the clean endpoint \(84\): a sparse-register
summation rule, a basis-permutation bridge, or a revised block-projection
layout.  After that convention is fixed as a theorem statement, lower work may
resume on a single product-to-coefficient target.

\subsection{2026-05-25 Gamma3 Sparse-Register Summation Convention}

The chosen convention statement now compiles as
\texttt{oneTermRobinGamma3SparseRegisterSummationConvention\_n3}, with
transcript theorem
\texttt{oneTermRobinGamma3SparseRegisterSummationConvention\_n3\_transcript}.
It selects sparse-register summation for the focused \(n=3\), system-entry
\((2,5)\) gamma3 audit because Eq.~\texttt{ROBIN clarified} writes the
\(\gamma_3\) branch as a sum over \(s=0,\dots,\kappa-1\).

The compiled record keeps the endpoint facts explicit:
\[
\begin{array}{c|ccc}
\text{endpoint} & \text{system row} & \text{indicator} & \text{sparse index} \\
\hline
84 & 2 & 0 & 5 \\
228 & 2 & 1 & 6
\end{array}
\]
The convention is therefore only a contract map.  It does not prove that the
current block projection implements sparse-register summation, and it does not
absorb the indicator mismatch.  The field
\texttt{indicatorMismatchObligation.proved = false} remains the next explicit
projection/register obligation before any gamma3 product-to-coefficient proof
search resumes.

\subsection{2026-05-25 Gamma3 Indicator-Field Gap}

The lower indicator-field check now compiles as
\texttt{oneTermRobinGamma3SparseRegisterSummation\_indicatorGap\_n3}.  It
reuses the sparse-register summation convention and records that this
convention does not explain the indicator bit: the clean endpoint \(84\) has
indicator \(0\), while the full endpoint \(228\) has indicator \(1\).  The
system row still agrees, and the sparse-index mismatch remains secondary
evidence.

The next proof-map target is an exact indicator projection/register convention
for the same endpoint pair.  It must be tied to Eq.~\texttt{ROBIN clarified},
Fig.~\texttt{1 term ROBIN}, and Definition~\texttt{def:block-encoding}.  Until
that convention is stated, product-to-coefficient search, LCU correctness,
block projection, block correctness, normalized block equality, and final
extraction remain false.

\subsection{2026-05-25 Gamma3 Indicator Projection Convention}

The indicator projection convention now compiles as
\texttt{oneTermRobinGamma3IndicatorProjectionConvention\_n3}, with transcript
theorem
\texttt{oneTermRobinGamma3IndicatorProjectionConvention\_n3\_transcript}.
The source audit reread Eq.~\texttt{ROBIN clarified},
Fig.~\texttt{1 term ROBIN}, the paragraph defining
$U_{\mathrm{indic}}$, and Definition~\texttt{def:block-encoding}.  These
anchors specify the clean $\gamma_3$ branch and the fact that
$U_{\mathrm{indic}}$ sets the bulk indicator, but they do not state a reset,
summation, or basis-permutation rule for the indicator field after
sparse-register summation.

\[
\begin{array}{c|cc}
\text{field} & \text{clean endpoint }84 & \text{full endpoint }228 \\
\hline
\text{indicator} & 0 & 1 \\
\text{source relation specified} & \multicolumn{2}{c}{\text{false}} \\
\end{array}
\]

The compiled record therefore keeps
\texttt{conventionObligation.proved = false}.  This is classified as a
source-contract gap plus an internal paper-step interface, not as a finite
matrix multiplication problem.  The next lower packet must either replace this
false obligation with one source-backed projection/register rule, or keep the
gamma3 product-to-coefficient route blocked.  No LCU, projection,
block-correctness, normalized equality, cleanup, unitarity, or final-extraction
flag is promoted by this convention record.

\subsection{2026-05-25 Gamma3 Human-Input Freeze}

The middle source-dependency audit now treats the indicator-field relation as a
human or source-backed convention decision.  The Lean record
\texttt{oneTermRobinGamma3IndicatorProjectionConvention\_n3} has
\texttt{indicatorRelationSpecifiedBySource = false},
\texttt{humanInputRequired = true}, and
\texttt{conventionObligation.proved = false}.  This states that the current
GHL2025 transcript does not specify how the theorem-level projection should
relate the full endpoint \(228\), where the indicator is \(1\), to the clean
endpoint \(84\), where the indicator is \(0\).

The next accepted lower step must supply one source-backed projection/register
rule for this indicator field or leave the gamma3 product-to-coefficient proof
blocked.  Sparse-register summation alone is not a proof of the indicator
relation.

\subsection{2026-05-25 Gamma3 Bulk-Indicator Source Audit}

The focused source-column audit now compiles as
\texttt{oneTermRobinGamma3BulkIndicatorSourceAudit\_n3}, with transcript
theorem
\texttt{oneTermRobinGamma3BulkIndicatorSourceAudit\_n3\_transcript}.  For the
focused \(n=3\) gamma3 entry, the source column is \(j=5\).  The paper's
\(U_{\mathrm{indic}}\) paragraph sets the indicator qubit to \(1\) for the
bulk window \(K_1 \leq j \leq K_2\).  In the active one-term Robin audit,
\(K_1=2\) and \(K_2=5\), so Lean records \(j=5\) as a bulk column.

\[
\begin{array}{c|ccc}
\text{endpoint} & \text{indicator} & \text{matches source bulk indicator} & \text{role} \\
\hline
228 & 1 & \text{true} & \text{full Fig.~\texttt{1 term ROBIN} endpoint} \\
84 & 0 & \text{false} & \text{clean Eq.~\texttt{ROBIN clarified} endpoint}
\end{array}
\]

This audit refines the source-contract gap.  The indicator \(1\) at endpoint
\(228\) is source-backed, not an artifact of the Lean encoding.  The remaining
missing ingredient is still the projection/register rule that relates the full
endpoint \(228\) to the clean endpoint \(84\).  The record keeps
\texttt{sourceStatesResetRule = false},
\texttt{conventionObligationProved = false}, and
\texttt{productSearchBlocked = true}.  Product-to-coefficient equality, LCU
composition, cleanup, unitarity, block projection, block correctness,
normalized equality, and final extraction remain false.

\subsection{Automation Rule}

Every future Lean change to this formalization must update at least one of:
\[
  \texttt{conversion-windows/QBE-AUTO-001.md},
  \qquad
  \texttt{paper-notes/GHL2025\_RobinOneTerm.tex},
  \qquad
  \texttt{proof-obligations/}.
\]
This keeps the Lean proof state, the paper notation, and the human explanation
synchronized.

\subsection{2026-05-26 Gamma3 Branch-Correct Transcript}

The previous focused audit used \(n=3\), \(K_1=2\), \(K_2=5\), and \(j=5\).
That value of \(j\) is a bulk column because \(K_1 \leq j \leq K_2\).  It
therefore belongs to the branch hidden by the \(+\dots\) in
Eq.~\texttt{ROBIN clarified}, not to the displayed boundary summand
\(0 \leq j < K_1\) or \(K_2 < j < 2^n\).

The Lean transcript now records the correction as
\texttt{oneTermRobinGamma3BranchCorrectSourceMap\_n3}, with theorem
\texttt{oneTermRobinGamma3BranchCorrectSourceMap\_n3\_transcript}.  The
boundary branch uses \(j=0\):
\[
\begin{array}{c|ccc}
\text{branch} & j & \text{indicator after }U_{\mathrm{indic}} & \text{Lean endpoint} \\
\hline
\text{displayed boundary} & 0 & 0 & 0 \\
\text{omitted bulk} & 5 & 1 & 228
\end{array}
\]
The clean endpoint \(84\) from the earlier slot-specific audit remains useful
bulk memory, but it is not a boundary endpoint for the displayed line.  The
fields
\texttt{oldEndpointMismatchWasBranchMismatch = true} and
\texttt{oldHumanProjectionFreezeSuperseded = true} make this correction
explicit.

No theorem-level semantic flag is promoted.  The boundary and bulk
product-to-coefficient obligations remain false, as do LCU composition, block
projection, block correctness, normalized equality, cleanup, unitarity, and
final extraction.  The next lower packet should target the boundary branch
first and may schedule the omitted bulk branch separately.

\subsection{2026-05-26 Gamma3 Boundary Product Interface}

The boundary product packet now compiles as
\texttt{oneTermRobinBlockEncodingProofRoute\_gamma3BoundaryProductToCoefficientInterface\_n3},
with transcript theorem
\texttt{oneTermRobinBlockEncodingProofRoute\_gamma3BoundaryProductToCoefficientInterface\_n3\_transcript}.
It specializes Eq.~\texttt{ROBIN clarified} to the displayed boundary branch
\(n=3\), \(j=0\), target entry \((0,0)\), and sparse slot \(s=2\).

The recorded ket-zero factors, in Fig.~\texttt{1 term ROBIN} order, are
\[
\begin{array}{c|c}
\text{gate} & \text{factor} \\
\hline
U_{\mathrm{indic}} & 1 \\
O_{D^T}^{S} & 1 \\
R_y^{\mathrm{boundary}} & \texttt{boundary\_cos\_half\_0\_2} \\
O_D^{BS} & 1 \\
O_f & \texttt{f\_3\_0}\cdot \texttt{N\_f\_inv} \\
\mathrm{SWAP} & 1 \\
(O_D^{BS})^\dagger & 1
\end{array}
\]
The exact seven-gate product equality is not promoted.  The new record keeps
\texttt{uniquePathSupportObligation.proved = false}; the next finite Lean
theorem must prove that all other evaluated paths from column \(32\) to row
\(32\) vanish before applying the reusable unique-path multiplication lemma.
Product-to-coefficient equality, projection-slot convention, boundary
normalizer semantics, LCU composition, block projection, block correctness,
normalized equality, cleanup, unitarity, and final extraction remain false.

\subsection{2026-05-26 Boundary Unique-Path Support Packet}

The next Lean target is a finite support theorem for the boundary product
interface.  The GHL2025 source gives the wavefunction route in
Eq.~\texttt{ROBIN clarified} and the gate order in Fig.~\texttt{1 term ROBIN},
but it does not contain a separate finite matrix-product proof.  The missing
ingredient is therefore a QBE-local classical Lean lemma.

For \(n=3\), \(j=0\), target entry \((0,0)\), and sparse slot \(s=2\), the
surviving ket-zero path is
\[
  32 \longrightarrow 32 \longrightarrow 32 \longrightarrow 32
  \longrightarrow 0 \longrightarrow 0 \longrightarrow 0
  \longrightarrow 32.
\]
The planned theorem
\texttt{oneTermRobinBlockEncodingProofRoute\_gamma3BoundaryUniquePathSupport\_n3}
should prove that every other evaluated contribution from column \(32\) to row
\(32\) is zero.  Only after that support theorem is available should lower work
apply \texttt{Matrix.evalWith\_mul\_unique\_path} to the boundary product
entry.  This packet does not promote product-to-coefficient equality,
projection-slot convention, boundary normalizer semantics, LCU composition,
block projection, block correctness, normalized equality, cleanup, unitarity,
or final extraction.

\subsection{2026-05-26 Boundary Support Audit Result}

The first support audit now compiles as
\texttt{oneTermRobinBlockEncodingProofRoute\_gamma3BoundaryUniquePathSupportAudit\_n3},
with theorem
\texttt{oneTermRobinBlockEncodingProofRoute\_gamma3BoundaryUniquePathSupport\_n3}.
It records the displayed boundary branch for \(n=3\), \(j=0\), target entry
\((0,0)\), and sparse slot \(s=2\).

The compiled state lists are
\[
\begin{array}{c|c}
\text{branch} & \text{state list} \\
\hline
\text{ket-zero} & [32,32,32,32,0,0,0,32] \\
\text{adjacent ket-one} & [32,32,32,33,1,1,1,33].
\end{array}
\]
The audit proves concrete zero entries for the adjacent ket-one contribution:
\[
  O_{D^T}^{S}[33,32] = 0,\qquad
  O_f[0,1] = 0,\qquad
  (O_D^{BS})^\dagger[32,1] = 0.
\]
These facts are not yet the full unique-path theorem.  The Lean record keeps
\texttt{supportComplete = false} and names the next missing interface as the
four-gate prefix support statement
\[
  (O_D^{BS} R_y^{\mathrm{boundary}} O_{D^T}^{S} U_{\mathrm{indic}})[k,32] = 0
\]
for all prefix rows \(k \notin \{0,1\}\).  This is a QBE-local classical Lean
lemma, not a cited oracle or LCU dependency.  Product-to-coefficient equality,
projection-slot convention, boundary normalizer semantics, LCU composition,
block projection, block correctness, normalized equality, cleanup, unitarity,
and final extraction remain false.

\subsection{2026-05-26 Boundary Prefix Support Result}

The prefix-support block now compiles.  The accepted theorem is
\[
\texttt{oneTermRobinBlockEncodingProofRoute\_gamma3BoundaryPrefixSupport\_n3}.
\]
It is an evaluated matrix statement for the four-gate prefix
\[
  O_D^{BS} R_y^{\mathrm{boundary}} O_{D^T}^{S} U_{\mathrm{indic}}.
\]
For every coefficient environment and every row outside \(\{0,1\}\), the
entry from source column \(32\) evaluates to zero:
\[
  \operatorname{eval}
  \bigl[
    (O_D^{BS} R_y^{\mathrm{boundary}} O_{D^T}^{S} U_{\mathrm{indic}})[k,32]
  \bigr] = 0.
\]
The proof uses the local matrix reducer
\texttt{Matrix.evalWith\_mul\_eq\_zero\_of\_all\_paths\_zero}, together with
the staged prefix declarations
\texttt{oneTermRobinGamma3BoundaryDUPrefixSupport\_n3} and
\texttt{oneTermRobinGamma3BoundaryRDUPrefixSupport\_n3}.  This is a
QBE-local finite matrix lemma.  It does not prove the raw symbolic product is
syntactically zero, and it does not promote product-to-coefficient equality,
projection-slot convention, boundary normalizer semantics, LCU composition,
block projection, block correctness, normalized equality, cleanup, unitarity,
or final extraction.

\subsection{2026-05-26 Seven-Gate Support Packet}

The next lower target is the suffix composition after the compiled prefix.
The GHL2025 source anchors, Eq.~\texttt{ROBIN clarified} and
Fig.~\texttt{1 term ROBIN}, specify the gamma3 wavefunction route and the gate
order, but they do not include a separate finite support proof for the row
\(32\), column \(32\) matrix entry.  The missing step is therefore a local
Lean lemma over the finite matrix backend.

The planned declarations are
\[
\begin{array}{l}
\texttt{oneTermRobinGamma3BoundarySuffixMatrix\_n3},\\
\texttt{oneTermRobinGamma3BoundarySevenGateMatrix\_n3},\\
\texttt{oneTermRobinBlockEncodingProofRoute\_gamma3BoundarySevenGateSupport\_n3}.
\end{array}
\]
The suffix matrix should represent
\[
  (O_D^{BS})^\dagger \,\mathrm{SWAP}\, O_f.
\]
The support theorem should show that in the product of this suffix with the
compiled prefix, every evaluated contribution into row \(32\) from column
\(32\) vanishes except the row-\(0\) prefix branch.  Rows outside
\(\{0,1\}\) should use the prefix-support theorem; the row-\(1\) branch should
use the compiled \(O_f\) and dagger zero entries already recorded by
\texttt{oneTermRobinBlockEncodingProofRoute\_gamma3BoundaryUniquePathSupport\_n3}.

Only after this support theorem compiles should the proof attempt apply
\texttt{Matrix.evalWith\_mul\_unique\_path} to the boundary product entry.
All theorem-level semantic flags remain false until an exact build-tested
theorem proves them.

\subsection{2026-05-26 Seven-Gate Support Result}

The seven-gate support block now compiles.  The accepted Lean declarations are
\[
\begin{array}{l}
\texttt{oneTermRobinGamma3BoundaryOfSwapMatrix\_n3},\\
\texttt{oneTermRobinGamma3BoundarySuffixMatrix\_n3},\\
\texttt{oneTermRobinGamma3BoundarySevenGateMatrix\_n3},\\
\texttt{oneTermRobinGamma3BoundaryOfSwapRow0Col1\_zero\_n3},\\
\texttt{oneTermRobinGamma3BoundarySuffixRow32Col1\_zero\_n3},\\
\texttt{oneTermRobinBlockEncodingProofRoute\_gamma3BoundarySevenGateSupport\_n3},\\
\texttt{oneTermRobinBlockEncodingProofRoute\_gamma3BoundarySevenGateUniquePath\_n3}.
\end{array}
\]
They apply to the displayed boundary branch of Eq.~\texttt{ROBIN clarified}
for \(n=3\), \(j=0\), target entry \((0,0)\), and sparse slot \(s=2\).
The support theorem proves that every evaluated suffix-prefix contribution
from column \(32\) to row \(32\) vanishes except the row-\(0\) branch.  Rows
outside \(\{0,1\}\) use the compiled prefix-support theorem; the row-\(1\)
branch uses the \(O_f\), SWAP, and dagger zero-support facts.

The unique-path theorem reduces the full seven-gate evaluated entry to the
row-\(0\) branch.  It does not evaluate that branch, and it does not prove the
gamma3 product-to-coefficient equality.  Projection-slot convention,
boundary-normalizer semantics, LCU composition, block projection, block
correctness, normalized equality, cleanup, unitarity, circuit unitarity, and
final extraction remain false obligations.

\subsection{2026-05-26 Boundary Product-Entry Packet}

The next lower packet is a local finite matrix evaluation.  It should use
\texttt{oneTermRobinBlockEncodingProofRoute\_gamma3BoundarySevenGateUniquePath\_n3}
to reduce the seven-gate entry to the row-\(0\) branch, then prove the prefix
and suffix factors:
\[
\begin{aligned}
  \operatorname{eval}\bigl[
    (O_D^{BS}R_y^{\mathrm{boundary}}O_{D^T}^{S}U_{\mathrm{indic}})[0,32]
  \bigr]
  &= \operatorname{eval}\bigl[\texttt{boundary\_cos\_half\_0\_2}\bigr],\\
  \operatorname{eval}\bigl[
    ((O_D^{BS})^\dagger\mathrm{SWAP}O_f)[32,0]
  \bigr]
  &= \operatorname{eval}\bigl[\texttt{f\_3\_0}\cdot\texttt{N\_f\_inv}\bigr].
\end{aligned}
\]
The intended product-entry theorem is
\[
\operatorname{eval}\bigl[U_{\mathrm{seven}}[32,32]\bigr]
=
\operatorname{eval}\bigl[
  (\texttt{f\_3\_0}\cdot\texttt{N\_f\_inv})
  \texttt{boundary\_cos\_half\_0\_2}
\bigr].
\]
This theorem is still below the paper-level coefficient claim: replacing
\texttt{boundary\_cos\_half\_0\_2} by the normalized derivative coefficient
uses the boundary-rotation analytic contract, which remains unproved in QBE.
The normalized \(A_k\) entry, sparse-slot projection, LCU composition, and
final block-extraction theorem remain explicit obligations.

\subsection{2026-05-26 Boundary Product-Entry Result and Coefficient Drift}

The boundary product-entry theorem now compiles as
\[
\texttt{oneTermRobinBlockEncodingProofRoute\_gamma3BoundaryProductEntryEval\_n3}.
\]
For the displayed boundary branch of Eq.~\texttt{ROBIN clarified} with
\(n=3\), \(j=0\), target entry \((0,0)\), sparse slot \(s=2\), source column
\(32\), and target row \(32\), Lean proves
\[
  \operatorname{eval}\bigl[U_{\mathrm{seven}}[32,32]\bigr]
  =
  \bigl(\operatorname{eval}\texttt{f\_3\_0}\bigr)
  \bigl(\operatorname{eval}\texttt{N\_f\_inv}\bigr)
  \operatorname{eval}\texttt{boundary\_cos\_half\_0\_2}.
\]
This theorem is a finite matrix-product calculation.  It uses the compiled
prefix support, suffix support, and unique-path reduction.  It does not prove
the coefficient identity in Eq.~\texttt{ROBIN clarified}.

The remaining blocker is a coefficient-source mismatch.  Lemma
\texttt{Banded-sparse-access-oracle} and the zero-inclusion paragraph before
Theorem~\texttt{1 term robin} use a global sparse slot \(s\) with
\[
  r_{si}=r_{s0}+i \pmod{2^n}.
\]
For the focused boundary branch, the active Lean table gives
\[
  \texttt{oneTermRobinGlobalSparseAddress 3 2 0 = 0}.
\]
Thus global sparse slot \(2\) selects the diagonal coefficient in row \(0\).
The current derivative and boundary-normalizer contracts instead use the
row-dependent helper
\[
  \texttt{robinSparseAmplitudeValue 3 2 0},
\]
which is the third row-local boundary entry, equal to \(-1/6\).  The diagonal
target coefficient is the row-local slot \(0\) value
\[
  -5/2 + (7/3)\,\mathcal{A}_1\Delta x.
\]
This is contract drift in the coefficient-source interface.  The next Lean
packet should introduce a global sparse-slot coefficient source and bridge it
to the \(O_{D^T}^{S}\) and \(R_y^{\mathrm{boundary}}\) normalizer routes before
any product-to-coefficient theorem is attempted.  All analytic normalizer,
projection, LCU, cleanup, unitarity, block-correctness, normalized-equality,
and final-extraction obligations remain false.

\subsection{2026-05-26 Global Sparse-Slot Coefficient Source}

The coefficient source is now defined before the next product-to-coefficient
claim.  The Lean declaration
\[
  \texttt{GHL2025.robinGlobalSparseAmplitudeValue}(n,s,i)
\]
uses the global sparse-slot table from
\(\texttt{GHL2025.oneTermRobinGlobalSparseOffset}\).  In this table, slot
\(2\) is the zero-offset diagonal slot.  Boundary slots that are present in
the global sparse register but absent from a boundary row are assigned
coefficient \(0\); the sparse slot remains part of the register domain.

For the displayed gamma3 boundary branch with \(n=3\), row \(0\), and sparse
slot \(s=2\), Lean proves
\[
  \texttt{GHL2025.robinGlobalSparseAmplitudeValue 3 2 0}
  =
  -5/2 + (7/3)\mathcal{A}_1\Delta x.
\]
It also proves that this value differs from
\(\texttt{GHL2025.robinSparseAmplitudeValue 3 2 0}\), the old row-local slot
value \(-1/6\).  A test checks matrix coherence against
\(\texttt{Examples.RobinHeat.robinDerivativeMatrix 3}\) at the column selected
by \(\texttt{oneTermRobinGlobalSparseAddress 3 2 0}\).

The typed normalizer contracts
\(\texttt{derivativeNormalizerNDContract}\),
\(\texttt{sparseAmplitudeOracleDTCoefficientNormalizerContract}\), and
\(\texttt{boundaryRotationAngleNormalizerContract}\) now read the global source.
The bridge
\[
  \texttt{GHL2025.robinGlobalSparseAmplitudeValue\_sharedNormalizerRoutes}
\]
records that the derivative and boundary normalizer routes share this source.
The bridge is only data wiring: the \(N_D\) bound, division semantics,
half-angle relation, Eq.~\texttt{ROBIN clarified} coefficient equality, LCU
composition, cleanup, unitarity, block correctness, normalized equality,
circuit unitarity, and final extraction remain false obligations.

\subsection{2026-05-26 Boundary \(R_y\) Angle-Convention Gap}

The compiled finite product theorem for the displayed boundary branch is
\texttt{oneTermRobinBlockEncodingProofRoute\_gamma3BoundaryProductEntryEval\_n3}.
It evaluates the row-\(32\), column-\(32\) seven-gate product as
\[
  f_{3,0}N_f^{-1}\,
  \texttt{boundary\_cos\_half\_0\_2}.
\]
The coefficient needed by Eq.~\texttt{eq: ROBIN clarified} is instead
\[
  f_{3,0}N_f^{-1}\,D_0^{(2)}N_D^{-1},
\]
where the focused global sparse-slot source is
\texttt{GHL2025.robinGlobalSparseAmplitudeValue 3 2 0}.

The remaining bridge is a source-contract gap.  Eq.~\texttt{eq:angles for Ry}
defines
\[
  \theta_j^s = \arccos(D_j^{(s)}/\mathcal{N}_D),
\]
whereas the Lean matrix entry
\texttt{boundary\_cos\_half\_0\_2} is the standard half-angle entry
\(\cos(\theta_0^2/2)\).  A standard half-angle identity gives a square-root
expression, not \(D_0^{(2)}/\mathcal{N}_D\).  The next Lean packet should
therefore add a focused false obligation, for example
\texttt{oneTermRobinGamma3BoundaryRyCoefficientBridge\_n3}, that records the
needed convention without promoting the product-to-coefficient theorem,
boundary-normalizer semantics, LCU composition, block correctness, or final
extraction.

\subsection{2026-05-26 Boundary \(R_y\) Coefficient Bridge}

The focused source-contract gap is now a Lean declaration:
\[
  \texttt{Examples.RobinHeat.oneTermRobinGamma3BoundaryRyCoefficientBridge\_n3}.
\]
For \(n=3\), system entry \((0,0)\), and global sparse slot \(s=2\), the
record stores the product factor
\[
  \texttt{boundary\_cos\_half\_0\_2}
\]
and the normalized coefficient
\[
  \texttt{GHL2025.boundaryRotationNormalizedCoefficient p 0 2}
  =
  \texttt{GHL2025.robinGlobalSparseAmplitudeValue 3 2 0}\,N_D^{-1},
\]
where \(p=\texttt{oneTermParameters 3}\).  The transcript theorem
\[
  \texttt{Examples.RobinHeat.oneTermRobinGamma3BoundaryRyCoefficientBridge\_n3\_transcript}
\]
checks this wiring and keeps the bridge obligation false.

This is not a theorem-level product-to-coefficient proof.  The boundary
division, arccos semantics, half-angle semantics, normalizer bound,
two-by-two unitarity, LCU composition, block projection, block correctness,
and final extraction remain explicit false obligations.

\subsection{2026-05-26 Boundary \(R_y\) Decision Packet}

The source audit now ends at a convention decision rather than a Lean product
calculation.  The paper defines
\[
  \theta_0^2=\arccos(D_0^{(2)}/\mathcal{N}_D)
\]
in Eq.~\texttt{eq:angles for Ry}.  The compiled boundary product entry for
the displayed branch of Eq.~\texttt{eq: ROBIN clarified} contains
\[
  f_{3,0}N_f^{-1}\,\texttt{boundary\_cos\_half\_0\_2},
\]
where \texttt{boundary\_cos\_half\_0\_2} is the standard
\(\cos(\theta_0^2/2)\) matrix entry.  The coefficient required by the
displayed formula is instead
\[
  f_{3,0}N_f^{-1}D_0^{(2)}N_D^{-1}.
\]

Lean records this boundary as
\[
  \texttt{Examples.RobinHeat.oneTermRobinGamma3BoundaryRyAngleConventionDecision\_n3}.
\]
Its transcript theorem checks that the source does not currently specify a
direct half-angle coefficient rule, that human or source input is required,
and that the product-to-coefficient search is blocked.  The bridge obligation,
boundary half-angle semantics, product-to-coefficient theorem, LCU
composition, block projection, block correctness, and final extraction remain
false.

There are only two faithful continuations.  One may keep the standard
\(R_y\) matrix convention and leave the displayed coefficient bridge as an
open theorem gap.  Alternatively, one may supply a source-backed boundary
amplitude convention explaining how the focused half-angle matrix entry
contributes \(D_0^{(2)}/\mathcal{N}_D\).  A tactic proof, a silent change to
the \(R_y\) matrix, or a return to the old bulk \(j=5\) boundary comparison is
not a faithful continuation.

\subsection{2026-05-26 Boundary \(R_y\) Lower-Packet Guard}

The lower-packet guard is now compiled as
\[
  \texttt{Examples.RobinHeat.oneTermRobinGamma3BoundaryRyLowerPacketGuard\_n3}.
\]
This record is not a new semantic claim.  It packages the decision packet so
that lower product-to-coefficient proof search stays disabled while the
boundary \(R_y\) convention is unresolved.

The transcript theorem
\[
  \texttt{Examples.RobinHeat.oneTermRobinGamma3BoundaryRyLowerPacketGuard\_n3\_transcript}
\]
checks the following guard state:
\[
\begin{aligned}
  \texttt{humanInputRequired} &= \texttt{true},\\
  \texttt{productSearchBlocked} &= \texttt{true},\\
  \texttt{lowerProductProofPacketAllowed} &= \texttt{false},\\
  \texttt{sourceBackedConventionPacketAllowed} &= \texttt{true}.
\end{aligned}
\]
It also checks that the bridge obligation, the decision obligation, the
boundary half-angle semantics, the gamma3 product-to-coefficient theorem, LCU
composition, block projection, block correctness, and final extraction remain
false.

The next faithful step is therefore a source-backed convention decision, or a
human/source decision to keep the standard \(R_y\) half-angle matrix entry and
leave the coefficient bridge open.  The guard forbids more local product
tactic search on
\(\texttt{oneTermRobinGamma3ProductToCoefficientObligation 3 0 0}\) until that
decision is present.

\subsection{2026-05-26 Source-Evidence Guard Sync}

The latest guard is also checked by a focused regression in
\texttt{Tests/Basic.lean}.  The test records the source-evidence boundary for
the unresolved \(R_y\) convention.  It checks that
\[
  \texttt{lowerProductProofPacketAllowed}=\texttt{false}
  \quad\text{and}\quad
  \texttt{sourceBackedConventionPacketAllowed}=\texttt{true}.
\]
It also checks that the accepted source-backed options are limited to keeping
the standard \(R_y\) half-angle convention with the coefficient bridge left
open, or supplying a paper-backed boundary amplitude convention.

Middle rechecked Eq.~\texttt{eq:angles for Ry}, Eq.~\texttt{eq: ROBIN
clarified}, Fig.~\texttt{fig:1 term ROBIN}, and
Theorem~\texttt{theorem: 1 term robin}.  The classification is unchanged.
The finite product entry and the global sparse-slot coefficient source are
compiled, but the paper source does not state a direct rule identifying
\(\cos(\theta_0^2/2)\) with \(D_0^{(2)}/\mathcal{N}_D\).  Therefore
\texttt{oneTermRobinGamma3ProductToCoefficientObligation 3 0 0} remains
blocked by a source-contract gap, not by another matrix-product lemma.

\subsection{2026-05-26 Persistent Boundary \(R_y\) Gap Closeout}

Middle rechecked the source neighborhood for Eq.~\texttt{eq:angles for Ry},
Eq.~\texttt{eq: ROBIN clarified}, Fig.~\texttt{fig:1 term ROBIN}, and
Theorem~\texttt{theorem: 1 term robin}.  The source defines
\[
  \theta_0^2=\arccos(D_0^{(2)}/\mathcal{N}_D).
\]
The compiled seven-gate boundary product remains
\[
  f_{3,0}N_f^{-1}\,
  \texttt{boundary\_cos\_half\_0\_2},
\]
where \(\texttt{boundary\_cos\_half\_0\_2}\) is the standard
\(\cos(\theta_0^2/2)\) matrix entry.  The displayed gamma3 coefficient instead
requires
\[
  f_{3,0}N_f^{-1}D_0^{(2)}N_D^{-1}.
\]
The audited source does not state a rule identifying these two boundary
factors.

The Lean closeout is therefore a guard state, not a proof of the displayed
coefficient.  The declarations
\[
  \texttt{oneTermRobinGamma3BoundaryRyCoefficientBridge\_n3},\qquad
  \texttt{oneTermRobinGamma3BoundaryRyAngleConventionDecision\_n3},\qquad
  \texttt{oneTermRobinGamma3BoundaryRyLowerPacketGuard\_n3}
\]
record the missing convention and keep the bridge, decision,
product-to-coefficient, LCU, block projection, block correctness, and final
extraction flags false.

The next faithful step is one of three actions: provide a source-backed
boundary amplitude convention, record a human or source decision to keep the
standard \(R_y\) matrix entry and leave the coefficient bridge open, or perform
a reviewer closeout with all semantic flags still false.  More local product
search on
\[
  \texttt{oneTermRobinGamma3ProductToCoefficientObligation 3 0 0}
\]
is not an accepted proof route from this state.

\subsection{2026-05-26 Corrected-Angle Boundary Interface}

The source audit now treats the displayed boundary angle line as a
source-backed correction route.  Eq.~\texttt{eq:angles for Ry} writes
\[
  \theta_0^2=\arccos(D_0^{(2)}/\mathcal{N}_D),
\]
but the standard \(R_y(\theta)\) convention has clean entry
\(\cos(\theta/2)\).  The prior sparse-amplitude construction and the companion
Robin heat implementation use the factor-two convention, so the Lean route
records
\[
  \theta_0^2=2\arccos(D_0^{(2)}/\mathcal{N}_D)
\]
as the corrected faithful continuation.

The compiled declarations
\[
  \texttt{oneTermRobinGamma3BoundaryRyCorrectedAngleSourceDecision\_n3},
  \qquad
  \texttt{oneTermRobinGamma3BoundaryCorrectedCoefficientInterface\_n3}
\]
record the correction and the conditional coefficient interface.  The theorem
\[
  \texttt{oneTermRobinBlockEncodingProofRoute\_gamma3BoundaryProductEntryEval\_correctedAngle\_n3}
\]
rewrites the finite seven-gate product using
\(\texttt{boundaryRotationNormalizedCoefficient (oneTermParameters 3) 0 2}\)
under the corrected-entry hypothesis.

The hypothesis remains an explicit obligation.  The focused
\(\texttt{oneTermRobinGamma3ProductToCoefficientObligation 3 0 0}\), LCU
composition, block projection, cleanup, unitarity, block correctness,
normalized equality, circuit unitarity, and final extraction all remain false.

\subsection{2026-05-27 Boundary Normalizer Split Target}

The boundary product route now has a compiled normalizer/projection packet
\[
  \texttt{oneTermRobinGamma3BoundaryNormalizerProjectionConvention\_n3}.
\]
Middle split this packet into
\[
  \texttt{oneTermRobinGamma3BoundaryNormalizerSplitTarget\_n3},
\]
which keeps the target
\[
  \texttt{oneTermRobinGamma3ProductToCoefficientObligation 3 0 0}
\]
fixed.  The split is:
\[
  \underbrace{N_D^{-1}N_f^{-1}}_{\text{symbolic inverse convention}}
  \qquad\text{and}\qquad
  \underbrace{\kappa^{-1}}_{\text{sparse-register projection factor}}.
\]

The symbolic inverse obligation is
\(\texttt{symbolicInverseObligation.proved = false}\).  The sparse-register
projection obligation is
\(\texttt{kappaProjectionObligation.proved = false}\).  The finite normalized
block equality and the product-to-coefficient obligation also remain false.
This packet only fixes the next lower target; it does not promote LCU
composition, block projection, block correctness, normalized equality, circuit
unitarity, or final extraction.

\subsection{2026-05-27 Boundary \(\kappa\)-Projection Source Contract}

The sparse-register part of the boundary product route now has a compiled
conditional algebra lemma
\[
  \texttt{oneTermRobinGamma3BoundaryKappaProjectionEval\_n3}.
\]
For a coefficient environment satisfying
\[
  \mathrm{env}(N_D^{-1})\mathrm{env}(N_D)=1,\qquad
  \mathrm{env}(N_f^{-1})\mathrm{env}(N_f)=1,\qquad
  \mathrm{env}(\kappa^{-1})\mathrm{env}(\kappa)=1,
\]
the lemma checks the cancellation
\[
  \mathrm{eval}(\texttt{branchLocalProduct}\cdot\kappa^{-1})
  \mathrm{eval}(N_DN_f\kappa)
  =
  \mathrm{eval}(\texttt{targetEntry}).
\]
This is coefficient algebra only.  It does not prove that the circuit supplies
the factor \(\kappa^{-1}\).

The paper source for that factor is Eq.~\texttt{arbitrary sparcity}, where
\[
  H_W^{(\kappa)}\ket{0}^{\lceil\log_2\kappa\rceil}
  =
  \frac{1}{\sqrt{\kappa}}
  \sum_{s=0}^{\kappa-1}\ket{s}^{\lceil\log_2\kappa\rceil}.
\]
GHL2025 cites Shukla--Vedula 2024 for the implementation of this uniform
superposition.  The matching block projection onto the same sparse slot should
contribute the second factor \(1/\sqrt{\kappa}\).  In Lean this source contract
is recorded by
\[
  \texttt{oneTermRobinGamma3BoundaryProjectionSourceContract\_n3}
\]
and its transcript theorem.  The focused data are \(\kappa=7\), sparse slot
\(2\), and clean basis index \(32\).

The remaining obligations are explicit:
\[
\begin{array}{lll}
\text{uniform sparse-register preparation} &
\texttt{uniformPreparationObligation.proved} & \text{false}\\
\text{matching sparse-register projection} &
\texttt{matchingProjectionObligation.proved} & \text{false}\\
\text{identification with }\texttt{Coeff.symbol "kappa\_inv"} &
\texttt{projectionFactorSemantics.proved} & \text{false}\\
\text{finite normalized block equality} &
\texttt{finiteCompositionNormalizedEquality.proved} & \text{false}\\
\text{focused product-to-coefficient theorem} &
\texttt{productObligation.proved} & \text{false}.
\end{array}
\]
The next lower target is the finite projection-factor semantics for the fixed
boundary route.  It should use the compiled source contract above and must not
promote product-to-coefficient, LCU composition, block projection, block
correctness, normalized equality, circuit unitarity, or final extraction.

\subsection{2026-05-27 Boundary Projection Factor Interface}

The finite projection-factor interface now compiles as
\[
  \texttt{oneTermRobinGamma3BoundaryProjectionFactorSemantics\_n3}.
\]
It reuses the source contract for the inserted
\(\texttt{kappa\_inv}\) factor and adds the finite index lemma
\[
  \texttt{oneTermRobinGamma3BoundaryProjectionFactorIndex\_n3}.
\]
The index lemma proves only finite bookkeeping: the prepared sparse slot and
the projected sparse slot are both slot \(2\), and both sides use the clean
basis index \(32\).

The semantic content is still an obligation.  The paper source says that
\(H_W^{(\kappa)}\) prepares each sparse slot with amplitude
\(1/\sqrt{\kappa}\), and the matching block projection onto the same slot
should contribute another \(1/\sqrt{\kappa}\).  QBE still has to prove or
contract-map that their product is the Lean factor
\[
  \texttt{Coeff.symbol "kappa\_inv"}.
\]

The current status is:
\[
\begin{array}{lll}
\text{finite slot and basis agreement} &
\texttt{oneTermRobinGamma3BoundaryProjectionFactorIndex\_n3} &
\text{compiled}\\
\text{projection-factor interface} &
\texttt{oneTermRobinGamma3BoundaryProjectionFactorSemantics\_n3} &
\text{compiled; semantic fields false}\\
\text{uniform sparse-register preparation} &
\texttt{uniformPreparationObligation.proved} & \text{false}\\
\text{matching sparse-register projection} &
\texttt{matchingProjectionObligation.proved} & \text{false}\\
\text{identification with }\texttt{Coeff.symbol "kappa\_inv"} &
\texttt{factorSemanticsObligation.proved} & \text{false}\\
\text{focused product-to-coefficient theorem} &
\texttt{productObligation.proved} & \text{false}.
\end{array}
\]
Thus the next lower target remains
\[
  \texttt{oneTermRobinGamma3ProductToCoefficientObligation 3 0 0},
\]
routed through the projection-factor interface.  The LCU composition, block
projection, block correctness, normalized equality, circuit unitarity, and
final extraction flags remain false.

\subsection{2026-05-27 Projection-Factor Obstruction Split}

The focused boundary route now has a compiled obstruction packet:
\[
  \texttt{oneTermRobinGamma3BoundaryProjectionFactorObstruction\_n3}.
\]
It separates the remaining \(\texttt{kappa\_inv}\) statement into two inputs.
The first input is the cited \(H_W^{(\kappa)}\) per-slot amplitude from
GHL2025 Eq.~\texttt{arbitrary sparcity} and the Shukla--Vedula uniform
superposition result.  The second input is QBE's local matching projection
onto the same sparse slot \(2\), as required by Definition~\texttt{def:block-encoding}.

The current status is:
\[
\begin{array}{lll}
\text{obstruction packet} &
\texttt{oneTermRobinGamma3BoundaryProjectionFactorObstruction\_n3} &
\text{compiled; semantic fields false}\\
\text{cited uniform preparation} &
\texttt{uniformPreparationObligation.proved} & \text{false}\\
\text{matching sparse-slot projection} &
\texttt{matchingProjectionObligation.proved} & \text{false}\\
\text{identification with }\texttt{Coeff.symbol "kappa\_inv"} &
\texttt{factorSemanticsObligation.proved} & \text{false}\\
\text{focused product-to-coefficient theorem} &
\texttt{productObligation.proved} & \text{false}.
\end{array}
\]
The next lower packet should state the local matching-projection convention,
not recursively formalize the cited state-preparation paper.  It should reuse
the obstruction packet, the finite slot-index lemma, and the conditional
\(\texttt{kappa\_inv}\) cancellation lemma, and it must keep LCU composition,
block projection, block correctness, normalized equality, circuit unitarity,
and final extraction false.

\subsection{2026-05-27 Matching-Projection Amplitude Obstruction}

The matching-projection convention has now been refined to the compiled
obstruction packet
\[
  \texttt{oneTermRobinGamma3BoundaryMatchingProjectionAmplitudeObstruction\_n3}.
\]
This packet keeps the same focused boundary branch, sparse slot \(2\), and
clean basis index \(32\).  It separates the ket-side amplitude supplied by the
cited \(H_W^{(\kappa)}\) preparation from the bra-side matching projection
amplitude required by the block-encoding definition.

The packet also names the local coefficient-algebra lemma
\[
  \texttt{oneTermRobinGamma3BoundaryProjectionFactorProductEval\_n3}.
\]
That lemma says only that, under an explicit coefficient-environment
hypothesis, the product of two symbols
\(\texttt{sqrt\_kappa\_inv}\) evaluates as
\(\texttt{kappa\_inv}\).  It does not prove either sparse-register amplitude.

The current status is:
\[
\begin{array}{lll}
\text{matching amplitude obstruction} &
\texttt{oneTermRobinGamma3BoundaryMatchingProjectionAmplitudeObstruction\_n3} &
\text{compiled; semantic fields false}\\
\text{ket-side uniform preparation} &
\texttt{uniformPreparationObligation.proved} & \text{false}\\
\text{bra-side matching projection} &
\texttt{matchingProjectionObligation.proved} & \text{false}\\
\text{identification with }\texttt{Coeff.symbol "kappa\_inv"} &
\texttt{factorSemanticsObligation.proved} & \text{false}\\
\text{focused product-to-coefficient theorem} &
\texttt{productObligation.proved} & \text{false}.
\end{array}
\]
The next lower packet should prove or contract-map the bra-side matching
projection amplitude for sparse slot \(2\), starting from this obstruction
packet.  It should not recursively formalize the Shukla--Vedula state
preparation paper and should not promote LCU composition, block projection,
block correctness, normalized equality, circuit unitarity, or final extraction.

\subsection{Projection-Amplitude Semantics Contract}

Lean now defines the Phase-1 semantics contract
\[
  \texttt{oneTermRobinGamma3BoundaryProjectionAmplitudeSemantics\_n3}.
\]
This is the current source-facing interface for the focused boundary
\(\gamma_3\) route.  It accepts the ket-side \(H_W^{(\kappa)}\) amplitude and
the bra-side matching projection amplitude only as explicit contracts, both
represented by \(\texttt{Coeff.symbol "sqrt\_kappa\_inv"}\).

The accompanying theorem
\[
  \texttt{oneTermRobinGamma3BoundaryProjectionAmplitudeContractProductEval\_n3}
\]
is only a coefficient-algebra lemma.  It assumes
\[
  \texttt{env "sqrt\_kappa\_inv" * env "sqrt\_kappa\_inv" = env "kappa\_inv"}
\]
and then rewrites the symbolic product of the accepted ket and bra factors to
\(\texttt{Coeff.symbol "kappa\_inv"}\).  It does not prove the cited
Shukla--Vedula uniform-state preparation theorem, the QBE matching projection
amplitude, or the final product-to-coefficient statement.

The current proof obligations are:
\[
\begin{array}{lll}
\text{ket-side uniform preparation} &
\texttt{uniformPreparationObligation.proved} & \text{false}\\
\text{bra-side matching projection amplitude} &
\texttt{braAmplitudeObligation.proved} & \text{false}\\
\text{factor semantics for }\texttt{kappa\_inv} &
\texttt{factorSemanticsObligation.proved} & \text{false}\\
\text{finite normalized block equality} &
\texttt{finiteCompositionNormalizedEquality.proved} & \text{false}\\
\text{focused product-to-coefficient theorem} &
\texttt{productObligation.proved} & \text{false}.
\end{array}
\]

The next Lean packet should target the factor-semantics obligation for sparse
slot \(2\), using the compiled transcript theorem, the conditional product
lemma, and \(\texttt{oneTermRobinGamma3BoundaryKappaProjectionEval\_n3}\).
No semantic flag should be promoted until a build-tested Lean theorem proves
that exact field.

\subsection{2026-05-27 Clean-Column To Bra-Route Contract}

The focused route now has a compiled clean-column-to-bra contract:
\[
  \texttt{oneTermRobinGamma3BoundaryCleanColumnBraRouteContract\_n3}.
\]
It connects the contract-only clean-column bridge for
\[
  H_W^{(\kappa)}[2,0] = 1/\sqrt{\kappa}
\]
to both active bra-side obligations:
\[
  \texttt{oneTermRobinGamma3BoundaryBraProjectionAmplitudeSourceMap\_n3}
  \quad\text{and}\quad
  \texttt{oneTermRobinGamma3BoundaryFactorSemanticsContractMap\_n3}.
\]
The theorem
\[
  \texttt{oneTermRobinGamma3BoundaryCleanColumnBraRouteContract\_feedsBraAmplitude\_n3}
\]
is conditional on the external Shukla--Vedula clean-column input.  It rewrites
that input through the compiled transpose bridge to the expected
\(\texttt{sqrt\_kappa\_inv}\) bra factor, but it does not prove the
state-preparation theorem, the QBE projection entry, or the product route.

\[
\begin{array}{lll}
\text{route contract} &
\texttt{oneTermRobinGamma3BoundaryCleanColumnBraRouteContract\_n3} &
\text{compiled; semantic fields false}\\
\text{conditional bra-factor theorem} &
\texttt{oneTermRobinGamma3BoundaryCleanColumnBraRouteContract\_feedsBraAmplitude\_n3} &
\text{compiled under external clean-column hypothesis}\\
\text{bra-side source map} &
\texttt{sourceMapBraAmplitudeObligation.proved} & \text{false}\\
\text{factor-map bra field} &
\texttt{factorMapBraAmplitudeObligation.proved} & \text{false}\\
\text{factor semantics} &
\texttt{factorSemanticsObligation.proved} & \text{false}\\
\text{focused product theorem} &
\texttt{oneTermRobinGamma3ProductToCoefficientObligation\ 3\ 0\ 0} &
\text{false}.
\end{array}
\]

\subsection{2026-05-27 Next Factor-Semantics Route Packet}

The next Lean packet should use
\(\texttt{oneTermRobinGamma3BoundaryCleanColumnBraRouteContract\_n3}\) as the
bra-factor input for
\(\texttt{oneTermRobinGamma3BoundaryFactorSemanticsContractMapEval\_n3}\).
It should be an under-contract bridge, not a proof of the cited
Shukla--Vedula state-preparation circuit and not a proof of standard LCU
composition.  The packet must keep the following fields false until Lean
proves them directly or records a stricter typed contract:
\[
\begin{array}{lll}
\text{external clean-column amplitude} & \texttt{uniformColumnObligation} & \text{false}\\
\text{ket-side sparse-register amplitude} & \texttt{ketAmplitudeObligation} & \text{false}\\
\text{square-root product identity} & \texttt{productHypothesisObligation} & \text{false}\\
\text{finite normalized block equality} & \texttt{finiteCompositionNormalizedEquality} & \text{false}\\
\text{focused product equality} & \texttt{productObligation} & \text{false}.
\end{array}
\]

\subsection{2026-05-27 Clean-Column To Factor-Semantics Route}

The focused route now has an under-contract factor-semantics bridge:
\[
  \texttt{oneTermRobinGamma3BoundaryCleanColumnFactorSemanticsRoute\_n3}.
\]
It uses
\[
  \texttt{oneTermRobinGamma3BoundaryCleanColumnBraRouteContract\_n3}
\]
as the bra-factor input for
\[
  \texttt{oneTermRobinGamma3BoundaryFactorSemanticsContractMapEval\_n3}.
\]
The Lean theorem
\[
  \texttt{oneTermRobinGamma3BoundaryCleanColumnFactorSemanticsRouteEval\_n3}
\]
is conditional on the external clean-column amplitude
\[
  H_W^{(\kappa)}[2,0] = 1/\sqrt{\kappa}
\]
and on the coefficient-environment identities for
\(\texttt{N\_D\_inv}\), \(\texttt{N\_f\_inv}\),
\(\texttt{kappa\_inv}\), and \(\texttt{sqrt\_kappa\_inv}\).  It therefore
compiles the route from the clean bra factor into the factor-semantics
calculation without proving the cited state-preparation theorem, finite LCU
composition, or the focused product-to-coefficient theorem.

\[
\begin{array}{lll}
\text{clean-column source contract} & \texttt{uniformColumnObligation} & \text{false}\\
\text{ket-side sparse-register amplitude} & \texttt{ketAmplitudeObligation} & \text{false}\\
\text{bra-factor route} & \texttt{conditionalBraFactorCompiled} & \text{true}\\
\text{factor-map evaluation} & \texttt{conditionalFactorEvalCompiled} & \text{true}\\
\text{factor semantics} & \texttt{factorSemanticsObligation.proved} & \text{false}\\
\text{finite normalized block equality} & \texttt{finiteCompositionNormalizedEquality.proved} & \text{false}\\
\text{focused product theorem} & \texttt{productObligation.proved} & \text{false}.
\end{array}
\]

\subsection{2026-05-27 Product Under Contracts Route}

The focused boundary branch now has a product-under-contracts route:
\[
  \texttt{oneTermRobinGamma3BoundaryProductUnderContractsRoute\_n3}.
\]
It consumes
\[
  \texttt{oneTermRobinGamma3BoundaryCleanColumnFactorSemanticsRoute\_n3}
\]
and points to the fixed theorem-facing obligation
\[
  \texttt{oneTermRobinGamma3ProductToCoefficientObligation\ 3\ 0\ 0}.
\]
The theorem
\[
  \texttt{oneTermRobinGamma3BoundaryProductUnderContractsEval\_n3}
\]
compiles only the local coefficient calculation under the external
clean-column amplitude and the coefficient-environment identities for
\(\texttt{N\_D\_inv}\), \(\texttt{N\_f\_inv}\),
\(\texttt{kappa\_inv}\), and \(\texttt{sqrt\_kappa\_inv}\).  It does not
prove Shukla--Vedula state preparation, finite LCU composition, block
projection, or the product-to-coefficient theorem.

\[
\begin{array}{lll}
\text{route object} & \texttt{oneTermRobinGamma3BoundaryProductUnderContractsRoute\_n3} & \text{compiled}\\
\text{conditional coefficient theorem} & \texttt{oneTermRobinGamma3BoundaryProductUnderContractsEval\_n3} & \text{compiled}\\
\text{projection/product bridge} & \texttt{productBridgeObligation.proved} & \text{false}\\
\text{finite normalized block equality} & \texttt{finiteCompositionNormalizedEquality.proved} & \text{false}\\
\text{fixed product theorem} & \texttt{productObligation.proved} & \text{false}.
\end{array}
\]

\subsection{2026-05-27 Finite Projection/Product Bridge}

The finite bridge now compiles as
\[
  \texttt{oneTermRobinGamma3BoundaryFiniteProjectionProductBridge\_n3}.
\]
It records the Definition~\texttt{def:block-encoding} index fact for the
focused system entry \((0,0)\): the signal-zero block entry is the full matrix
entry \([0,0]\).  The branch-correct slot from Eq.~\texttt{ROBIN clarified}
is still sparse slot \(2\), whose embedded basis index is \(32\).  Therefore
the branch-local product at \([32,32]\) still needs a separate
branch-decomposition/projection theorem before it can be identified with the
block entry.

\[
\begin{array}{lll}
\text{finite index lemma} & \texttt{oneTermRobinGamma3BoundaryFiniteProjectionBlockEntryIndex\_n3} & \text{compiled}\\
\text{finite bridge} & \texttt{oneTermRobinGamma3BoundaryFiniteProjectionProductBridge\_n3} & \text{compiled; flags false}\\
\text{transcript theorem} & \texttt{oneTermRobinGamma3BoundaryFiniteProjectionProductBridge\_n3\_transcript} & \text{compiled}\\
\text{branch equals block index} & \texttt{branchBasisMatchesSignalBlockIndex} & \text{false}\\
\text{branch-decomposition theorem} & \texttt{branchDecompositionObligation.proved} & \text{false}\\
\text{product bridge} & \texttt{productBridgeObligation.proved} & \text{false}\\
\text{fixed product theorem} & \texttt{oneTermRobinGamma3ProductToCoefficientObligation\ 3\ 0\ 0} & \text{false}.
\end{array}
\]

The next Lean target should be the branch-decomposition/projection interface
for slot \(2\).  Shukla--Vedula and LCU remain contract-only dependencies, and
no downstream block-encoding flag should be promoted until an exact
build-tested theorem supplies the missing field.

\subsection{2026-05-28 Branch-Decomposition Packet}

The active theorem-map object remains
\[
  \texttt{oneTermRobinGamma3BoundaryFiniteProjectionProductBridge\_n3}.
\]
Middle re-audited Theorem~\texttt{1 term robin},
Eq.~\texttt{ROBIN clarified}, Fig.~\texttt{1 term ROBIN}, and
Definition~\texttt{def:block-encoding}.  The source fixes the focused boundary
branch data for system entry \((0,0)\) and sparse slot \(2\), while the
block-encoding definition fixes the signal-zero block entry.  It does not
state the finite QBE matrix theorem that expands the signal-zero block entry
as a branch sum and selects the slot-\(2\) contribution from the full branch
entry \([32,32]\).

The next lower packet should introduce a typed interface such as
\[
  \texttt{OneTermRobinGamma3BoundaryBranchDecompositionSlot2},\qquad
  \texttt{oneTermRobinGamma3BoundaryBranchDecompositionSlot2\_n3},
\]
or a smaller obstruction record.  It must consume the finite bridge, preserve
the compiled facts \([0,0]\) and \([32,32]\), and keep
\[
\begin{array}{lll}
\text{branch decomposition} & \texttt{branchDecompositionObligation} & \text{false}\\
\text{product bridge} & \texttt{productBridgeObligation} & \text{false}\\
\text{finite normalized equality} & \texttt{finiteCompositionNormalizedEquality} & \text{false}\\
\text{focused product theorem} & \texttt{oneTermRobinGamma3ProductToCoefficientObligation\ 3\ 0\ 0} & \text{false}.
\end{array}
\]

No Shukla--Vedula, LCU, block-projection, block-correctness, circuit-unitarity,
or final-extraction flag is promoted by this packet.

\subsection{2026-05-28 Branch-Decomposition Interface}

The branch-decomposition packet now compiles as
\[
  \texttt{oneTermRobinGamma3BoundaryBranchDecompositionSlot2\_n3}.
\]
It consumes the finite bridge
\[
  \texttt{oneTermRobinGamma3BoundaryFiniteProjectionProductBridge\_n3}
\]
and records the exact source-to-Lean indices for the displayed boundary branch
of Eq.~\texttt{ROBIN clarified}.  The focused system entry is \((0,0)\), the
sparse slot is \(2\), the signal-zero block entry from
Definition~\texttt{def:block-encoding} is \([0,0]\), and the branch-local
product entry is \([32,32]\).

The transcript theorem
\[
  \texttt{oneTermRobinGamma3BoundaryBranchDecompositionSlot2\_n3\_transcript}
\]
checks that these two entries are still distinct:
\[
\begin{array}{lll}
\text{signal-zero block entry} & [0,0] & \text{compiled}\\
\text{slot-2 branch entry} & [32,32] & \text{compiled}\\
\text{branch equals block entry} & \texttt{signalBlockEntryMatchesBranchEntry} & \text{false}\\
\text{projection summation} & \texttt{projectionSummationObligation.proved} & \text{false}\\
\text{branch decomposition} & \texttt{branchDecompositionObligation.proved} & \text{false}\\
\text{fixed product theorem} & \texttt{oneTermRobinGamma3ProductToCoefficientObligation\ 3\ 0\ 0} & \text{false}.
\end{array}
\]

The remaining proof obligation is therefore internal QBE finite matrix
bookkeeping: state and prove the projection-summation theorem that expands the
signal-zero block entry as a sparse-branch sum and selects the slot-\(2\)
contribution.  This packet does not prove Shukla--Vedula state preparation,
standard LCU composition, normalized block equality, block correctness,
circuit unitarity, or final extraction.

\subsection{2026-05-28 Typed Projection-Summation Target}

The typed projection-summation target now compiles as
\[
  \texttt{oneTermRobinGamma3BoundaryProjectionSummationTarget\_n3}.
\]
It refines the branch-decomposition interface by naming the actual coefficient
terms that the next theorem must relate.  Definition~\texttt{def:block-encoding}
selects the signal-zero block entry, while Eq.~\texttt{ROBIN clarified} and
Fig.~\texttt{1 term ROBIN} keep the focused boundary contribution at sparse
slot \(2\).

The compiled theorem
\[
  \texttt{oneTermRobinGamma3BoundaryProjectionSummationTarget\_blockEntry\_eq\_unitary\_n3}
\]
proves only the finite index equality between the block entry and the full
unitary entry at \([0,0]\).  The branch-local product remains a distinct typed
entry at \([32,32]\).

\[
\begin{array}{lll}
\text{signal block entry} & \texttt{signalBlockEntry} & [0,0]\ \text{typed}\\
\text{full unitary entry} & \texttt{signalUnitaryEntry} & [0,0]\ \text{proved equal to the signal block entry}\\
\text{branch matrix entry} & \texttt{branchMatrixEntry} & [32,32]\ \text{typed}\\
\text{branch-entry selection} & \texttt{branchEntrySelectionObligation.proved} & \text{false}\\
\text{projection summation} & \texttt{projectionSummationObligation.proved} & \text{false}\\
\text{fixed product theorem} & \texttt{oneTermRobinGamma3ProductToCoefficientObligation\ 3\ 0\ 0} & \text{false}.
\end{array}
\]

The next lower proof should first connect the branch matrix entry at
\([32,32]\) to the route's projected branch product under the existing
coefficient contracts.  Only after that bridge is typed should it prove, or
further refine, the finite projection-summation theorem selecting the slot-\(2\)
contribution in the signal-zero block entry.  Shukla--Vedula uniform
preparation and standard LCU composition remain contract-only dependencies.

\subsection{2026-05-28 Projection-Summation Obstruction}

The focused projection bridge has now been sharpened to the compiled
obstruction
\[
  \texttt{oneTermRobinGamma3BoundaryProjectionSummationObstruction\_n3}.
\]
It starts from the block-entry target above and records the sparse-slot domain
\(\{0,\ldots,6\}\) with focused slot \(2\).  The selected slot contribution is
the coefficient
\[
  \texttt{oneTermRobinGamma3BoundarySevenGateMatrix\_n3[32,32]}
  \cdot \texttt{sqrt\_kappa\_inv}\cdot \texttt{sqrt\_kappa\_inv}.
\]

The theorem
\[
  \texttt{oneTermRobinGamma3BoundaryProjectionSummationObstruction\_selectedSlotEval\_n3}
\]
reuses the branch-entry selection calculation: under the corrected
\(\texttt{Ry\_boundary}\) entry hypothesis, the selected slot contribution
evaluates to the route's projected branch product.  This is a conditional
local coefficient step, not the block-projection theorem.

The remaining missing interface is a finite branch-contribution family for the
signal-zero block entry:
\[
  \texttt{branchContribution : Fin\ 7 -> Coeff}.
\]
The next Lean target should prove or type as an obstruction the two statements
\[
  \texttt{branchContribution[2] = selectedSlotContribution}
  \qquad\text{and}\qquad
  \texttt{signalBlockEntry = Finset.univ.sum branchContribution}.
\]
Until those statements and the finite normalized block equality are proved,
\(\texttt{oneTermRobinGamma3ProductToCoefficientObligation\ 3\ 0\ 0}\)
remains false.  No Shukla--Vedula state-preparation theorem, LCU theorem,
unitarity claim, block-projection flag, block-correctness flag, or final
extraction flag is promoted by this obstruction.

\subsection{2026-05-28 Branch-Contribution Family Interface}

The branch-contribution interface now compiles as
\[
  \texttt{oneTermRobinGamma3BoundaryBranchContributionFamily\_n3}.
\]
It introduces a typed family
\[
  \texttt{branchContribution : Fin\ 7 -> Coeff}
\]
and the project-local fold
\[
  \texttt{oneTermRobinGamma3BoundaryBranchContributionSum}.
\]
The selected-slot theorem
\[
  \texttt{oneTermRobinGamma3BoundaryBranchContribution\_selectedSlot\_n3}
\]
proves only that slot \(2\) is the previously accepted selected contribution.
The signal-zero branch-sum theorem remains absent.

\[
\begin{array}{lll}
\text{branch family} &
\texttt{oneTermRobinGamma3BoundaryBranchContributionFamily\_n3} &
\text{compiled}\\
\text{selected slot} &
\texttt{oneTermRobinGamma3BoundaryBranchContribution\_selectedSlot\_n3} &
\text{compiled}\\
\text{branch sum} &
\texttt{oneTermRobinGamma3BoundaryBranchContribution\_sum\_n3} &
\text{absent}\\
\text{fixed product theorem} &
\texttt{oneTermRobinGamma3ProductToCoefficientObligation\ 3\ 0\ 0} &
\text{false}.
\end{array}
\]

\subsection{2026-05-28 Backend Projection-Summation Field Target}

The backend target now compiles as
\[
  \texttt{oneTermRobinGamma3BoundaryBackendProjectionSummationFieldTarget\_n3}.
\]
It records the smallest remaining finite projection obstruction.  The
placeholder branch family is not backend-sourced and may not be used to prove
the signal-zero branch sum.

The acceptance predicate is
\[
  \texttt{oneTermRobinGamma3BoundaryBackendBranchContributionPredicate\_n3}.
\]
A backend-sourced family must select slot \(2\) as the accepted contribution
and must sum to the signal-zero block entry chosen by
Definition~\texttt{def:block-encoding}.

\[
\begin{array}{lll}
\text{backend predicate} &
\texttt{oneTermRobinGamma3BoundaryBackendBranchContributionPredicate\_n3} &
\text{typed}\\
\text{backend target} &
\texttt{oneTermRobinGamma3BoundaryBackendProjectionSummationFieldTarget\_n3} &
\text{compiled}\\
\text{transcript} &
\texttt{oneTermRobinGamma3BoundaryBackendProjectionSummationFieldTarget\_n3\_transcript} &
\text{compiled}\\
\text{placeholder source} &
\texttt{placeholderFamilyIsBackendSourced} &
\text{false}\\
\text{placeholder closure} &
\texttt{placeholderMayCloseProjectionSummation} &
\text{false}\\
\text{backend field} &
\texttt{branchContribution : Fin\ 7 -> Coeff} &
\text{absent}.
\end{array}
\]

Middle classifies this as QBE-local finite projection bookkeeping, not a
Shukla--Vedula state-preparation theorem and not standard LCU composition.
The next lower packet should either expose a backend-sourced branch family
satisfying the predicate above or return a smaller typed obstruction explaining
why \(\texttt{BlockExtractionTarget}\) cannot yet expose that field.  Product
equality, LCU composition, block projection, normalized equality, block
correctness, circuit unitarity, and final extraction remain false.

\subsection{2026-05-28 Block-Extraction Backend Gap}

The narrow backend gap now compiles as
\[
  \texttt{oneTermRobinGamma3BoundaryBlockExtractionBackendGap\_n3}.
\]
It ties the obstruction to the actual \(\texttt{BlockExtractionTarget}\)
fields: the signal-zero block entry and its full-unitary entry are available
and equal for the focused entry, but the backend does not expose a sparse-slot
summand family for that entry.

\[
\begin{array}{lll}
\text{backend gap} &
\texttt{oneTermRobinGamma3BoundaryBlockExtractionBackendGap\_n3} &
\text{compiled}\\
\text{transcript} &
\texttt{oneTermRobinGamma3BoundaryBlockExtractionBackendGap\_n3\_transcript} &
\text{compiled}\\
\text{block-entry bridge} &
\texttt{gap.signalBlockEntry = gap.unitaryEntry} &
\text{compiled}\\
\text{sparse-slot summand family} &
\texttt{exposesBranchContributionField} &
\text{false}\\
\text{backend predicate} &
\texttt{oneTermRobinGamma3BoundaryBackendBranchContributionPredicate\_n3} &
\text{typed}\\
\text{fixed product theorem} &
\texttt{oneTermRobinGamma3ProductToCoefficientObligation\ 3\ 0\ 0} &
\text{false}.
\end{array}
\]

The next local interface is a backend-sourced
\(\texttt{branchContribution : Fin\ 7 -> Coeff}\) for
\(\texttt{contract.expectedTarget.blockMatrix[0,0]}\) satisfying
\[
  \texttt{oneTermRobinGamma3BoundaryBackendBranchContributionPredicate\_n3}.
\]
This remains QBE-local projection bookkeeping, not a Shukla--Vedula theorem
and not standard LCU composition.

\subsection{2026-05-28 Generic Block-Extraction Branch Target}

The backend now has a generic one-entry branch target,
\[
  \texttt{BlockExtractionBranchContributionTarget},
\]
with branch fold
\[
  \texttt{blockExtractionBranchContributionSum}.
\]
For the focused \(n=3\) Robin boundary entry, this interface is instantiated as
\[
  \texttt{oneTermRobinGamma3BoundaryBlockExtractionBranchContributionTarget\_n3}.
\]
It pairs the signal-zero entry \(\texttt{contract.expectedTarget.blockMatrix[0,0]}\)
with a seven-slot candidate family.  The selected branch is slot \(2\), and
Lean proves the local selected-slot identity in
\[
  \texttt{oneTermRobinGamma3BoundaryBlockExtractionBranchContributionTarget\_selected\_n3}
\]
and the selected contribution equality in
\[
  \texttt{oneTermRobinGamma3BoundaryBlockExtractionBranchContributionTarget\_selectedContribution\_eq\_n3}.
\]

The interface does not prove that the candidate family is produced by the
finite projection backend.  It also does not prove that the signal-zero block
entry equals the seven-branch fold.  These remain the two focused backend
obligations:
\[
\begin{array}{lll}
\text{backend source} &
\texttt{target.backendSource.proved} &
\text{false}\\
\text{branch summation} &
\texttt{target.branchSummationCorrect.proved} &
\text{false}.
\end{array}
\]
Thus the exact obstruction has been narrowed: \(\texttt{BlockExtractionTarget}\)
can be paired with a typed seven-slot branch target, but the backend-source
proof and signal-block branch-sum theorem are still absent.  The theorem
\[
  \texttt{oneTermRobinGamma3ProductToCoefficientObligation\ 3\ 0\ 0}
\]
and the LCU, block projection, normalized equality, block correctness, circuit
unitarity, and final extraction flags remain false.

\subsection{2026-05-28 All-Slot Backend Summand Formula}

The all-slot backend family now compiles as
\[
  \texttt{oneTermRobinGamma3BoundaryBackendBranchContribution\_n3}
  \colon \texttt{Fin 7 -> Coeff}.
\]
For each sparse slot, this family reads the focused seven-gate matrix at the
clean full-basis branch index and multiplies by the two sparse-register
projection amplitudes.  The backend target using this family is
\[
  \texttt{oneTermRobinGamma3BoundaryBackendBranchContributionTarget\_n3}.
\]

Lean proves the selected slot-\(2\) clause:
\[
  \texttt{oneTermRobinGamma3BoundaryBackendBranchContribution\_selected\_n3}.
\]
It also proves that the backend target's selected contribution is the accepted
slot-\(2\) contribution:
\[
  \texttt{oneTermRobinGamma3BoundaryBackendBranchContributionTarget\_selectedContribution\_eq\_n3}.
\]

\[
\begin{array}{lll}
\text{all-slot family} &
\texttt{oneTermRobinGamma3BoundaryBackendBranchContribution\_n3} &
\text{compiled}\\
\text{selected slot theorem} &
\texttt{oneTermRobinGamma3BoundaryBackendBranchContribution\_selected\_n3} &
\text{compiled}\\
\text{backend target} &
\texttt{oneTermRobinGamma3BoundaryBackendBranchContributionTarget\_n3} &
\text{compiled; semantic flags false}\\
\text{branch-sum theorem} &
\texttt{BlockExtractionBranchContributionTarget.projectionSummationStatement} &
\text{absent}\\
\text{fixed product theorem} &
\texttt{oneTermRobinGamma3ProductToCoefficientObligation\ 3\ 0\ 0} &
\text{false}.
\end{array}
\]

The remaining theorem is the signal-zero branch-sum statement
\[
  \texttt{BlockExtractionBranchContributionTarget.projectionSummationStatement\
  oneTermRobinGamma3BoundaryBackendBranchContributionTarget\_n3}.
\]
Equivalently, lower may prove the second conjunct of
\[
  \texttt{oneTermRobinGamma3BoundaryBackendBranchContributionPredicate\_n3\
  oneTermRobinGamma3BoundaryBackendBranchContribution\_n3}.
\]
This is QBE-local finite projection bookkeeping.  It is not supplied by the
Shukla--Vedula sparse-register preparation contract and not by standard LCU
composition.  Product-to-coefficient, LCU composition, block projection,
normalized equality, block correctness, circuit unitarity, and final
extraction remain false.

\subsection{2026-05-28 Projection-Statement Bridge}

The backend branch-sum target is now aligned with the generic
\(\texttt{BlockExtractionBranchContributionTarget}\) interface.  Lean proves
\[
  \texttt{oneTermRobinGamma3BoundaryBackendProjectionStatement\_unfold\_n3},
\]
which unfolds the generic projection statement into the equality between the
backend target block entry and the folded backend branch family.  Lean also
proves
\[
  \texttt{oneTermRobinGamma3BoundaryBackendProjectionStatement\_signalEntry\_n3},
\]
which identifies the Robin-local signal entry with the block entry of
\(\texttt{oneTermRobinGamma3BoundaryBackendBranchContributionTarget\_n3}\).

The conditional bridge
\[
  \texttt{oneTermRobinGamma3BoundaryBackendBranchContributionPredicate\_of\_targetProjection\_n3}
\]
states that a proof of
\[
  \texttt{BlockExtractionBranchContributionTarget.projectionSummationStatement\
  oneTermRobinGamma3BoundaryBackendBranchContributionTarget\_n3}
\]
is sufficient to close
\[
  \texttt{oneTermRobinGamma3BoundaryBackendBranchContributionPredicate\_n3\
  oneTermRobinGamma3BoundaryBackendBranchContribution\_n3}.
\]
This is only a conditional proof-DAG bridge.  It does not prove the
projection/summation equality, and it does not promote
\(\texttt{oneTermRobinGamma3ProductToCoefficientObligation\ 3\ 0\ 0}\) or any
LCU, block projection, normalized equality, block correctness, circuit
unitarity, or final extraction flag.

\subsection{2026-05-28 Fold-Support Packet}

The active full-entry target is still
\[
  \Lean{oneTermRobinGamma3BoundaryProjectionSummationTarget\_n3.signalUnitaryEntry}
  =
  \Lean{blockExtractionBranchContributionSum\
  oneTermRobinGamma3BoundaryBackendBranchContribution\_n3}.
\]
The fold-support packet proves only finite support data for this target.  Lean
now has
\[
  \Lean{oneTermRobinGamma3BoundaryBackendSelectedBranch\_mem\_fold\_n3},
\]
which says that the focused sparse slot \(2\) lies in the seven-slot fold
domain, and
\[
  \Lean{oneTermRobinGamma3BoundaryBackendUnitaryEntryFoldSupportTarget\_n3}.
\]
The transcript theorem for that packet checks that slot \(2\) still gives the
accepted \([32,32]\) branch contribution.

\[
\begin{array}{lll}
\text{Item} & \text{Lean declaration or field} & \text{Status}\\
\text{selected slot membership} &
\Lean{oneTermRobinGamma3BoundaryBackendSelectedBranch\_mem\_fold\_n3} &
\text{compiled}\\
\text{fold-support packet} &
\Lean{oneTermRobinGamma3BoundaryBackendUnitaryEntryFoldSupportTarget\_n3} &
\text{compiled; theorem-facing fields false}\\
\text{transcript check} &
\Lean{oneTermRobinGamma3BoundaryBackendUnitaryEntryFoldSupportTarget\_n3\_transcript} &
\text{compiled}\\
\text{full unitary-entry fold} &
\Lean{signalUnitaryEntry = blockExtractionBranchContributionSum\
oneTermRobinGamma3BoundaryBackendBranchContribution\_n3} &
\text{absent}\\
\text{fixed product theorem} &
\Lean{oneTermRobinGamma3ProductToCoefficientObligation\ 3\ 0\ 0} &
\text{false}.
\end{array}
\]

This remains QBE-local finite projection bookkeeping.  The Shukla--Vedula
uniform sparse-register result supplies only the external amplitude contract,
and standard LCU supplies only the later composition context.  Neither result
proves the full unitary-entry fold.

\subsection{2026-05-28 Prepared-Branch Expansion Target}

For a sparse slot \(s\), the prepared branch contribution is now defined as
the focused seven-gate diagonal entry at the paper branch basis index for
\(s\), multiplied by the ket-side and bra-side sparse-register amplitudes.
Lean proves the uniform summand formula
\[
  \Lean{oneTermRobinGamma3BoundaryPreparedBranchContribution\_formula\_n3}.
\]
The formula is packaged in
\[
  \Lean{oneTermRobinGamma3BoundaryPreparedBranchExpansionTarget\_n3}
\]
and checked by
\[
  \Lean{oneTermRobinGamma3BoundaryPreparedBranchExpansionTarget\_n3\_transcript}.
\]

\[
\begin{array}{lll}
\text{Item} & \text{Lean declaration or field} & \text{Status}\\
\text{prepared summand formula} &
\Lean{oneTermRobinGamma3BoundaryPreparedBranchContribution\_formula\_n3} &
\text{compiled}\\
\text{prepared branch expansion packet} &
\Lean{oneTermRobinGamma3BoundaryPreparedBranchExpansionTarget\_n3} &
\text{compiled; theorem-facing fields false}\\
\text{target formula proof} &
\Lean{oneTermRobinGamma3BoundaryPreparedBranchExpansionTarget\_formula\_n3} &
\text{compiled}\\
\text{raw signal-zero entry as prepared fold} &
\Lean{oneTermRobinGamma3BoundaryPreparedBranchExpansionTarget\_n3.preparedProjectionBackendStatement} &
\text{absent}\\
\text{fixed product theorem} &
\Lean{oneTermRobinGamma3ProductToCoefficientObligation\ 3\ 0\ 0} &
\text{false}.
\end{array}
\]

The remaining target is still
\[
  \Lean{oneTermRobinGamma3BoundaryProjectionSummationTarget\_n3.signalUnitaryEntry}
  =
  \Lean{blockExtractionBranchContributionSum\
  oneTermRobinGamma3BoundaryBackendBranchContribution\_n3}.
\]
The new obstruction is narrower than the previous fold-support gap: the
summands themselves are no longer ambiguous.  The missing proof is the
prepared projection backend connecting the raw signal-zero entry exposed by
\(\Lean{CircuitMatrixSemantics}\) to the seven prepared summands.  The
Shukla--Vedula row supplies only the sparse-register amplitude contract, and
standard LCU supplies only later composition context.  Product-to-coefficient,
LCU composition, block projection, normalized equality, block correctness,
circuit unitarity, and final extraction remain false.

\subsection{2026-06-02 Source-Prepared Route Witness}

Define the source-prepared projection entry to be the clean entry of
\[
  \Lean{oneTermRobinGamma3BoundaryPreparedCompositeCircuitSemantics\_n3\ H}
\]
at \(\Lean{oneTermRobinGamma3BoundarySparseCleanIndex\_n3}\).  This entry is
the Lean representative of the
\(H_W^{(\kappa)\dagger} U_{\gamma_3,\mathrm{bdry}} H_W^{(\kappa)}\)
singleton route supplied by Eq.~\texttt{arbitrary sparcity} and
Fig.~\texttt{fig:1 term ROBIN}.

Lean now proves the route witness
\[
  \Lean{oneTermRobinGamma3BoundarySourcePreparedProjectionTarget\_preparedBackendEval\_n3}.
\]
It exposes the field
\[
  \Lean{preparedSingletonToBackendEvalStatement}
\]
of
\[
  \Lean{oneTermRobinGamma3BoundarySourcePreparedProjectionTarget\_n3\ H\ env}.
\]
Under the clean-column contract
\[
  \Lean{oneTermRobinGamma3BoundaryHWKappaUniformColumnAllSlotsStatement\_n3\ H},
\]
this field states that the prepared singleton clean entry evaluates to the
backend branch fold.  The companion theorem
\[
  \Lean{oneTermRobinGamma3BoundarySourcePreparedProjectionTarget\_feedsProductRoute\_n3}
\]
checks that this prepared backend bridge is available to the focused product
route and that the product, LCU, block-projection, block-correctness, and final
extraction flags remain false.

The remaining theorem is not a Shukla--Vedula theorem and not an LCU theorem.
It is the QBE-local finite composition statement
\[
  \Lean{(oneTermRobinGamma3BoundarySourcePreparedProjectionTarget\_n3\ H\ env).activeToPreparedSingletonEvalStatement}.
\]
Equivalently, lower may target
\[
  \Lean{(oneTermRobinGamma3BoundaryActivePreparedCircuitFieldTarget\_n3\ H\ env).activePreparedCompositeEvalStatement}.
\]
After this statement is proved, the evaluated backend fold can be recovered by
\[
  \Lean{oneTermRobinGamma3BoundaryEvaluatedBackendFoldStatement\_iff\_activePreparedCircuitField\_n3}.
\]
The raw \(H\)-free fold remains a stronger diagnostic statement until it is
derived through the source-prepared route.  The theorem
\[
  \Lean{oneTermRobinGamma3ProductToCoefficientObligation\ 3\ 0\ 0}
\]
and all theorem-facing semantic flags remain false.

\subsection{2026-06-03 Backend Expansion Packet}

Define the backend expansion target to be
\[
  \Lean{oneTermRobinGamma3BoundaryBackendBranchContributionTarget\_n3.backendExpansionStatement}.
\]
By the compiled bridge
\[
  \Lean{oneTermRobinGamma3BoundaryBackendExpansionStatement\_equivUnitaryEntryFold\_n3},
\]
this is equivalent to the raw symbolic fold
\[
  \Lean{oneTermRobinGamma3BoundaryProjectionSummationTarget\_n3.signalUnitaryEntry}
  =
  \Lean{blockExtractionBranchContributionSum\
  oneTermRobinGamma3BoundaryBackendBranchContribution\_n3}.
\]

The source anchors are Theorem~\texttt{theorem: 1 term robin},
Eq.~\texttt{ROBIN clarified}, Eq.~\texttt{arbitrary sparcity},
Fig.~\texttt{fig:1 term ROBIN}, and Definition~\texttt{def:block-encoding}.
Equation~\texttt{arbitrary sparcity} supplies only the clean-column contract
for \(H_W^{(\kappa)}\).  The finite expansion of the signal-zero entry as the
seven sparse-slot backend fold is a QBE-local matrix-projection theorem.

Lean now has the conditional route
\[
  \Lean{oneTermRobinGamma3BoundaryUncastPreparedSandwichEval\_of\_backendExpansion\_n3}.
\]
It says that the backend expansion, together with the all-slot
\(H_W^{(\kappa)}\) clean-column contract, closes the prepared-sandwich
evaluated target
\[
  \Lean{oneTermRobinGamma3BoundaryUncastPreparedSandwichEvalStatement\_n3\ H\ env}.
\]
The alternative route
\[
  \Lean{oneTermRobinGamma3BoundaryUncastPreparedSandwichEval\_of\_unitaryEntryFold\_n3}
\]
uses the equivalent raw unitary-entry fold.  These theorems are route wiring
only.  The backend expansion, product-to-coefficient theorem, LCU composition,
block projection, normalized equality, block correctness, circuit unitarity,
and final extraction remain unproved.

\subsection{2026-06-05 Focused Product-Obligation Prepared Bridge}

Definition first: the theorem-facing projection entry for the focused
\(\gamma_3\) route is
\[
  \Lean{(oneTermRobinGamma3BoundarySourcePreparedProjectionTarget\_n3\ H\ env).preparedProjectionEntry}.
\]
This is the prepared singleton clean entry of
\[
  \Lean{oneTermRobinGamma3BoundaryPreparedCompositeCircuitSemantics\_n3\ H}.
\]
It is not the standalone \(H\)-free active entry.

Lean now exposes this source-prepared backend equality at the fixed product
obligation by
\[
  \Lean{oneTermRobinGamma3BoundaryFocusedProductObligation\_preparedProjectionBackendEval\_n3}.
\]
The theorem uses
\[
  \Lean{oneTermRobinGamma3BoundaryHWKappaUniformColumnAllSlotsStatement\_n3\ H}
\]
only as the contract for the \(H_W^{(\kappa)}\) clean-column amplitudes from
Eq.~\texttt{arbitrary sparcity}.  It then routes the prepared backend fold
through the finite projection/product bridge for
\[
  \Lean{oneTermRobinGamma3ProductToCoefficientObligation\ 3\ 0\ 0}.
\]

\[
\begin{array}{lll}
\text{Block} & \text{Lean declaration or field} & \text{Status}\\
\text{prepared target-field backend bridge} &
\Lean{oneTermRobinGamma3BoundarySourcePreparedProjectionTarget\_preparedProjectionEntryEval\_eq\_backend\_n3} &
\text{compiled under the all-slot contract}\\
\text{focused product-obligation bridge} &
\Lean{oneTermRobinGamma3BoundaryFocusedProductObligation\_preparedProjectionBackendEval\_n3} &
\text{compiled; theorem-facing flags false}\\
\text{active/prepared selected-entry equality} &
\Lean{(oneTermRobinGamma3BoundarySourcePreparedProjectionTarget\_n3\ H\ env).activeToPreparedSingletonEvalStatement} &
\text{absent}\\
\text{active/prepared circuit field} &
\Lean{(oneTermRobinGamma3BoundaryActivePreparedCircuitFieldTarget\_n3\ H\ env).activePreparedCompositeEvalStatement} &
\text{equivalent next target}\\
\text{\(H\)-free active fold} &
\Lean{oneTermRobinGamma3BoundaryBackendBranchContributionTarget\_n3.backendExpansionStatement} &
\text{diagnostic/backlog unless recovered through the prepared route}.
\end{array}
\]

The source-dependency classification is unchanged.  Shukla--Vedula is an
external contract for uniform sparse-register preparation.  The remaining
active/prepared selected-entry equality is QBE-local finite
\(\Lean{CircuitMatrixSemantics}\) bookkeeping.  Standard LCU and final
block-composition contracts are downstream and do not prove this local
projection entry.  Product-to-coefficient, branch decomposition, LCU
composition, block projection, normalized equality, block correctness, circuit
unitarity, and final extraction remain false.
